% Options for packages loaded elsewhere
% Options for packages loaded elsewhere
\PassOptionsToPackage{unicode}{hyperref}
\PassOptionsToPackage{hyphens}{url}
\PassOptionsToPackage{dvipsnames,svgnames,x11names}{xcolor}
%
\documentclass[
  letterpaper,
]{book}
\usepackage{xcolor}
\usepackage{amsmath,amssymb}
\setcounter{secnumdepth}{5}
\usepackage{iftex}
\ifPDFTeX
  \usepackage[T1]{fontenc}
  \usepackage[utf8]{inputenc}
  \usepackage{textcomp} % provide euro and other symbols
\else % if luatex or xetex
  \usepackage{unicode-math} % this also loads fontspec
  \defaultfontfeatures{Scale=MatchLowercase}
  \defaultfontfeatures[\rmfamily]{Ligatures=TeX,Scale=1}
\fi
\usepackage{lmodern}
\ifPDFTeX\else
  % xetex/luatex font selection
\fi
% Use upquote if available, for straight quotes in verbatim environments
\IfFileExists{upquote.sty}{\usepackage{upquote}}{}
\IfFileExists{microtype.sty}{% use microtype if available
  \usepackage[]{microtype}
  \UseMicrotypeSet[protrusion]{basicmath} % disable protrusion for tt fonts
}{}
\makeatletter
\@ifundefined{KOMAClassName}{% if non-KOMA class
  \IfFileExists{parskip.sty}{%
    \usepackage{parskip}
  }{% else
    \setlength{\parindent}{0pt}
    \setlength{\parskip}{6pt plus 2pt minus 1pt}}
}{% if KOMA class
  \KOMAoptions{parskip=half}}
\makeatother
% Make \paragraph and \subparagraph free-standing
\makeatletter
\ifx\paragraph\undefined\else
  \let\oldparagraph\paragraph
  \renewcommand{\paragraph}{
    \@ifstar
      \xxxParagraphStar
      \xxxParagraphNoStar
  }
  \newcommand{\xxxParagraphStar}[1]{\oldparagraph*{#1}\mbox{}}
  \newcommand{\xxxParagraphNoStar}[1]{\oldparagraph{#1}\mbox{}}
\fi
\ifx\subparagraph\undefined\else
  \let\oldsubparagraph\subparagraph
  \renewcommand{\subparagraph}{
    \@ifstar
      \xxxSubParagraphStar
      \xxxSubParagraphNoStar
  }
  \newcommand{\xxxSubParagraphStar}[1]{\oldsubparagraph*{#1}\mbox{}}
  \newcommand{\xxxSubParagraphNoStar}[1]{\oldsubparagraph{#1}\mbox{}}
\fi
\makeatother

\usepackage{color}
\usepackage{fancyvrb}
\newcommand{\VerbBar}{|}
\newcommand{\VERB}{\Verb[commandchars=\\\{\}]}
\DefineVerbatimEnvironment{Highlighting}{Verbatim}{commandchars=\\\{\}}
% Add ',fontsize=\small' for more characters per line
\usepackage{framed}
\definecolor{shadecolor}{RGB}{241,243,245}
\newenvironment{Shaded}{\begin{snugshade}}{\end{snugshade}}
\newcommand{\AlertTok}[1]{\textcolor[rgb]{0.68,0.00,0.00}{#1}}
\newcommand{\AnnotationTok}[1]{\textcolor[rgb]{0.37,0.37,0.37}{#1}}
\newcommand{\AttributeTok}[1]{\textcolor[rgb]{0.40,0.45,0.13}{#1}}
\newcommand{\BaseNTok}[1]{\textcolor[rgb]{0.68,0.00,0.00}{#1}}
\newcommand{\BuiltInTok}[1]{\textcolor[rgb]{0.00,0.23,0.31}{#1}}
\newcommand{\CharTok}[1]{\textcolor[rgb]{0.13,0.47,0.30}{#1}}
\newcommand{\CommentTok}[1]{\textcolor[rgb]{0.37,0.37,0.37}{#1}}
\newcommand{\CommentVarTok}[1]{\textcolor[rgb]{0.37,0.37,0.37}{\textit{#1}}}
\newcommand{\ConstantTok}[1]{\textcolor[rgb]{0.56,0.35,0.01}{#1}}
\newcommand{\ControlFlowTok}[1]{\textcolor[rgb]{0.00,0.23,0.31}{\textbf{#1}}}
\newcommand{\DataTypeTok}[1]{\textcolor[rgb]{0.68,0.00,0.00}{#1}}
\newcommand{\DecValTok}[1]{\textcolor[rgb]{0.68,0.00,0.00}{#1}}
\newcommand{\DocumentationTok}[1]{\textcolor[rgb]{0.37,0.37,0.37}{\textit{#1}}}
\newcommand{\ErrorTok}[1]{\textcolor[rgb]{0.68,0.00,0.00}{#1}}
\newcommand{\ExtensionTok}[1]{\textcolor[rgb]{0.00,0.23,0.31}{#1}}
\newcommand{\FloatTok}[1]{\textcolor[rgb]{0.68,0.00,0.00}{#1}}
\newcommand{\FunctionTok}[1]{\textcolor[rgb]{0.28,0.35,0.67}{#1}}
\newcommand{\ImportTok}[1]{\textcolor[rgb]{0.00,0.46,0.62}{#1}}
\newcommand{\InformationTok}[1]{\textcolor[rgb]{0.37,0.37,0.37}{#1}}
\newcommand{\KeywordTok}[1]{\textcolor[rgb]{0.00,0.23,0.31}{\textbf{#1}}}
\newcommand{\NormalTok}[1]{\textcolor[rgb]{0.00,0.23,0.31}{#1}}
\newcommand{\OperatorTok}[1]{\textcolor[rgb]{0.37,0.37,0.37}{#1}}
\newcommand{\OtherTok}[1]{\textcolor[rgb]{0.00,0.23,0.31}{#1}}
\newcommand{\PreprocessorTok}[1]{\textcolor[rgb]{0.68,0.00,0.00}{#1}}
\newcommand{\RegionMarkerTok}[1]{\textcolor[rgb]{0.00,0.23,0.31}{#1}}
\newcommand{\SpecialCharTok}[1]{\textcolor[rgb]{0.37,0.37,0.37}{#1}}
\newcommand{\SpecialStringTok}[1]{\textcolor[rgb]{0.13,0.47,0.30}{#1}}
\newcommand{\StringTok}[1]{\textcolor[rgb]{0.13,0.47,0.30}{#1}}
\newcommand{\VariableTok}[1]{\textcolor[rgb]{0.07,0.07,0.07}{#1}}
\newcommand{\VerbatimStringTok}[1]{\textcolor[rgb]{0.13,0.47,0.30}{#1}}
\newcommand{\WarningTok}[1]{\textcolor[rgb]{0.37,0.37,0.37}{\textit{#1}}}

\usepackage{longtable,booktabs,array}
\usepackage{calc} % for calculating minipage widths
% Correct order of tables after \paragraph or \subparagraph
\usepackage{etoolbox}
\makeatletter
\patchcmd\longtable{\par}{\if@noskipsec\mbox{}\fi\par}{}{}
\makeatother
% Allow footnotes in longtable head/foot
\IfFileExists{footnotehyper.sty}{\usepackage{footnotehyper}}{\usepackage{footnote}}
\makesavenoteenv{longtable}
\usepackage{graphicx}
\makeatletter
\newsavebox\pandoc@box
\newcommand*\pandocbounded[1]{% scales image to fit in text height/width
  \sbox\pandoc@box{#1}%
  \Gscale@div\@tempa{\textheight}{\dimexpr\ht\pandoc@box+\dp\pandoc@box\relax}%
  \Gscale@div\@tempb{\linewidth}{\wd\pandoc@box}%
  \ifdim\@tempb\p@<\@tempa\p@\let\@tempa\@tempb\fi% select the smaller of both
  \ifdim\@tempa\p@<\p@\scalebox{\@tempa}{\usebox\pandoc@box}%
  \else\usebox{\pandoc@box}%
  \fi%
}
% Set default figure placement to htbp
\def\fps@figure{htbp}
\makeatother





\setlength{\emergencystretch}{3em} % prevent overfull lines

\providecommand{\tightlist}{%
  \setlength{\itemsep}{0pt}\setlength{\parskip}{0pt}}



 


\makeatletter
\@ifpackageloaded{tcolorbox}{}{\usepackage[skins,breakable]{tcolorbox}}
\@ifpackageloaded{fontawesome5}{}{\usepackage{fontawesome5}}
\definecolor{quarto-callout-color}{HTML}{909090}
\definecolor{quarto-callout-note-color}{HTML}{0758E5}
\definecolor{quarto-callout-important-color}{HTML}{CC1914}
\definecolor{quarto-callout-warning-color}{HTML}{EB9113}
\definecolor{quarto-callout-tip-color}{HTML}{00A047}
\definecolor{quarto-callout-caution-color}{HTML}{FC5300}
\definecolor{quarto-callout-color-frame}{HTML}{acacac}
\definecolor{quarto-callout-note-color-frame}{HTML}{4582ec}
\definecolor{quarto-callout-important-color-frame}{HTML}{d9534f}
\definecolor{quarto-callout-warning-color-frame}{HTML}{f0ad4e}
\definecolor{quarto-callout-tip-color-frame}{HTML}{02b875}
\definecolor{quarto-callout-caution-color-frame}{HTML}{fd7e14}
\makeatother
\makeatletter
\@ifpackageloaded{bookmark}{}{\usepackage{bookmark}}
\makeatother
\makeatletter
\@ifpackageloaded{caption}{}{\usepackage{caption}}
\AtBeginDocument{%
\ifdefined\contentsname
  \renewcommand*\contentsname{Table of contents}
\else
  \newcommand\contentsname{Table of contents}
\fi
\ifdefined\listfigurename
  \renewcommand*\listfigurename{List of Figures}
\else
  \newcommand\listfigurename{List of Figures}
\fi
\ifdefined\listtablename
  \renewcommand*\listtablename{List of Tables}
\else
  \newcommand\listtablename{List of Tables}
\fi
\ifdefined\figurename
  \renewcommand*\figurename{Figure}
\else
  \newcommand\figurename{Figure}
\fi
\ifdefined\tablename
  \renewcommand*\tablename{Table}
\else
  \newcommand\tablename{Table}
\fi
}
\@ifpackageloaded{float}{}{\usepackage{float}}
\floatstyle{ruled}
\@ifundefined{c@chapter}{\newfloat{codelisting}{h}{lop}}{\newfloat{codelisting}{h}{lop}[chapter]}
\floatname{codelisting}{Listing}
\newcommand*\listoflistings{\listof{codelisting}{List of Listings}}
\makeatother
\makeatletter
\makeatother
\makeatletter
\@ifpackageloaded{caption}{}{\usepackage{caption}}
\@ifpackageloaded{subcaption}{}{\usepackage{subcaption}}
\makeatother
\usepackage{bookmark}
\IfFileExists{xurl.sty}{\usepackage{xurl}}{} % add URL line breaks if available
\urlstyle{same}
\hypersetup{
  pdftitle={MoBPS Cookbook},
  pdfauthor={MoBPS Development Team},
  colorlinks=true,
  linkcolor={Maroon},
  filecolor={Maroon},
  citecolor={Blue},
  urlcolor={Blue},
  pdfcreator={LaTeX via pandoc}}


\title{MoBPS Cookbook}
\usepackage{etoolbox}
\makeatletter
\providecommand{\subtitle}[1]{% add subtitle to \maketitle
  \apptocmd{\@title}{\par {\large #1 \par}}{}{}
}
\makeatother
\subtitle{A Practical Guide to Modular Breeding Program Simulator}
\author{MoBPS Development Team}
\date{2025-11-16}
\begin{document}
\frontmatter
\maketitle

\renewcommand*\contentsname{Table of contents}
{
\hypersetup{linkcolor=}
\setcounter{tocdepth}{2}
\tableofcontents
}

\mainmatter
\bookmarksetup{startatroot}

\chapter*{Welcome}\label{welcome}
\addcontentsline{toc}{chapter}{Welcome}

\markboth{Welcome}{Welcome}

Welcome to the \textbf{MoBPS Cookbook} - a practical, hands-on guide to
using the Modular Breeding Program Simulator!

\includegraphics[width=2.08333in,height=\textheight,keepaspectratio]{images/mobps-logo.png}

\section*{What is MoBPS?}\label{what-is-mobps}
\addcontentsline{toc}{section}{What is MoBPS?}

\markright{What is MoBPS?}

MoBPS (Modular Breeding Program Simulator) is a powerful R package for
stochastic simulation of breeding programs. Whether you're working with:

\begin{itemize}
\tightlist
\item
  \textbf{Animal breeding} (dairy cattle, pigs, chickens, etc.)
\item
  \textbf{Plant breeding} (maize, wheat, crops with selfing/DH)
\item
  \textbf{Conservation programs}
\item
  \textbf{Research on breeding strategies}
\end{itemize}

MoBPS provides the flexibility and tools you need to simulate, analyze,
and optimize breeding schemes.

\section*{About This Cookbook}\label{about-this-cookbook}
\addcontentsline{toc}{section}{About This Cookbook}

\markright{About This Cookbook}

Unlike traditional reference manuals, this \textbf{cookbook} is designed
to:

\begin{itemize}
\tightlist
\item
  ✅ \textbf{Show you HOW} to accomplish specific tasks
\item
  ✅ \textbf{Provide working examples} you can copy and adapt
\item
  ✅ \textbf{Build progressively} from simple to complex scenarios
\item
  ✅ \textbf{Focus on practical workflows} rather than exhaustive
  parameter lists
\end{itemize}

\section*{Who Should Use This Book?}\label{who-should-use-this-book}
\addcontentsline{toc}{section}{Who Should Use This Book?}

\markright{Who Should Use This Book?}

This cookbook is for:

\begin{itemize}
\tightlist
\item
  \textbf{Beginners} learning MoBPS for the first time
\item
  \textbf{Practitioners} designing breeding programs
\item
  \textbf{Researchers} evaluating breeding strategies
\item
  \textbf{Students} studying quantitative genetics and breeding
\end{itemize}

\section*{How to Use This Book}\label{how-to-use-this-book}
\addcontentsline{toc}{section}{How to Use This Book}

\markright{How to Use This Book}

The book is organized into 5 main parts:

\begin{enumerate}
\def\labelenumi{\arabic{enumi}.}
\tightlist
\item
  \textbf{Getting Started} - Installation, core concepts, your first
  simulation
\item
  \textbf{Building Populations \& Traits} - Creating realistic founder
  populations
\item
  \textbf{Selection, Mating \& Breeding} - Simulating breeding programs
\item
  \textbf{Genotyping \& Genomic Analysis} - Working with SNP data and
  predictions
\item
  \textbf{Advanced Topics \& Reference} - Expert techniques and function
  references
\end{enumerate}

\textbf{Recommended approach:} - Start with Part 1 if you're new to
MoBPS - Use Ctrl+F / Search to find specific topics - Run the code
examples in your own R session - Modify examples to fit your needs

\section*{Getting Help}\label{getting-help}
\addcontentsline{toc}{section}{Getting Help}

\markright{Getting Help}

If you get stuck:

\begin{itemize}
\tightlist
\item
  📧 Contact Torsten Pook:
  \href{mailto:Torsten.pook@wur.nl}{\nolinkurl{Torsten.pook@wur.nl}}
\item
  🌐 Visit \href{https://www.mobps.de}{www.mobps.de}
\item
  📦 GitHub repository:
  \href{https://github.com/tpook92/MoBPS}{github.com/tpook92/MoBPS}
\end{itemize}

\section*{Authors}\label{authors}
\addcontentsline{toc}{section}{Authors}

\markright{Authors}

\textbf{MoBPS Development Team:}

\begin{itemize}
\tightlist
\item
  Torsten Pook (Wageningen University \& Research)
\item
  Martin Schlather (University of Mannheim)
\item
  Henner Simianer (University of Göttingen)
\end{itemize}

\begin{center}\rule{0.5\linewidth}{0.5pt}\end{center}

\textbf{Ready to get started?} Head to
\href{chapters/01-introduction.qmd}{Chapter 1: Introduction} or jump to
any topic using the navigation menu!

\section*{The MoBPS Mascot}\label{the-mobps-mascot}
\addcontentsline{toc}{section}{The MoBPS Mascot}

\markright{The MoBPS Mascot}

\begin{figure}[H]

{\centering \includegraphics[width=3.125in,height=\textheight,keepaspectratio]{images/pug.png}

}

\caption{This little pug (German: Mobps) is our name-giver!}

\end{figure}%

\part{Getting Started}

\chapter{Introduction to MoBPS}\label{sec-introduction}

\section{Learning Objectives}\label{learning-objectives}

By the end of this chapter, you will:

\begin{itemize}
\tightlist
\item
  Understand what MoBPS is and what it can do
\item
  Know how to install MoBPS and its dependencies
\item
  Be able to load the package and verify installation
\item
  Understand the basic workflow of a MoBPS simulation
\end{itemize}

\section{What is MoBPS?}\label{what-is-mobps-1}

\textbf{MoBPS} (Modular Breeding Program Simulator) is an R package for
stochastic simulation of breeding programs. It provides a flexible
framework to:

\begin{itemize}
\tightlist
\item
  \textbf{Create} realistic founder populations with complex genetic
  architectures
\item
  \textbf{Simulate} multiple generations of breeding actions (selection,
  mating, culling)
\item
  \textbf{Analyze} genetic gain, diversity, inbreeding, and population
  structure
\item
  \textbf{Compare} different breeding strategies
\item
  \textbf{Optimize} breeding programs for specific goals
\end{itemize}

\subsection{Why Simulate Breeding
Programs?}\label{why-simulate-breeding-programs}

Breeding program simulation allows you to:

\begin{enumerate}
\def\labelenumi{\arabic{enumi}.}
\tightlist
\item
  \textbf{Test strategies} before implementing them in reality
\item
  \textbf{Understand trade-offs} between genetic gain and diversity
\item
  \textbf{Optimize selection} and mating decisions
\item
  \textbf{Predict long-term} outcomes over many generations
\item
  \textbf{Train} future breeders in a risk-free environment
\end{enumerate}

\subsection{Key Features}\label{key-features}

MoBPS offers exceptional flexibility:

\begin{itemize}
\tightlist
\item
  ✅ Works with \textbf{both plant and animal breeding}
\item
  ✅ Supports \textbf{single or multiple traits}
\item
  ✅ Handles \textbf{additive, dominance, and epistatic effects}
\item
  ✅ Simulates \textbf{realistic genomic data} or uses real data
\item
  ✅ Integrates with \textbf{external prediction software} (BLUP, GBLUP,
  Bayesian)
\item
  ✅ Provides \textbf{rich visualization} and analysis tools
\item
  ✅ Scales from small experiments to \textbf{large commercial programs}
\end{itemize}

\section{Installation}\label{sec-installation}

\subsection{System Requirements}\label{system-requirements}

MoBPS requires:

\begin{itemize}
\tightlist
\item
  \textbf{R version 3.0 or higher} (R 4.0+ recommended)
\item
  For large-scale simulations: \texttt{RandomFieldsUtils} and
  \texttt{miraculix} (optional but highly recommended)
\end{itemize}

\subsection{Step 1: Platform-Specific
Prerequisites}\label{step-1-platform-specific-prerequisites}

\subsubsection{Windows Users}\label{windows-users}

Some Windows systems require \textbf{Rtools}:

\begin{enumerate}
\def\labelenumi{\arabic{enumi}.}
\tightlist
\item
  Download from: \url{https://cran.r-project.org/bin/windows/Rtools/}
\item
  Install the version matching your R version
\item
  Verify with: \texttt{pkgbuild::find\_rtools()}
\end{enumerate}

\subsubsection{Mac Users}\label{mac-users}

Some Mac systems require \textbf{gfortran/homebrew}:

\begin{enumerate}
\def\labelenumi{\arabic{enumi}.}
\tightlist
\item
  Download from: \url{https://mac.r-project.org/tools/}
\item
  Follow installation instructions for your macOS version
\end{enumerate}

\subsubsection{Linux Users}\label{linux-users}

Usually no special prerequisites needed. Ensure build tools are
installed:

\begin{Shaded}
\begin{Highlighting}[]
\CommentTok{\# Ubuntu/Debian}
\FunctionTok{sudo}\NormalTok{ apt{-}get install build{-}essential gfortran}

\CommentTok{\# Fedora/RHEL}
\FunctionTok{sudo}\NormalTok{ dnf install gcc{-}gfortran}
\end{Highlighting}
\end{Shaded}

\subsection{Step 2: Install Performance Packages (Optional but
Recommended)}\label{step-2-install-performance-packages-optional-but-recommended}

For \textbf{significant performance improvements} in large simulations:

\begin{Shaded}
\begin{Highlighting}[]
\CommentTok{\# Install devtools if not already installed}
\ControlFlowTok{if}\NormalTok{ (}\SpecialCharTok{!}\FunctionTok{require}\NormalTok{(}\StringTok{"devtools"}\NormalTok{)) }\FunctionTok{install.packages}\NormalTok{(}\StringTok{"devtools"}\NormalTok{)}

\CommentTok{\# Install RandomFieldsUtils (install this FIRST)}
\NormalTok{devtools}\SpecialCharTok{::}\FunctionTok{install\_github}\NormalTok{(}\StringTok{"tpook92/RandomFieldsUtils"}\NormalTok{)}

\CommentTok{\# Then install miraculix (install SECOND)}
\NormalTok{devtools}\SpecialCharTok{::}\FunctionTok{install\_github}\NormalTok{(}\StringTok{"tpook92/miraculix"}\NormalTok{)}
\end{Highlighting}
\end{Shaded}

\textbf{Important:} - Install \texttt{RandomFieldsUtils} \textbf{before}
\texttt{miraculix} - Use compatible versions (both v1.6.0.2 recommended)
- These packages dramatically speed up genomic computations

\subsection{Step 3: Install MoBPS}\label{step-3-install-mobps}

Install the latest MoBPS from GitHub (recommended over CRAN for latest
features):

\begin{Shaded}
\begin{Highlighting}[]
\CommentTok{\# Install MoBPS from GitHub}
\NormalTok{devtools}\SpecialCharTok{::}\FunctionTok{install\_github}\NormalTok{(}\StringTok{"tpook92/MoBPS"}\NormalTok{, }\AttributeTok{subdir=}\StringTok{"pkg"}\NormalTok{)}
\end{Highlighting}
\end{Shaded}

\textbf{Why GitHub over CRAN?} - GitHub has the latest features and bug
fixes - Regular updates and improvements - CRAN version may be outdated
(last update: November 2021)

\subsection{Step 4: Install Genetic Maps
(Optional)}\label{step-4-install-genetic-maps-optional}

For working with real species data:

\begin{Shaded}
\begin{Highlighting}[]
\CommentTok{\# Install MoBPS maps package (includes Ensembl maps)}
\NormalTok{devtools}\SpecialCharTok{::}\FunctionTok{install\_github}\NormalTok{(}\StringTok{"tpook92/MoBPS"}\NormalTok{, }\AttributeTok{subdir=}\StringTok{"pkg{-}maps"}\NormalTok{)}
\end{Highlighting}
\end{Shaded}

This package includes genetic maps for common species from Ensembl.

\subsection{Step 5: Verify
Installation}\label{step-5-verify-installation}

Load the package and check version:

\begin{Shaded}
\begin{Highlighting}[]
\CommentTok{\# Load MoBPS}
\FunctionTok{library}\NormalTok{(MoBPS)}

\CommentTok{\# Check version}
\FunctionTok{packageVersion}\NormalTok{(}\StringTok{"MoBPS"}\NormalTok{)}

\CommentTok{\# Should see something like: \textquotesingle{}1.13.0\textquotesingle{} or higher}
\end{Highlighting}
\end{Shaded}

\subsection{Additional Packages}\label{additional-packages}

MoBPS can integrate with various prediction software. These are
\textbf{optional} and only needed for specific analyses:

\begin{Shaded}
\begin{Highlighting}[]
\CommentTok{\# For genomic prediction}
\FunctionTok{install.packages}\NormalTok{(}\StringTok{"rrBLUP"}\NormalTok{)    }\CommentTok{\# GBLUP}
\FunctionTok{install.packages}\NormalTok{(}\StringTok{"BGLR"}\NormalTok{)      }\CommentTok{\# Bayesian methods}
\FunctionTok{install.packages}\NormalTok{(}\StringTok{"sommer"}\NormalTok{)    }\CommentTok{\# Mixed models}

\CommentTok{\# For visualization}
\FunctionTok{install.packages}\NormalTok{(}\StringTok{"ggplot2"}\NormalTok{)   }\CommentTok{\# Advanced plotting}
\FunctionTok{install.packages}\NormalTok{(}\StringTok{"viridis"}\NormalTok{)   }\CommentTok{\# Color palettes}
\end{Highlighting}
\end{Shaded}

\section{Basic Workflow Overview}\label{sec-workflow}

Every MoBPS simulation follows the same basic structure:

\subsection{1. Create a Founder
Population}\label{create-a-founder-population}

Use \texttt{creating.diploid()} to initialize your population:

\begin{Shaded}
\begin{Highlighting}[]
\NormalTok{population }\OtherTok{\textless{}{-}} \FunctionTok{creating.diploid}\NormalTok{(}
  \AttributeTok{nsnp =} \DecValTok{1000}\NormalTok{,          }\CommentTok{\# Number of SNP markers}
  \AttributeTok{nindi =} \DecValTok{100}\NormalTok{,          }\CommentTok{\# Number of individuals}
  \AttributeTok{n.additive =} \DecValTok{100}\NormalTok{,     }\CommentTok{\# Number of QTLs}
  \AttributeTok{mean.target =} \DecValTok{100}     \CommentTok{\# Mean breeding value}
\NormalTok{)}
\end{Highlighting}
\end{Shaded}

\subsection{2. Simulate Breeding
Actions}\label{simulate-breeding-actions}

Use \texttt{breeding.diploid()} to perform selection, mating, and other
actions:

\begin{Shaded}
\begin{Highlighting}[]
\NormalTok{population }\OtherTok{\textless{}{-}} \FunctionTok{breeding.diploid}\NormalTok{(}
\NormalTok{  population,}
  \AttributeTok{selection.size =} \FunctionTok{c}\NormalTok{(}\DecValTok{10}\NormalTok{, }\DecValTok{10}\NormalTok{),  }\CommentTok{\# 10 males, 10 females}
  \AttributeTok{breeding.size =} \FunctionTok{c}\NormalTok{(}\DecValTok{50}\NormalTok{, }\DecValTok{50}\NormalTok{)     }\CommentTok{\# Generate 50 male, 50 female offspring}
\NormalTok{)}
\end{Highlighting}
\end{Shaded}

\subsection{3. Extract and Analyze
Results}\label{extract-and-analyze-results}

Use \texttt{get.*()} functions to extract information:

\begin{Shaded}
\begin{Highlighting}[]
\CommentTok{\# Get breeding values}
\NormalTok{bv }\OtherTok{\textless{}{-}} \FunctionTok{get.bv}\NormalTok{(population, }\AttributeTok{gen =} \DecValTok{2}\NormalTok{)}

\CommentTok{\# Get phenotypes}
\NormalTok{pheno }\OtherTok{\textless{}{-}} \FunctionTok{get.pheno}\NormalTok{(population, }\AttributeTok{gen =} \DecValTok{2}\NormalTok{)}

\CommentTok{\# Calculate inbreeding}
\NormalTok{inbreeding }\OtherTok{\textless{}{-}} \FunctionTok{inbreeding.exp}\NormalTok{(population, }\AttributeTok{gen =} \DecValTok{2}\NormalTok{)}
\end{Highlighting}
\end{Shaded}

\subsection{4. Visualize Results}\label{visualize-results}

Use built-in plotting functions:

\begin{Shaded}
\begin{Highlighting}[]
\CommentTok{\# Track genetic gain}
\FunctionTok{bv.development}\NormalTok{(population)}

\CommentTok{\# PCA plot}
\FunctionTok{get.pca}\NormalTok{(population, }\AttributeTok{gen =} \DecValTok{1}\SpecialCharTok{:}\DecValTok{5}\NormalTok{)}

\CommentTok{\# Kinship development}
\FunctionTok{kinship.development}\NormalTok{(population, }\AttributeTok{gen =} \DecValTok{1}\SpecialCharTok{:}\DecValTok{5}\NormalTok{)}
\end{Highlighting}
\end{Shaded}

\section{The Population Object}\label{the-population-object}

The core of MoBPS is the \textbf{population object} (usually called
\texttt{population} or \texttt{pop}):

\begin{itemize}
\tightlist
\item
  It's an \textbf{R list} containing all simulation data
\item
  Stores genotypes, phenotypes, pedigrees, trait architectures
\item
  Updated by each \texttt{breeding.diploid()} call
\item
  Can be saved/loaded for reproducibility
\end{itemize}

Think of it as a \textbf{complete database} of your breeding program.

\section{Getting Help}\label{getting-help-1}

\subsection{Within R}\label{within-r}

\begin{Shaded}
\begin{Highlighting}[]
\CommentTok{\# Function documentation}
\NormalTok{?creating.diploid}
\NormalTok{?breeding.diploid}

\CommentTok{\# See all MoBPS functions}
\FunctionTok{help}\NormalTok{(}\AttributeTok{package =} \StringTok{"MoBPS"}\NormalTok{)}
\end{Highlighting}
\end{Shaded}

\subsection{Online Resources}\label{online-resources}

\begin{itemize}
\tightlist
\item
  \textbf{GitHub:} \url{https://github.com/tpook92/MoBPS}
\item
  \textbf{Web Interface:} \url{https://www.mobps.de}
\item
  \textbf{Email Support:} Torsten.pook@wur.nl
\end{itemize}

\subsection{Important Notes}\label{important-notes}

\begin{tcolorbox}[enhanced jigsaw, rightrule=.15mm, bottomrule=.15mm, bottomtitle=1mm, titlerule=0mm, colbacktitle=quarto-callout-warning-color!10!white, colframe=quarto-callout-warning-color-frame, opacityback=0, title=\textcolor{quarto-callout-warning-color}{\faExclamationTriangle}\hspace{0.5em}{Take Simulation Results with Caution!}, coltitle=black, left=2mm, leftrule=.75mm, breakable, toptitle=1mm, toprule=.15mm, opacitybacktitle=0.6, arc=.35mm, colback=white]

MoBPS won't stop you from making biologically impossible settings (e.g.,
collecting milk records from bulls). Always validate your simulation
setup makes biological sense!

\end{tcolorbox}

\begin{tcolorbox}[enhanced jigsaw, rightrule=.15mm, bottomrule=.15mm, bottomtitle=1mm, titlerule=0mm, colbacktitle=quarto-callout-tip-color!10!white, colframe=quarto-callout-tip-color-frame, opacityback=0, title=\textcolor{quarto-callout-tip-color}{\faLightbulb}\hspace{0.5em}{Don't Get Overwhelmed!}, coltitle=black, left=2mm, leftrule=.75mm, breakable, toptitle=1mm, toprule=.15mm, opacitybacktitle=0.6, arc=.35mm, colback=white]

MoBPS has many parameters, but you only need a few for most simulations.
Start simple, then add complexity as needed. Use Ctrl+F to search the
documentation for specific features.

\end{tcolorbox}

\section{What's Next?}\label{whats-next}

In the next chapter, we'll dive into \textbf{core concepts} that are
essential for understanding how MoBPS works:

\begin{itemize}
\tightlist
\item
  Individual grouping (generations, databases, cohorts)
\item
  The population list structure
\item
  How traits are represented
\item
  Time flow in simulations
\end{itemize}

Let's continue to \href{02-core-concepts.qmd}{Chapter 2: Core Concepts}!

\section{Summary}\label{summary}

\begin{itemize}
\tightlist
\item
  MoBPS is a flexible R package for breeding program simulation
\item
  Install from GitHub for latest features
\item
  Optional performance packages (\texttt{miraculix}) are highly
  recommended
\item
  Basic workflow: Create → Breed → Analyze → Visualize
\item
  The population object stores all your simulation data
\item
  Help is available through documentation and direct contact
\end{itemize}

\chapter{Core Concepts}\label{sec-core-concepts}

\section{Learning Objectives}\label{learning-objectives-1}

By the end of this chapter, you will understand:

\begin{itemize}
\tightlist
\item
  The gen/database/cohorts system for grouping individuals
\item
  How sex is handled in MoBPS
\item
  The structure of the population object
\item
  How time flows through generations
\item
  Key terminology and concepts
\end{itemize}

\section{Individual Grouping}\label{sec-grouping}

One of the biggest challenges in breeding simulation is having the
flexibility to perform operations on \textbf{specific groups} of
individuals. MoBPS provides three powerful ways to select groups:

\begin{enumerate}
\def\labelenumi{\arabic{enumi}.}
\tightlist
\item
  \textbf{Generations} (\texttt{gen}) - Select all individuals from
  specific generation(s)
\item
  \textbf{Cohorts} (\texttt{cohorts}) - Select named groups with
  specific characteristics
\item
  \textbf{Database} (\texttt{database}) - Precise selection by
  generation, sex, and individual range
\end{enumerate}

\subsection{\texorpdfstring{Understanding Generations
(\texttt{gen})}{Understanding Generations (gen)}}\label{understanding-generations-gen}

Every time you create offspring with \texttt{breeding.diploid()}, they
are assigned to a \textbf{new generation}. Generations are numbered
sequentially starting from 1.

\begin{Shaded}
\begin{Highlighting}[]
\CommentTok{\# Create founder population (generation 1)}
\NormalTok{population }\OtherTok{\textless{}{-}} \FunctionTok{creating.diploid}\NormalTok{(}\AttributeTok{nsnp =} \DecValTok{1000}\NormalTok{, }\AttributeTok{nindi =} \DecValTok{100}\NormalTok{)}

\CommentTok{\# Create generation 2}
\NormalTok{population }\OtherTok{\textless{}{-}} \FunctionTok{breeding.diploid}\NormalTok{(population,}
                               \AttributeTok{selection.size =} \FunctionTok{c}\NormalTok{(}\DecValTok{10}\NormalTok{, }\DecValTok{10}\NormalTok{),}
                               \AttributeTok{breeding.size =} \FunctionTok{c}\NormalTok{(}\DecValTok{50}\NormalTok{, }\DecValTok{50}\NormalTok{))}

\CommentTok{\# Select generation 2}
\NormalTok{bv\_gen2 }\OtherTok{\textless{}{-}} \FunctionTok{get.bv}\NormalTok{(population, }\AttributeTok{gen =} \DecValTok{2}\NormalTok{)}

\CommentTok{\# Select multiple generations}
\NormalTok{bv\_multi }\OtherTok{\textless{}{-}} \FunctionTok{get.bv}\NormalTok{(population, }\AttributeTok{gen =} \DecValTok{2}\SpecialCharTok{:}\DecValTok{5}\NormalTok{)  }\CommentTok{\# Generations 2, 3, 4, 5}
\end{Highlighting}
\end{Shaded}

\textbf{Key points:} - Generations track \textbf{when} individuals were
born - Useful for age-structured populations - Easy to analyze trends
over time

\subsection{Named Groups: Cohorts}\label{sec-cohorts}

\textbf{Cohorts} are named groups of individuals you define. They're
incredibly useful for tracking:

\begin{itemize}
\tightlist
\item
  Different selection lines (e.g., ``HighYield'', ``LowFat'')
\item
  Breeding groups (e.g., ``NucleusHerd'', ``CommercialLine'')
\item
  Treatment groups (e.g., ``Tested'', ``Controls'')
\end{itemize}

\begin{Shaded}
\begin{Highlighting}[]
\CommentTok{\# Create offspring and assign to named cohort}
\NormalTok{population }\OtherTok{\textless{}{-}} \FunctionTok{breeding.diploid}\NormalTok{(}
\NormalTok{  population,}
  \AttributeTok{selection.size =} \FunctionTok{c}\NormalTok{(}\DecValTok{10}\NormalTok{, }\DecValTok{10}\NormalTok{),}
  \AttributeTok{breeding.size =} \FunctionTok{c}\NormalTok{(}\DecValTok{50}\NormalTok{, }\DecValTok{50}\NormalTok{),}
  \AttributeTok{name.cohort =} \StringTok{"SelectedLine"}  \CommentTok{\# Give this group a name}
\NormalTok{)}

\CommentTok{\# Later, select only this cohort}
\NormalTok{bv\_selected }\OtherTok{\textless{}{-}} \FunctionTok{get.bv}\NormalTok{(population, }\AttributeTok{cohorts =} \StringTok{"SelectedLine"}\NormalTok{)}

\CommentTok{\# Select multiple cohorts}
\NormalTok{multi\_cohorts }\OtherTok{\textless{}{-}} \FunctionTok{get.bv}\NormalTok{(population,}
                         \AttributeTok{cohorts =} \FunctionTok{c}\NormalTok{(}\StringTok{"SelectedLine"}\NormalTok{, }\StringTok{"ControlLine"}\NormalTok{))}
\end{Highlighting}
\end{Shaded}

\textbf{Best practices:} - Use descriptive names: ``TopSires'',
``TestGroup1'', ``Line\_A'' - Keep naming consistent across simulations
- Use cohorts when generation number alone isn't enough

\subsection{Precise Selection: Database}\label{sec-database}

The \textbf{database} parameter gives you surgical precision. It's a
matrix where each row specifies:

\begin{enumerate}
\def\labelenumi{\arabic{enumi}.}
\tightlist
\item
  \textbf{Generation} number
\item
  \textbf{Sex} (1 = male, 2 = female)
\item
  \textbf{First individual} to include (optional)
\item
  \textbf{Last individual} to include (optional)
\end{enumerate}

\begin{Shaded}
\begin{Highlighting}[]
\CommentTok{\# Select males 1{-}20 from generation 3}
\NormalTok{database }\OtherTok{\textless{}{-}} \FunctionTok{matrix}\NormalTok{(}\FunctionTok{c}\NormalTok{(}\DecValTok{3}\NormalTok{, }\DecValTok{1}\NormalTok{, }\DecValTok{1}\NormalTok{, }\DecValTok{20}\NormalTok{), }\AttributeTok{ncol =} \DecValTok{4}\NormalTok{)}
\NormalTok{males }\OtherTok{\textless{}{-}} \FunctionTok{get.bv}\NormalTok{(population, }\AttributeTok{database =}\NormalTok{ database)}

\CommentTok{\# Select multiple groups}
\NormalTok{database }\OtherTok{\textless{}{-}} \FunctionTok{rbind}\NormalTok{(}
  \FunctionTok{c}\NormalTok{(}\DecValTok{3}\NormalTok{, }\DecValTok{1}\NormalTok{, }\DecValTok{1}\NormalTok{, }\DecValTok{20}\NormalTok{),    }\CommentTok{\# Males 1{-}20 from gen 3}
  \FunctionTok{c}\NormalTok{(}\DecValTok{4}\NormalTok{, }\DecValTok{2}\NormalTok{, }\DecValTok{5}\NormalTok{, }\DecValTok{15}\NormalTok{),    }\CommentTok{\# Females 5{-}15 from gen 4}
  \FunctionTok{c}\NormalTok{(}\DecValTok{5}\NormalTok{, }\DecValTok{1}\NormalTok{, }\ConstantTok{NA}\NormalTok{, }\ConstantTok{NA}\NormalTok{)    }\CommentTok{\# All males from gen 5}
\NormalTok{)}
\NormalTok{subset }\OtherTok{\textless{}{-}} \FunctionTok{get.bv}\NormalTok{(population, }\AttributeTok{database =}\NormalTok{ database)}
\end{Highlighting}
\end{Shaded}

\textbf{When to use database:} - Need specific individual ranges -
Complex selection criteria - Combining multiple precise selections

\subsection{Combining Selection
Methods}\label{combining-selection-methods}

You can mix and match these methods:

\begin{Shaded}
\begin{Highlighting}[]
\CommentTok{\# Select from generations 4{-}5, specific males from gen 3, AND a cohort}
\NormalTok{database }\OtherTok{\textless{}{-}} \FunctionTok{matrix}\NormalTok{(}\FunctionTok{c}\NormalTok{(}\DecValTok{3}\NormalTok{, }\DecValTok{1}\NormalTok{, }\DecValTok{21}\NormalTok{, }\DecValTok{50}\NormalTok{), }\AttributeTok{ncol =} \DecValTok{4}\NormalTok{)}

\NormalTok{bv }\OtherTok{\textless{}{-}} \FunctionTok{get.bv}\NormalTok{(population,}
             \AttributeTok{gen =} \DecValTok{4}\SpecialCharTok{:}\DecValTok{5}\NormalTok{,              }\CommentTok{\# All of gen 4 \& 5}
             \AttributeTok{database =}\NormalTok{ database,     }\CommentTok{\# Males 21{-}50 from gen 3}
             \AttributeTok{cohorts =} \StringTok{"Founders"}\NormalTok{)    }\CommentTok{\# Plus the Founders cohort}
\end{Highlighting}
\end{Shaded}

This gives you incredible flexibility to work with exactly the
individuals you need!

\section{Sex Handling}\label{sec-sex}

MoBPS has a flexible approach to sex:

\subsection{Traditional Two-Sex
Systems}\label{traditional-two-sex-systems}

By default, individuals are assigned male (1) or female (2):

\begin{Shaded}
\begin{Highlighting}[]
\CommentTok{\# Control sex ratio in founders}
\NormalTok{population }\OtherTok{\textless{}{-}} \FunctionTok{creating.diploid}\NormalTok{(}
  \AttributeTok{nsnp =} \DecValTok{1000}\NormalTok{,}
  \AttributeTok{nindi =} \DecValTok{100}\NormalTok{,}
  \AttributeTok{sex.quota =} \FloatTok{0.5}  \CommentTok{\# 50\% female}
\NormalTok{)}

\CommentTok{\# Or specify exactly}
\NormalTok{population }\OtherTok{\textless{}{-}} \FunctionTok{creating.diploid}\NormalTok{(}
  \AttributeTok{nsnp =} \DecValTok{1000}\NormalTok{,}
  \AttributeTok{nindi =} \DecValTok{100}\NormalTok{,}
  \AttributeTok{sex.s =} \FunctionTok{c}\NormalTok{(}\FunctionTok{rep}\NormalTok{(}\DecValTok{1}\NormalTok{, }\DecValTok{40}\NormalTok{), }\FunctionTok{rep}\NormalTok{(}\DecValTok{2}\NormalTok{, }\DecValTok{60}\NormalTok{))  }\CommentTok{\# 40 males, 60 females}
\NormalTok{)}
\end{Highlighting}
\end{Shaded}

\subsection{Flexible Sex Usage}\label{flexible-sex-usage}

\textbf{Important:} Sex assignments are \textbf{not binding} for
breeding operations!

\begin{itemize}
\tightlist
\item
  An individual stored as ``female'' can be used as father
\item
  Useful for plant breeding where sex may not be fixed
\item
  Useful for modeling hermaphrodites or aquaculture
\end{itemize}

\subsection{One-Sex Mode}\label{one-sex-mode}

For organisms without meaningful sex distinctions:

\begin{Shaded}
\begin{Highlighting}[]
\CommentTok{\# Deactivate two{-}sex system}
\NormalTok{population }\OtherTok{\textless{}{-}} \FunctionTok{creating.diploid}\NormalTok{(}
  \AttributeTok{nsnp =} \DecValTok{1000}\NormalTok{,}
  \AttributeTok{nindi =} \DecValTok{100}\NormalTok{,}
  \AttributeTok{one.sex.mode =} \ConstantTok{TRUE}  \CommentTok{\# All individuals in "sex 1"}
\NormalTok{)}
\end{Highlighting}
\end{Shaded}

This automatically adjusts \texttt{breeding.size},
\texttt{selection.size}, etc. to work with a single group.

\subsection{Using Sex as Structure}\label{using-sex-as-structure}

Even in plants, you can use ``sex'' to organize populations:

\begin{itemize}
\tightlist
\item
  Sex 1 = Gene pool A
\item
  Sex 2 = Gene pool B
\end{itemize}

This provides convenient structure for tracking different groups!

\section{The Population Object}\label{sec-population-object}

The population object is the \textbf{heart of MoBPS}. It's an R list
containing everything about your breeding program.

\subsection{Main Components}\label{main-components}

The population object has two major sections:

\begin{Shaded}
\begin{Highlighting}[]
\CommentTok{\# Examine structure}
\FunctionTok{names}\NormalTok{(population)}
\CommentTok{\# [1] "info"     "breeding"}

\CommentTok{\# General information}
\FunctionTok{names}\NormalTok{(population}\SpecialCharTok{$}\NormalTok{info)}
\CommentTok{\# QTL effects, genetic maps, trait info, etc.}

\CommentTok{\# Individual{-}level data}
\FunctionTok{names}\NormalTok{(population}\SpecialCharTok{$}\NormalTok{breeding)}
\CommentTok{\# Genotypes, phenotypes, pedigrees by generation}
\end{Highlighting}
\end{Shaded}

\subsection{What's Stored?}\label{whats-stored}

\textbf{In \texttt{\$info} (population-level):} - Genetic map
(chromosome structure, marker positions) - QTL effects and locations -
Trait names and architectures - Correlation structures - Fixed effects

\textbf{In \texttt{\$breeding} (individual-level):} -
Genotypes/haplotypes for each individual - Breeding values (true genetic
values) - Phenotypes (observed values) - Pedigree information -
Genotyping status - Cohort assignments

\subsection{Accessing the Population
Object}\label{accessing-the-population-object}

While you \emph{can} access elements directly, it's better to use
\texttt{get.*()} functions:

\begin{Shaded}
\begin{Highlighting}[]
\CommentTok{\# DON\textquotesingle{}T DO THIS (fragile, complex):}
\NormalTok{bv }\OtherTok{\textless{}{-}}\NormalTok{ population}\SpecialCharTok{$}\NormalTok{breeding[[}\DecValTok{2}\NormalTok{]][[}\DecValTok{1}\NormalTok{]][[}\DecValTok{6}\NormalTok{]]}

\CommentTok{\# DO THIS INSTEAD (clear, robust):}
\NormalTok{bv }\OtherTok{\textless{}{-}} \FunctionTok{get.bv}\NormalTok{(population, }\AttributeTok{gen =} \DecValTok{2}\NormalTok{)}
\end{Highlighting}
\end{Shaded}

\textbf{Available getter functions:} - \texttt{get.bv()} - Breeding
values - \texttt{get.pheno()} - Phenotypes - \texttt{get.geno()} -
Genotypes - \texttt{get.pedigree()} - Pedigree - \texttt{get.map()} -
Genetic map - Many more! (See \hyperref[sec-function-reference]{Chapter
17})

\section{Classes: Advanced Grouping}\label{sec-classes}

\textbf{Classes} provide another layer of organization:

\begin{Shaded}
\begin{Highlighting}[]
\CommentTok{\# Assign individuals to class 1}
\NormalTok{population }\OtherTok{\textless{}{-}} \FunctionTok{set.class}\NormalTok{(population,}
                        \AttributeTok{gen =} \DecValTok{2}\NormalTok{,}
                        \AttributeTok{class =} \DecValTok{1}\NormalTok{)}

\CommentTok{\# Only phenotype class 1}
\NormalTok{population }\OtherTok{\textless{}{-}} \FunctionTok{breeding.diploid}\NormalTok{(}
\NormalTok{  population,}
  \AttributeTok{phenotyping.gen =} \DecValTok{2}\NormalTok{,}
  \AttributeTok{phenotyping.class =} \DecValTok{1}  \CommentTok{\# Only these get phenotyped}
\NormalTok{)}
\end{Highlighting}
\end{Shaded}

\textbf{Special classes:} - Class 0 (default): Normal active individuals
- Class -1: Culled/dead animals (automatically assigned) - Class 1+:
User-defined groups

Use classes for: - Active vs.~culled animals - Test groups vs.~controls
- Multiple herds/flocks with different management

\section{Time and Generation Flow}\label{sec-time-flow}

Understanding how time works in MoBPS is crucial:

\subsection{Sequential Generations}\label{sequential-generations}

Each \texttt{breeding.diploid()} call creates the \textbf{next
generation}:

\begin{Shaded}
\begin{Highlighting}[]
\CommentTok{\# Gen 1: Founders}
\NormalTok{pop }\OtherTok{\textless{}{-}} \FunctionTok{creating.diploid}\NormalTok{(}\AttributeTok{nsnp =} \DecValTok{1000}\NormalTok{, }\AttributeTok{nindi =} \DecValTok{100}\NormalTok{)}

\CommentTok{\# Gen 2: First offspring}
\NormalTok{pop }\OtherTok{\textless{}{-}} \FunctionTok{breeding.diploid}\NormalTok{(pop,}
                        \AttributeTok{selection.size =} \FunctionTok{c}\NormalTok{(}\DecValTok{10}\NormalTok{, }\DecValTok{10}\NormalTok{),}
                        \AttributeTok{breeding.size =} \FunctionTok{c}\NormalTok{(}\DecValTok{50}\NormalTok{, }\DecValTok{50}\NormalTok{))}

\CommentTok{\# Gen 3: Second offspring generation}
\NormalTok{pop }\OtherTok{\textless{}{-}} \FunctionTok{breeding.diploid}\NormalTok{(pop,}
                        \AttributeTok{selection.size =} \FunctionTok{c}\NormalTok{(}\DecValTok{10}\NormalTok{, }\DecValTok{10}\NormalTok{),}
                        \AttributeTok{breeding.size =} \FunctionTok{c}\NormalTok{(}\DecValTok{50}\NormalTok{, }\DecValTok{50}\NormalTok{))}
\end{Highlighting}
\end{Shaded}

Generation numbers are \textbf{automatic} and \textbf{sequential}.

\subsection{Overlapping Generations}\label{overlapping-generations}

You can have overlapping generations by selecting parents from multiple
generations:

\begin{Shaded}
\begin{Highlighting}[]
\CommentTok{\# Use individuals from generation 2 AND 3 as parents}
\NormalTok{pop }\OtherTok{\textless{}{-}} \FunctionTok{breeding.diploid}\NormalTok{(pop,}
                        \AttributeTok{selection.m.gen =} \DecValTok{2}\NormalTok{,      }\CommentTok{\# Males from gen 2}
                        \AttributeTok{selection.f.gen =} \DecValTok{3}\NormalTok{,      }\CommentTok{\# Females from gen 3}
                        \AttributeTok{breeding.size =} \FunctionTok{c}\NormalTok{(}\DecValTok{50}\NormalTok{, }\DecValTok{50}\NormalTok{))}
\end{Highlighting}
\end{Shaded}

This creates \textbf{generation 4} from a mix of gen 2 and 3 parents.

\subsection{Age Structure}\label{age-structure}

MoBPS doesn't explicitly track age, but you can model it:

\begin{itemize}
\tightlist
\item
  Use \textbf{generations} as age cohorts
\item
  Use \textbf{cohorts} to track birth years
\item
  Select parents from specific generations to control age
\end{itemize}

See Section 6.20 in the full manual for detailed age structure examples.

\section{Genetic Pools and Crossbreeding}\label{sec-pools}

\textbf{Founder pools} track the origin of genome segments:

\begin{Shaded}
\begin{Highlighting}[]
\CommentTok{\# Create two founder populations}
\NormalTok{pop1 }\OtherTok{\textless{}{-}} \FunctionTok{creating.diploid}\NormalTok{(}\AttributeTok{nsnp =} \DecValTok{1000}\NormalTok{, }\AttributeTok{nindi =} \DecValTok{50}\NormalTok{,}
                         \AttributeTok{founder.pool =} \DecValTok{1}\NormalTok{)  }\CommentTok{\# Pool 1}

\NormalTok{pop2 }\OtherTok{\textless{}{-}} \FunctionTok{creating.diploid}\NormalTok{(}\AttributeTok{nsnp =} \DecValTok{1000}\NormalTok{, }\AttributeTok{nindi =} \DecValTok{50}\NormalTok{,}
                         \AttributeTok{founder.pool =} \DecValTok{2}\NormalTok{)  }\CommentTok{\# Pool 2}
\end{Highlighting}
\end{Shaded}

\textbf{Why use pools?} - Track breed composition in crossbreeding -
Assign breed-specific QTL effects - Analyze admixture and introgression
- Model heterosis and breed complementarity

You can later query which parts of the genome came from which pool using
\texttt{get.pool()}.

\section{Key Terminology Recap}\label{key-terminology-recap}

\begin{longtable}[]{@{}
  >{\raggedright\arraybackslash}p{(\linewidth - 2\tabcolsep) * \real{0.3333}}
  >{\raggedright\arraybackslash}p{(\linewidth - 2\tabcolsep) * \real{0.6667}}@{}}
\toprule\noalign{}
\begin{minipage}[b]{\linewidth}\raggedright
Term
\end{minipage} & \begin{minipage}[b]{\linewidth}\raggedright
Definition
\end{minipage} \\
\midrule\noalign{}
\endhead
\bottomrule\noalign{}
\endlastfoot
\textbf{Generation} & Time point when individuals were born
(sequential) \\
\textbf{Cohort} & Named group of individuals with similar
characteristics \\
\textbf{Database} & Precise selection by gen, sex, and individual
range \\
\textbf{Class} & Category for management actions (0=active, -1=culled,
1+=custom) \\
\textbf{Sex} & Male (1) or female (2), flexibly used \\
\textbf{Pool} & Founder population origin for tracking breed
composition \\
\textbf{Population object} & R list containing all simulation data \\
\end{longtable}

\section{Practical Tips}\label{practical-tips}

\begin{tcolorbox}[enhanced jigsaw, rightrule=.15mm, bottomrule=.15mm, bottomtitle=1mm, titlerule=0mm, colbacktitle=quarto-callout-tip-color!10!white, colframe=quarto-callout-tip-color-frame, opacityback=0, title=\textcolor{quarto-callout-tip-color}{\faLightbulb}\hspace{0.5em}{Start Simple}, coltitle=black, left=2mm, leftrule=.75mm, breakable, toptitle=1mm, toprule=.15mm, opacitybacktitle=0.6, arc=.35mm, colback=white]

Don't try to use all features at once! Start with just \texttt{gen} for
selecting individuals, then add cohorts and classes as needed.

\end{tcolorbox}

\begin{tcolorbox}[enhanced jigsaw, rightrule=.15mm, bottomrule=.15mm, bottomtitle=1mm, titlerule=0mm, colbacktitle=quarto-callout-tip-color!10!white, colframe=quarto-callout-tip-color-frame, opacityback=0, title=\textcolor{quarto-callout-tip-color}{\faLightbulb}\hspace{0.5em}{Naming Conventions}, coltitle=black, left=2mm, leftrule=.75mm, breakable, toptitle=1mm, toprule=.15mm, opacitybacktitle=0.6, arc=.35mm, colback=white]

Use consistent naming: - Cohorts: ``Line\_A'', ``Test\_2024'',
``HighYield'' - Variables: \texttt{pop} or \texttt{population} for the
population object - Clear generation references in comments

\end{tcolorbox}

\begin{tcolorbox}[enhanced jigsaw, rightrule=.15mm, bottomrule=.15mm, bottomtitle=1mm, titlerule=0mm, colbacktitle=quarto-callout-warning-color!10!white, colframe=quarto-callout-warning-color-frame, opacityback=0, title=\textcolor{quarto-callout-warning-color}{\faExclamationTriangle}\hspace{0.5em}{Don't Lose Your Population Object!}, coltitle=black, left=2mm, leftrule=.75mm, breakable, toptitle=1mm, toprule=.15mm, opacitybacktitle=0.6, arc=.35mm, colback=white]

Always save important population objects:

\begin{Shaded}
\begin{Highlighting}[]
\FunctionTok{saveRDS}\NormalTok{(population, }\StringTok{"my\_population\_gen10.rds"}\NormalTok{)}
\CommentTok{\# Later:}
\NormalTok{population }\OtherTok{\textless{}{-}} \FunctionTok{readRDS}\NormalTok{(}\StringTok{"my\_population\_gen10.rds"}\NormalTok{)}
\end{Highlighting}
\end{Shaded}

\end{tcolorbox}

\section{Summary}\label{summary-1}

\begin{itemize}
\tightlist
\item
  \textbf{Three selection systems:} gen (generations), cohorts (named
  groups), database (precise)
\item
  \textbf{Flexible sex:} Can be biological sex, gene pools, or
  organizational structure
\item
  \textbf{Population object:} Central data structure storing all
  simulation information
\item
  \textbf{Classes:} Additional grouping for management (0=active,
  -1=culled, custom)
\item
  \textbf{Generations flow sequentially:} Each
  \texttt{breeding.diploid()} creates the next generation
\item
  \textbf{Pools track origins:} Useful for crossbreeding and admixture
\end{itemize}

\section{What's Next?}\label{whats-next-1}

Now that you understand the core concepts, let's put them into practice!

In \href{03-first-simulation.qmd}{Chapter 3: Your First Simulation},
you'll create a complete breeding program from start to finish.

\chapter{Your First Simulation}\label{sec-first-simulation}

\section{Learning Objectives}\label{learning-objectives-2}

By the end of this chapter, you will be able to:

\begin{itemize}
\tightlist
\item
  Create a complete breeding program simulation from scratch
\item
  Phenotype and genotype individuals
\item
  Perform genomic prediction
\item
  Select superior individuals
\item
  Generate offspring from selected parents
\item
  Analyze genetic gain and inbreeding
\end{itemize}

\section{The Complete Workflow}\label{the-complete-workflow}

Let's simulate a realistic breeding program with:

\begin{itemize}
\tightlist
\item
  100 founder animals
\item
  Genomic selection based on GBLUP
\item
  10 generations of selection
\item
  Analysis of genetic gain and diversity
\end{itemize}

\section{Step 1: Load MoBPS}\label{step-1-load-mobps}

\begin{Shaded}
\begin{Highlighting}[]
\CommentTok{\# Load the package}
\FunctionTok{library}\NormalTok{(MoBPS)}

\CommentTok{\# Optional: Set seed for reproducibility}
\FunctionTok{set.seed}\NormalTok{(}\DecValTok{12345}\NormalTok{)}
\end{Highlighting}
\end{Shaded}

\section{Step 2: Create Founder Population}\label{sec-create-founders}

\begin{Shaded}
\begin{Highlighting}[]
\CommentTok{\# Create a founder population}
\NormalTok{population }\OtherTok{\textless{}{-}} \FunctionTok{creating.diploid}\NormalTok{(}
  \AttributeTok{nsnp =} \DecValTok{10000}\NormalTok{,                }\CommentTok{\# 10,000 SNP markers}
  \AttributeTok{nindi =} \DecValTok{100}\NormalTok{,                 }\CommentTok{\# 100 individuals (50 male, 50 female)}
  \AttributeTok{chr.nr =} \DecValTok{5}\NormalTok{,                  }\CommentTok{\# 5 chromosomes}
  \AttributeTok{chromosome.length =} \DecValTok{2}\NormalTok{,       }\CommentTok{\# 2 Morgan each (10M total)}
  \AttributeTok{n.additive =} \DecValTok{100}\NormalTok{,            }\CommentTok{\# 100 additive QTLs}
  \AttributeTok{n.dominant =} \DecValTok{20}\NormalTok{,             }\CommentTok{\# 20 dominant QTLs}
  \AttributeTok{share.genotyped =} \DecValTok{1}\NormalTok{,         }\CommentTok{\# 100\% genotyped}
  \AttributeTok{var.target =} \DecValTok{1}\NormalTok{,              }\CommentTok{\# Genetic variance = 1}
  \AttributeTok{name.cohort =} \StringTok{"Founders"}     \CommentTok{\# Name this cohort}
\NormalTok{)}
\end{Highlighting}
\end{Shaded}

\textbf{What did we just create?}

\begin{itemize}
\tightlist
\item
  A population with realistic genomic structure (5 chromosomes)
\item
  100 animals with 10,000 SNPs each
\item
  A complex trait with both additive and dominance effects
\item
  All individuals are genotyped (ready for genomic selection)
\end{itemize}

\section{Step 3: Generate Phenotypes}\label{sec-phenotype}

\begin{Shaded}
\begin{Highlighting}[]
\CommentTok{\# Phenotype all individuals in generation 1}
\NormalTok{population }\OtherTok{\textless{}{-}} \FunctionTok{breeding.diploid}\NormalTok{(}
\NormalTok{  population,}
  \AttributeTok{phenotyping =} \StringTok{"all"}\NormalTok{,         }\CommentTok{\# Phenotype everyone}
  \AttributeTok{phenotyping.gen =} \DecValTok{1}\NormalTok{,         }\CommentTok{\# From generation 1}
  \AttributeTok{heritability =} \FloatTok{0.5}           \CommentTok{\# h² = 0.5 (moderate heritability)}
\NormalTok{)}
\end{Highlighting}
\end{Shaded}

\textbf{What happened?}

\begin{itemize}
\tightlist
\item
  Phenotypes were simulated for all 100 founders
\item
  Residual variance was automatically calculated to achieve h² = 0.5
\item
  Phenotypes = True Breeding Value + Environmental Effect
\end{itemize}

\section{Step 4: Genomic Prediction}\label{sec-prediction}

\begin{Shaded}
\begin{Highlighting}[]
\CommentTok{\# Perform genomic prediction (GBLUP)}
\NormalTok{population }\OtherTok{\textless{}{-}} \FunctionTok{breeding.diploid}\NormalTok{(}
\NormalTok{  population,}
  \AttributeTok{bve =} \ConstantTok{TRUE}\NormalTok{,                  }\CommentTok{\# Calculate breeding value estimates}
  \AttributeTok{bve.gen =} \DecValTok{1}                  \CommentTok{\# Use generation 1 as training set}
\NormalTok{)}
\end{Highlighting}
\end{Shaded}

\textbf{What happened?}

\begin{itemize}
\tightlist
\item
  A genomic relationship matrix (G) was calculated
\item
  Breeding values were estimated using the internal GBLUP
\item
  Each individual now has an estimated breeding value (EBV)
\end{itemize}

\begin{tcolorbox}[enhanced jigsaw, rightrule=.15mm, bottomrule=.15mm, bottomtitle=1mm, titlerule=0mm, colbacktitle=quarto-callout-note-color!10!white, colframe=quarto-callout-note-color-frame, opacityback=0, title=\textcolor{quarto-callout-note-color}{\faInfo}\hspace{0.5em}{Internal BVE Method}, coltitle=black, left=2mm, leftrule=.75mm, breakable, toptitle=1mm, toprule=.15mm, opacitybacktitle=0.6, arc=.35mm, colback=white]

The \texttt{bve\ =\ TRUE} uses MoBPS's internal GBLUP assuming
\textbf{known heritability}. For more realistic scenarios, use external
software like \texttt{rrBLUP}, \texttt{BGLR}, or \texttt{BLUPF90}
(covered in \hyperref[sec-genomic-prediction]{Chapter 13}).

\end{tcolorbox}

\section{Step 5: Selection of Parents}\label{sec-selection}

\begin{Shaded}
\begin{Highlighting}[]
\CommentTok{\# Select top individuals as parents}
\NormalTok{population }\OtherTok{\textless{}{-}} \FunctionTok{breeding.diploid}\NormalTok{(}
\NormalTok{  population,}
  \AttributeTok{selection.size =} \FunctionTok{c}\NormalTok{(}\DecValTok{10}\NormalTok{, }\DecValTok{10}\NormalTok{),         }\CommentTok{\# 10 males, 10 females}
  \AttributeTok{selection.criteria =} \StringTok{"bve"}\NormalTok{,         }\CommentTok{\# Select on estimated BV}
  \AttributeTok{selection.m.cohorts =} \StringTok{"Founders"}\NormalTok{,   }\CommentTok{\# Males from Founders}
  \AttributeTok{selection.f.cohorts =} \StringTok{"Founders"}\NormalTok{,   }\CommentTok{\# Females from Founders}
  \AttributeTok{breeding.size =} \FunctionTok{c}\NormalTok{(}\DecValTok{50}\NormalTok{, }\DecValTok{50}\NormalTok{),          }\CommentTok{\# Generate 50M + 50F offspring}
  \AttributeTok{name.cohort =} \StringTok{"Generation\_2"}        \CommentTok{\# Name the offspring}
\NormalTok{)}
\end{Highlighting}
\end{Shaded}

\textbf{What happened?}

\begin{itemize}
\tightlist
\item
  The top 10 males (by EBV) were selected as sires
\item
  The top 10 females (by EBV) were selected as dams
\item
  100 offspring (50 male, 50 female) were created via random mating
\item
  Selection intensity: 10\% of males, 10\% of females
\end{itemize}

\section{Step 6: Repeat for Multiple Generations}\label{sec-repeat}

\begin{Shaded}
\begin{Highlighting}[]
\CommentTok{\# Simulate 10 generations}
\ControlFlowTok{for}\NormalTok{ (gen }\ControlFlowTok{in} \DecValTok{2}\SpecialCharTok{:}\DecValTok{10}\NormalTok{) \{}

  \CommentTok{\# Phenotype the current generation}
\NormalTok{  population }\OtherTok{\textless{}{-}} \FunctionTok{breeding.diploid}\NormalTok{(}
\NormalTok{    population,}
    \AttributeTok{phenotyping =} \StringTok{"all"}\NormalTok{,}
    \AttributeTok{phenotyping.gen =}\NormalTok{ gen,}
    \AttributeTok{heritability =} \FloatTok{0.5}
\NormalTok{  )}

  \CommentTok{\# Update genomic predictions using all data}
\NormalTok{  population }\OtherTok{\textless{}{-}} \FunctionTok{breeding.diploid}\NormalTok{(}
\NormalTok{    population,}
    \AttributeTok{bve =} \ConstantTok{TRUE}\NormalTok{,}
    \AttributeTok{bve.gen =} \DecValTok{1}\SpecialCharTok{:}\NormalTok{gen                   }\CommentTok{\# Cumulative training set}
\NormalTok{  )}

  \CommentTok{\# Select and breed}
\NormalTok{  population }\OtherTok{\textless{}{-}} \FunctionTok{breeding.diploid}\NormalTok{(}
\NormalTok{    population,}
    \AttributeTok{selection.size =} \FunctionTok{c}\NormalTok{(}\DecValTok{10}\NormalTok{, }\DecValTok{10}\NormalTok{),}
    \AttributeTok{selection.criteria =} \StringTok{"bve"}\NormalTok{,}
    \AttributeTok{selection.gen =}\NormalTok{ gen,               }\CommentTok{\# Select from current generation}
    \AttributeTok{breeding.size =} \FunctionTok{c}\NormalTok{(}\DecValTok{50}\NormalTok{, }\DecValTok{50}\NormalTok{),}
    \AttributeTok{name.cohort =} \FunctionTok{paste0}\NormalTok{(}\StringTok{"Generation\_"}\NormalTok{, gen}\SpecialCharTok{+}\DecValTok{1}\NormalTok{)}
\NormalTok{  )}
\NormalTok{\}}
\end{Highlighting}
\end{Shaded}

\textbf{What's happening in the loop?}

\begin{enumerate}
\def\labelenumi{\arabic{enumi}.}
\tightlist
\item
  Each generation is phenotyped
\item
  Predictions are updated with cumulative data (generations 1 to
  current)
\item
  Top 10 males and 10 females are selected
\item
  100 offspring are produced
\end{enumerate}

\section{Step 7: Analyze Results}\label{sec-analyze}

\subsection{Genetic Gain}\label{genetic-gain}

\begin{Shaded}
\begin{Highlighting}[]
\CommentTok{\# Plot breeding value trends}
\FunctionTok{bv.development}\NormalTok{(population)}
\end{Highlighting}
\end{Shaded}

This creates a plot showing: - Mean breeding value by generation -
Genetic gain over time - Standard deviation within generation

\subsection{Extract Breeding Values}\label{extract-breeding-values}

\begin{Shaded}
\begin{Highlighting}[]
\CommentTok{\# Get true breeding values for all generations}
\NormalTok{bv\_all }\OtherTok{\textless{}{-}} \FunctionTok{get.bv}\NormalTok{(population, }\AttributeTok{gen =} \DecValTok{1}\SpecialCharTok{:}\DecValTok{11}\NormalTok{)}

\CommentTok{\# Calculate mean BV per generation}
\NormalTok{mean\_bv }\OtherTok{\textless{}{-}} \FunctionTok{apply}\NormalTok{(bv\_all, }\DecValTok{2}\NormalTok{, mean)}
\NormalTok{mean\_bv}

\CommentTok{\# Calculate genetic gain from generation 1 to 11}
\NormalTok{genetic\_gain }\OtherTok{\textless{}{-}}\NormalTok{ mean\_bv[}\DecValTok{11}\NormalTok{] }\SpecialCharTok{{-}}\NormalTok{ mean\_bv[}\DecValTok{1}\NormalTok{]}
\FunctionTok{cat}\NormalTok{(}\StringTok{"Total genetic gain:"}\NormalTok{, }\FunctionTok{round}\NormalTok{(genetic\_gain, }\DecValTok{3}\NormalTok{), }\StringTok{"genetic SD}\SpecialCharTok{\textbackslash{}n}\StringTok{"}\NormalTok{)}
\end{Highlighting}
\end{Shaded}

\subsection{Inbreeding Analysis}\label{inbreeding-analysis}

\begin{Shaded}
\begin{Highlighting}[]
\CommentTok{\# Calculate genomic inbreeding for generation 11}
\NormalTok{inbreeding\_gen11 }\OtherTok{\textless{}{-}} \FunctionTok{inbreeding.emp}\NormalTok{(population, }\AttributeTok{gen =} \DecValTok{11}\NormalTok{)}
\FunctionTok{mean}\NormalTok{(inbreeding\_gen11)}

\CommentTok{\# Track inbreeding across generations}
\NormalTok{inbreeding\_trend }\OtherTok{\textless{}{-}} \FunctionTok{numeric}\NormalTok{(}\DecValTok{11}\NormalTok{)}
\ControlFlowTok{for}\NormalTok{ (gen }\ControlFlowTok{in} \DecValTok{1}\SpecialCharTok{:}\DecValTok{11}\NormalTok{) \{}
\NormalTok{  inbreeding\_trend[gen] }\OtherTok{\textless{}{-}} \FunctionTok{mean}\NormalTok{(}\FunctionTok{inbreeding.emp}\NormalTok{(population, }\AttributeTok{gen =}\NormalTok{ gen))}
\NormalTok{\}}

\CommentTok{\# Plot inbreeding trend}
\FunctionTok{plot}\NormalTok{(}\DecValTok{1}\SpecialCharTok{:}\DecValTok{11}\NormalTok{, inbreeding\_trend,}
     \AttributeTok{type =} \StringTok{"b"}\NormalTok{,}
     \AttributeTok{xlab =} \StringTok{"Generation"}\NormalTok{,}
     \AttributeTok{ylab =} \StringTok{"Mean Inbreeding Coefficient"}\NormalTok{,}
     \AttributeTok{main =} \StringTok{"Inbreeding Over Time"}\NormalTok{,}
     \AttributeTok{col =} \StringTok{"red"}\NormalTok{, }\AttributeTok{lwd =} \DecValTok{2}\NormalTok{)}
\end{Highlighting}
\end{Shaded}

\subsection{Kinship Between
Individuals}\label{kinship-between-individuals}

\begin{Shaded}
\begin{Highlighting}[]
\CommentTok{\# Calculate kinship matrix for generation 11}
\NormalTok{kinship\_matrix }\OtherTok{\textless{}{-}} \FunctionTok{kinship.emp}\NormalTok{(population, }\AttributeTok{gen =} \DecValTok{11}\NormalTok{)}

\CommentTok{\# Average kinship (off{-}diagonal elements)}
\NormalTok{n }\OtherTok{\textless{}{-}} \FunctionTok{nrow}\NormalTok{(kinship\_matrix)}
\NormalTok{avg\_kinship }\OtherTok{\textless{}{-}} \FunctionTok{sum}\NormalTok{(kinship\_matrix }\SpecialCharTok{{-}} \FunctionTok{diag}\NormalTok{(}\FunctionTok{diag}\NormalTok{(kinship\_matrix))) }\SpecialCharTok{/}\NormalTok{ (n }\SpecialCharTok{*}\NormalTok{ (n }\SpecialCharTok{{-}} \DecValTok{1}\NormalTok{))}
\FunctionTok{cat}\NormalTok{(}\StringTok{"Average kinship in gen 11:"}\NormalTok{, }\FunctionTok{round}\NormalTok{(avg\_kinship, }\DecValTok{4}\NormalTok{), }\StringTok{"}\SpecialCharTok{\textbackslash{}n}\StringTok{"}\NormalTok{)}
\end{Highlighting}
\end{Shaded}

\subsection{Population Structure (PCA)}\label{population-structure-pca}

\begin{Shaded}
\begin{Highlighting}[]
\CommentTok{\# PCA of generations 1, 5, and 11}
\FunctionTok{get.pca}\NormalTok{(population, }\AttributeTok{gen =} \FunctionTok{c}\NormalTok{(}\DecValTok{1}\NormalTok{, }\DecValTok{5}\NormalTok{, }\DecValTok{11}\NormalTok{))}
\end{Highlighting}
\end{Shaded}

This creates a PCA plot showing how the population has diverged from
founders.

\section{Complete Script}\label{complete-script}

Here's the entire simulation in one script:

\begin{Shaded}
\begin{Highlighting}[]
\FunctionTok{library}\NormalTok{(MoBPS)}
\FunctionTok{set.seed}\NormalTok{(}\DecValTok{12345}\NormalTok{)}

\CommentTok{\# 1. Create founders}
\NormalTok{population }\OtherTok{\textless{}{-}} \FunctionTok{creating.diploid}\NormalTok{(}
  \AttributeTok{nsnp =} \DecValTok{10000}\NormalTok{, }\AttributeTok{nindi =} \DecValTok{100}\NormalTok{, }\AttributeTok{chr.nr =} \DecValTok{5}\NormalTok{, }\AttributeTok{chromosome.length =} \DecValTok{2}\NormalTok{,}
  \AttributeTok{n.additive =} \DecValTok{100}\NormalTok{, }\AttributeTok{n.dominant =} \DecValTok{20}\NormalTok{, }\AttributeTok{share.genotyped =} \DecValTok{1}\NormalTok{,}
  \AttributeTok{var.target =} \DecValTok{1}\NormalTok{, }\AttributeTok{name.cohort =} \StringTok{"Founders"}
\NormalTok{)}

\CommentTok{\# 2. Phenotype founders}
\NormalTok{population }\OtherTok{\textless{}{-}} \FunctionTok{breeding.diploid}\NormalTok{(}
\NormalTok{  population, }\AttributeTok{phenotyping =} \StringTok{"all"}\NormalTok{, }\AttributeTok{phenotyping.gen =} \DecValTok{1}\NormalTok{, }\AttributeTok{heritability =} \FloatTok{0.5}
\NormalTok{)}

\CommentTok{\# 3. Initial genomic prediction}
\NormalTok{population }\OtherTok{\textless{}{-}} \FunctionTok{breeding.diploid}\NormalTok{(population, }\AttributeTok{bve =} \ConstantTok{TRUE}\NormalTok{, }\AttributeTok{bve.gen =} \DecValTok{1}\NormalTok{)}

\CommentTok{\# 4. Create first offspring generation}
\NormalTok{population }\OtherTok{\textless{}{-}} \FunctionTok{breeding.diploid}\NormalTok{(}
\NormalTok{  population, }\AttributeTok{selection.size =} \FunctionTok{c}\NormalTok{(}\DecValTok{10}\NormalTok{, }\DecValTok{10}\NormalTok{), }\AttributeTok{selection.criteria =} \StringTok{"bve"}\NormalTok{,}
  \AttributeTok{selection.m.cohorts =} \StringTok{"Founders"}\NormalTok{, }\AttributeTok{selection.f.cohorts =} \StringTok{"Founders"}\NormalTok{,}
  \AttributeTok{breeding.size =} \FunctionTok{c}\NormalTok{(}\DecValTok{50}\NormalTok{, }\DecValTok{50}\NormalTok{), }\AttributeTok{name.cohort =} \StringTok{"Generation\_2"}
\NormalTok{)}

\CommentTok{\# 5. Repeat for 9 more generations}
\ControlFlowTok{for}\NormalTok{ (gen }\ControlFlowTok{in} \DecValTok{2}\SpecialCharTok{:}\DecValTok{10}\NormalTok{) \{}
\NormalTok{  population }\OtherTok{\textless{}{-}} \FunctionTok{breeding.diploid}\NormalTok{(}
\NormalTok{    population, }\AttributeTok{phenotyping =} \StringTok{"all"}\NormalTok{, }\AttributeTok{phenotyping.gen =}\NormalTok{ gen, }\AttributeTok{heritability =} \FloatTok{0.5}
\NormalTok{  )}
\NormalTok{  population }\OtherTok{\textless{}{-}} \FunctionTok{breeding.diploid}\NormalTok{(population, }\AttributeTok{bve =} \ConstantTok{TRUE}\NormalTok{, }\AttributeTok{bve.gen =} \DecValTok{1}\SpecialCharTok{:}\NormalTok{gen)}
\NormalTok{  population }\OtherTok{\textless{}{-}} \FunctionTok{breeding.diploid}\NormalTok{(}
\NormalTok{    population, }\AttributeTok{selection.size =} \FunctionTok{c}\NormalTok{(}\DecValTok{10}\NormalTok{, }\DecValTok{10}\NormalTok{), }\AttributeTok{selection.criteria =} \StringTok{"bve"}\NormalTok{,}
    \AttributeTok{selection.gen =}\NormalTok{ gen, }\AttributeTok{breeding.size =} \FunctionTok{c}\NormalTok{(}\DecValTok{50}\NormalTok{, }\DecValTok{50}\NormalTok{),}
    \AttributeTok{name.cohort =} \FunctionTok{paste0}\NormalTok{(}\StringTok{"Generation\_"}\NormalTok{, gen}\SpecialCharTok{+}\DecValTok{1}\NormalTok{)}
\NormalTok{  )}
\NormalTok{\}}

\CommentTok{\# 6. Analyze results}
\FunctionTok{bv.development}\NormalTok{(population)}
\NormalTok{bv\_all }\OtherTok{\textless{}{-}} \FunctionTok{get.bv}\NormalTok{(population, }\AttributeTok{gen =} \DecValTok{1}\SpecialCharTok{:}\DecValTok{11}\NormalTok{)}
\NormalTok{mean\_bv }\OtherTok{\textless{}{-}} \FunctionTok{apply}\NormalTok{(bv\_all, }\DecValTok{2}\NormalTok{, mean)}
\FunctionTok{cat}\NormalTok{(}\StringTok{"Genetic gain:"}\NormalTok{, }\FunctionTok{round}\NormalTok{(mean\_bv[}\DecValTok{11}\NormalTok{] }\SpecialCharTok{{-}}\NormalTok{ mean\_bv[}\DecValTok{1}\NormalTok{], }\DecValTok{3}\NormalTok{), }\StringTok{"}\SpecialCharTok{\textbackslash{}n}\StringTok{"}\NormalTok{)}

\NormalTok{inbreeding\_gen11 }\OtherTok{\textless{}{-}} \FunctionTok{inbreeding.emp}\NormalTok{(population, }\AttributeTok{gen =} \DecValTok{11}\NormalTok{)}
\FunctionTok{cat}\NormalTok{(}\StringTok{"Mean inbreeding gen 11:"}\NormalTok{, }\FunctionTok{round}\NormalTok{(}\FunctionTok{mean}\NormalTok{(inbreeding\_gen11), }\DecValTok{4}\NormalTok{), }\StringTok{"}\SpecialCharTok{\textbackslash{}n}\StringTok{"}\NormalTok{)}
\end{Highlighting}
\end{Shaded}

\section{Understanding the Results}\label{understanding-the-results}

\subsection{Expected Outcomes}\label{expected-outcomes}

After 10 generations of selection with these parameters, you should see:

\begin{itemize}
\tightlist
\item
  \textbf{Genetic gain:} \textasciitilde2-4 genetic standard deviations
\item
  \textbf{Inbreeding:} \textasciitilde5-15\% (due to small effective
  population size)
\item
  \textbf{Selection response:} Decreasing per generation (plateauing)
\item
  \textbf{Prediction accuracy:} Improving with cumulative training data
\end{itemize}

\subsection{Why These Results?}\label{why-these-results}

\begin{enumerate}
\def\labelenumi{\arabic{enumi}.}
\tightlist
\item
  \textbf{Strong selection} (10\% selected) = high genetic gain
\item
  \textbf{Small Ne} (10 males) = rapid inbreeding accumulation
\item
  \textbf{Moderate h²} (0.5) = good response to selection
\item
  \textbf{Finite genome} = some exhaustion of genetic variance
\end{enumerate}

\section{Variations to Try}\label{variations-to-try}

Modify the simulation to explore different scenarios:

\subsection{Higher Selection
Intensity}\label{higher-selection-intensity}

\begin{Shaded}
\begin{Highlighting}[]
\CommentTok{\# More extreme selection}
\NormalTok{selection.size }\OtherTok{=} \FunctionTok{c}\NormalTok{(}\DecValTok{5}\NormalTok{, }\DecValTok{5}\NormalTok{)  }\CommentTok{\# Top 5\% of each sex}
\end{Highlighting}
\end{Shaded}

\subsection{Larger Population}\label{larger-population}

\begin{Shaded}
\begin{Highlighting}[]
\CommentTok{\# Reduce inbreeding}
\NormalTok{nindi }\OtherTok{=} \DecValTok{500}\NormalTok{,                }\CommentTok{\# 500 founders}
\NormalTok{selection.size }\OtherTok{=} \FunctionTok{c}\NormalTok{(}\DecValTok{25}\NormalTok{, }\DecValTok{25}\NormalTok{),  }\CommentTok{\# But still select top 10\%}
\NormalTok{breeding.size }\OtherTok{=} \FunctionTok{c}\NormalTok{(}\DecValTok{250}\NormalTok{, }\DecValTok{250}\NormalTok{)  }\CommentTok{\# Maintain pop size}
\end{Highlighting}
\end{Shaded}

\subsection{Multiple Traits}\label{multiple-traits}

\begin{Shaded}
\begin{Highlighting}[]
\CommentTok{\# Add a second trait}
\NormalTok{n.additive }\OtherTok{=} \FunctionTok{c}\NormalTok{(}\DecValTok{100}\NormalTok{, }\DecValTok{80}\NormalTok{),           }\CommentTok{\# Trait 1 \& 2}
\NormalTok{trait.cor }\OtherTok{=} \FunctionTok{matrix}\NormalTok{(}\FunctionTok{c}\NormalTok{(}\DecValTok{1}\NormalTok{, }\SpecialCharTok{{-}}\FloatTok{0.3}\NormalTok{,      }\CommentTok{\# Negative correlation}
                      \SpecialCharTok{{-}}\FloatTok{0.3}\NormalTok{, }\DecValTok{1}\NormalTok{), }\AttributeTok{nrow =} \DecValTok{2}\NormalTok{)}
\end{Highlighting}
\end{Shaded}

\subsection{Different Mating
Strategies}\label{different-mating-strategies}

\begin{Shaded}
\begin{Highlighting}[]
\CommentTok{\# Avoid inbreeding}
\NormalTok{avoid.mating.halfsib }\OtherTok{=} \ConstantTok{TRUE}\NormalTok{,}
\NormalTok{avoid.mating.fullsib }\OtherTok{=} \ConstantTok{TRUE}
\end{Highlighting}
\end{Shaded}

\section{Common Issues and Solutions}\label{common-issues-and-solutions}

\begin{tcolorbox}[enhanced jigsaw, rightrule=.15mm, bottomrule=.15mm, bottomtitle=1mm, titlerule=0mm, colbacktitle=quarto-callout-warning-color!10!white, colframe=quarto-callout-warning-color-frame, opacityback=0, title=\textcolor{quarto-callout-warning-color}{\faExclamationTriangle}\hspace{0.5em}{``Not enough males/females''}, coltitle=black, left=2mm, leftrule=.75mm, breakable, toptitle=1mm, toprule=.15mm, opacitybacktitle=0.6, arc=.35mm, colback=white]

\textbf{Problem:} Selection.size too large for available individuals.

\textbf{Solution:} Check population size:

\begin{Shaded}
\begin{Highlighting}[]
\FunctionTok{get.size}\NormalTok{(population, }\AttributeTok{gen =} \DecValTok{2}\NormalTok{)}
\end{Highlighting}
\end{Shaded}

Ensure \texttt{selection.size} doesn't exceed population size.

\end{tcolorbox}

\begin{tcolorbox}[enhanced jigsaw, rightrule=.15mm, bottomrule=.15mm, bottomtitle=1mm, titlerule=0mm, colbacktitle=quarto-callout-warning-color!10!white, colframe=quarto-callout-warning-color-frame, opacityback=0, title=\textcolor{quarto-callout-warning-color}{\faExclamationTriangle}\hspace{0.5em}{``No phenotypes available''}, coltitle=black, left=2mm, leftrule=.75mm, breakable, toptitle=1mm, toprule=.15mm, opacitybacktitle=0.6, arc=.35mm, colback=white]

\textbf{Problem:} Forgot to phenotype before BVE.

\textbf{Solution:} Always phenotype before prediction:

\begin{Shaded}
\begin{Highlighting}[]
\NormalTok{population }\OtherTok{\textless{}{-}} \FunctionTok{breeding.diploid}\NormalTok{(population,}
                                \AttributeTok{phenotyping =} \StringTok{"all"}\NormalTok{,}
                                \AttributeTok{heritability =} \FloatTok{0.5}\NormalTok{)}
\end{Highlighting}
\end{Shaded}

\end{tcolorbox}

\begin{tcolorbox}[enhanced jigsaw, rightrule=.15mm, bottomrule=.15mm, bottomtitle=1mm, titlerule=0mm, colbacktitle=quarto-callout-tip-color!10!white, colframe=quarto-callout-tip-color-frame, opacityback=0, title=\textcolor{quarto-callout-tip-color}{\faLightbulb}\hspace{0.5em}{Speeding Up Simulations}, coltitle=black, left=2mm, leftrule=.75mm, breakable, toptitle=1mm, toprule=.15mm, opacitybacktitle=0.6, arc=.35mm, colback=white]

For large-scale simulations: - Reduce \texttt{nsnp} (1000 SNPs often
sufficient) - Use \texttt{miraculix} package for fast matrix operations
- Set \texttt{copy.breeding.values\ =\ FALSE} to reduce memory

\end{tcolorbox}

\section{Summary}\label{summary-2}

In this chapter, you:

\begin{itemize}
\tightlist
\item
  ✅ Created a complete breeding program simulation
\item
  ✅ Implemented genomic selection with GBLUP
\item
  ✅ Tracked genetic gain over 10 generations
\item
  ✅ Analyzed inbreeding and population structure
\item
  ✅ Learned the core MoBPS workflow
\end{itemize}

\section{What's Next?}\label{whats-next-2}

Now that you can run a basic simulation, let's dive deeper into each
component:

\begin{itemize}
\tightlist
\item
  \textbf{Part 2:} Learn how to create more realistic founder
  populations with LD structure, import real data, and design complex
  trait architectures
\item
  \textbf{\hyperref[sec-creating-populations]{Chapter 4}:} Advanced
  population creation
\item
  \textbf{\hyperref[sec-trait-architecture]{Chapter 5}:} Sophisticated
  trait modeling
\end{itemize}

Continue to \href{04-creating-populations.qmd}{Chapter 4: Creating
Populations}!

\part{Building Populations \& Traits}

\chapter{Creating Founder Populations}\label{sec-creating-populations}

\section{Learning Objectives}\label{learning-objectives-3}

By the end of this chapter, you will be able to:

\begin{itemize}
\tightlist
\item
  Generate simulated genomic data with different patterns
\item
  Import real genomic data (VCF, PLINK formats)
\item
  Set up genetic maps and chromosome structures
\item
  Control allele frequencies and genetic diversity
\item
  Create multiple breeds/populations
\end{itemize}

\section{\texorpdfstring{Overview of
\texttt{creating.diploid()}}{Overview of creating.diploid()}}\label{overview-of-creating.diploid}

The \texttt{creating.diploid()} function is your starting point for
every MoBPS simulation. It creates the \textbf{founder population} with:

\begin{itemize}
\tightlist
\item
  Genomic data (markers/haplotypes)
\item
  Genetic map (chromosomes, positions)
\item
  Trait architecture (QTLs, effects)
\item
  Population structure (cohorts, sex, pools)
\end{itemize}

\section{Generating Simulated Data}\label{sec-simulated-data}

\subsection{Basic Random Population}\label{basic-random-population}

The simplest approach: generate random genotypes.

\begin{Shaded}
\begin{Highlighting}[]
\CommentTok{\# Create a simple founder population}
\NormalTok{population }\OtherTok{\textless{}{-}} \FunctionTok{creating.diploid}\NormalTok{(}
  \AttributeTok{nsnp =} \DecValTok{5000}\NormalTok{,          }\CommentTok{\# 5,000 SNP markers}
  \AttributeTok{nindi =} \DecValTok{200}\NormalTok{,          }\CommentTok{\# 200 individuals}
  \AttributeTok{n.additive =} \DecValTok{100}      \CommentTok{\# 100 QTLs}
\NormalTok{)}
\end{Highlighting}
\end{Shaded}

\textbf{Default behavior:} - Single chromosome (5 Morgan length) -
Random genotypes with allele frequency \textasciitilde{} Uniform(0, 1) -
50\% male, 50\% female - All individuals diploid and unrelated

\subsection{Controlling Allele
Frequencies}\label{controlling-allele-frequencies}

Allele frequencies dramatically affect genetic architecture:

\begin{Shaded}
\begin{Highlighting}[]
\CommentTok{\# Uniform allele frequencies (default)}
\NormalTok{pop\_uniform }\OtherTok{\textless{}{-}} \FunctionTok{creating.diploid}\NormalTok{(}
  \AttributeTok{nsnp =} \DecValTok{1000}\NormalTok{,}
  \AttributeTok{nindi =} \DecValTok{100}\NormalTok{,}
  \AttributeTok{beta\_shape1 =} \DecValTok{1}\NormalTok{,      }\CommentTok{\# Beta(1,1) = Uniform(0,1)}
  \AttributeTok{beta\_shape2 =} \DecValTok{1}
\NormalTok{)}

\CommentTok{\# Common variants (higher MAF)}
\NormalTok{pop\_common }\OtherTok{\textless{}{-}} \FunctionTok{creating.diploid}\NormalTok{(}
  \AttributeTok{nsnp =} \DecValTok{1000}\NormalTok{,}
  \AttributeTok{nindi =} \DecValTok{100}\NormalTok{,}
  \AttributeTok{beta\_shape1 =} \DecValTok{2}\NormalTok{,      }\CommentTok{\# Beta(2,2) concentrates near 0.5}
  \AttributeTok{beta\_shape2 =} \DecValTok{2}
\NormalTok{)}

\CommentTok{\# Rare variants (low MAF)}
\NormalTok{pop\_rare }\OtherTok{\textless{}{-}} \FunctionTok{creating.diploid}\NormalTok{(}
  \AttributeTok{nsnp =} \DecValTok{1000}\NormalTok{,}
  \AttributeTok{nindi =} \DecValTok{100}\NormalTok{,}
  \AttributeTok{beta\_shape1 =} \FloatTok{0.5}\NormalTok{,    }\CommentTok{\# Beta(0.5,0.5) concentrates near 0 and 1}
  \AttributeTok{beta\_shape2 =} \FloatTok{0.5}
\NormalTok{)}
\end{Highlighting}
\end{Shaded}

\begin{figure}[H]

{\centering \pandocbounded{\includegraphics[keepaspectratio]{index_files/mediabag/Beta_distribution.pdf}}

}

\caption{Allele frequency distributions from Beta(α, β)}

\end{figure}%

\begin{tcolorbox}[enhanced jigsaw, rightrule=.15mm, bottomrule=.15mm, bottomtitle=1mm, titlerule=0mm, colbacktitle=quarto-callout-tip-color!10!white, colframe=quarto-callout-tip-color-frame, opacityback=0, title=\textcolor{quarto-callout-tip-color}{\faLightbulb}\hspace{0.5em}{Realistic Allele Frequencies}, coltitle=black, left=2mm, leftrule=.75mm, breakable, toptitle=1mm, toprule=.15mm, opacitybacktitle=0.6, arc=.35mm, colback=white]

Real populations often have many rare variants and fewer common
variants. Use \texttt{beta\_shape1\ =\ 0.3,\ beta\_shape2\ =\ 1.5} to
approximate this.

\end{tcolorbox}

\subsection{Initialization Modes}\label{initialization-modes}

Control starting genotypes with \texttt{dataset} parameter:

\begin{Shaded}
\begin{Highlighting}[]
\CommentTok{\# Mode 1: All zeros (000.../000...)}
\NormalTok{pop\_zero }\OtherTok{\textless{}{-}} \FunctionTok{creating.diploid}\NormalTok{(}\AttributeTok{nsnp =} \DecValTok{1000}\NormalTok{, }\AttributeTok{nindi =} \DecValTok{100}\NormalTok{,}
                              \AttributeTok{dataset =} \StringTok{"all0"}\NormalTok{)}

\CommentTok{\# Mode 2: Fully heterozygous (000.../111...)}
\NormalTok{pop\_het }\OtherTok{\textless{}{-}} \FunctionTok{creating.diploid}\NormalTok{(}\AttributeTok{nsnp =} \DecValTok{1000}\NormalTok{, }\AttributeTok{nindi =} \DecValTok{100}\NormalTok{,}
                             \AttributeTok{dataset =} \StringTok{"allhetero"}\NormalTok{)}

\CommentTok{\# Mode 3: Random (default, X₁X₂X₃.../X₄X₅X₆...)}
\NormalTok{pop\_random }\OtherTok{\textless{}{-}} \FunctionTok{creating.diploid}\NormalTok{(}\AttributeTok{nsnp =} \DecValTok{1000}\NormalTok{, }\AttributeTok{nindi =} \DecValTok{100}\NormalTok{,}
                                \AttributeTok{dataset =} \StringTok{"random"}\NormalTok{)}

\CommentTok{\# Mode 4: Homozygous random (X₁X₂X₃.../X₁X₂X₃... same haplotypes)}
\NormalTok{pop\_homo }\OtherTok{\textless{}{-}} \FunctionTok{creating.diploid}\NormalTok{(}\AttributeTok{nsnp =} \DecValTok{1000}\NormalTok{, }\AttributeTok{nindi =} \DecValTok{100}\NormalTok{,}
                              \AttributeTok{dataset =} \StringTok{"homorandom"}\NormalTok{)}
\end{Highlighting}
\end{Shaded}

\textbf{When to use each mode:}

\begin{itemize}
\tightlist
\item
  \texttt{"random"} - General purpose, unrelated founders
\item
  \texttt{"homorandom"} - Create inbred lines or DH lines
\item
  \texttt{"allhetero"} - Maximum heterozygosity (F1s)
\item
  \texttt{"all0"} - Specific starting conditions or testing
\end{itemize}

\section{Chromosome Structure}\label{sec-chromosomes}

\subsection{Single vs.~Multiple
Chromosomes}\label{single-vs.-multiple-chromosomes}

\begin{Shaded}
\begin{Highlighting}[]
\CommentTok{\# Single chromosome (default)}
\NormalTok{pop\_single }\OtherTok{\textless{}{-}} \FunctionTok{creating.diploid}\NormalTok{(}
  \AttributeTok{nsnp =} \DecValTok{5000}\NormalTok{,}
  \AttributeTok{chromosome.length =} \DecValTok{5}  \CommentTok{\# 5 Morgan}
\NormalTok{)}

\CommentTok{\# Multiple chromosomes}
\NormalTok{pop\_multi }\OtherTok{\textless{}{-}} \FunctionTok{creating.diploid}\NormalTok{(}
  \AttributeTok{nsnp =} \DecValTok{10000}\NormalTok{,}
  \AttributeTok{chr.nr =} \FunctionTok{c}\NormalTok{(}\DecValTok{2000}\NormalTok{, }\DecValTok{2000}\NormalTok{, }\DecValTok{3000}\NormalTok{, }\DecValTok{3000}\NormalTok{),  }\CommentTok{\# SNPs per chromosome}
  \AttributeTok{chromosome.length =} \FunctionTok{c}\NormalTok{(}\DecValTok{2}\NormalTok{, }\DecValTok{2}\NormalTok{, }\FloatTok{1.5}\NormalTok{, }\FloatTok{1.5}\NormalTok{) }\CommentTok{\# Length in Morgan}
\NormalTok{)}
\end{Highlighting}
\end{Shaded}

\textbf{Consequences of chromosome structure:} - More chromosomes = more
recombination = faster LD decay - Longer chromosomes = stronger linkage
between distant markers - Realistic structure improves simulation
accuracy

\subsection{Using Template Species
Maps}\label{using-template-species-maps}

MoBPS includes common species maps:

\begin{Shaded}
\begin{Highlighting}[]
\CommentTok{\# Cattle map}
\NormalTok{pop\_cattle }\OtherTok{\textless{}{-}} \FunctionTok{creating.diploid}\NormalTok{(}
  \AttributeTok{nsnp =} \DecValTok{50000}\NormalTok{,}
  \AttributeTok{nindi =} \DecValTok{100}\NormalTok{,}
  \AttributeTok{template.chip =} \StringTok{"cattle"}   \CommentTok{\# 29 autosomes}
\NormalTok{)}

\CommentTok{\# Other available templates}
\NormalTok{template.chip }\OtherTok{=} \StringTok{"pig"}       \CommentTok{\# Sus scrofa}
\NormalTok{template.chip }\OtherTok{=} \StringTok{"chicken"}   \CommentTok{\# Gallus gallus}
\NormalTok{template.chip }\OtherTok{=} \StringTok{"sheep"}     \CommentTok{\# Ovis aries}
\NormalTok{template.chip }\OtherTok{=} \StringTok{"maize"}     \CommentTok{\# Zea mays}
\end{Highlighting}
\end{Shaded}

These provide realistic: - Number of chromosomes - Relative chromosome
lengths (in Morgan) - Approximate recombination rates

\begin{tcolorbox}[enhanced jigsaw, rightrule=.15mm, bottomrule=.15mm, bottomtitle=1mm, titlerule=0mm, colbacktitle=quarto-callout-note-color!10!white, colframe=quarto-callout-note-color-frame, opacityback=0, title=\textcolor{quarto-callout-note-color}{\faInfo}\hspace{0.5em}{What Templates DON'T Include}, coltitle=black, left=2mm, leftrule=.75mm, breakable, toptitle=1mm, toprule=.15mm, opacitybacktitle=0.6, arc=.35mm, colback=white]

Templates provide chromosome \textbf{structure only}, not: - Real marker
positions (SNPs are evenly spaced) - Real allele frequencies - Real LD
patterns

For these, import real data (see next section).

\end{tcolorbox}

\subsection{Custom Genetic Maps}\label{custom-genetic-maps}

Provide a full genetic map:

\begin{Shaded}
\begin{Highlighting}[]
\CommentTok{\# Create a custom map}
\CommentTok{\# Columns: chr, snp\_name, position\_Morgan, position\_bp, allele\_freq}
\NormalTok{my\_map }\OtherTok{\textless{}{-}} \FunctionTok{data.frame}\NormalTok{(}
  \AttributeTok{chr =} \FunctionTok{c}\NormalTok{(}\DecValTok{1}\NormalTok{, }\DecValTok{1}\NormalTok{, }\DecValTok{1}\NormalTok{, }\DecValTok{2}\NormalTok{, }\DecValTok{2}\NormalTok{, }\DecValTok{2}\NormalTok{),}
  \AttributeTok{snp =} \FunctionTok{c}\NormalTok{(}\StringTok{"rs001"}\NormalTok{, }\StringTok{"rs002"}\NormalTok{, }\StringTok{"rs003"}\NormalTok{, }\StringTok{"rs004"}\NormalTok{, }\StringTok{"rs005"}\NormalTok{, }\StringTok{"rs006"}\NormalTok{),}
  \AttributeTok{pos\_M =} \FunctionTok{c}\NormalTok{(}\FloatTok{0.0}\NormalTok{, }\FloatTok{0.5}\NormalTok{, }\FloatTok{1.0}\NormalTok{, }\FloatTok{0.0}\NormalTok{, }\FloatTok{0.3}\NormalTok{, }\FloatTok{0.6}\NormalTok{),}
  \AttributeTok{pos\_bp =} \FunctionTok{c}\NormalTok{(}\DecValTok{1000}\NormalTok{, }\DecValTok{50000000}\NormalTok{, }\DecValTok{100000000}\NormalTok{, }\DecValTok{1000}\NormalTok{, }\DecValTok{30000000}\NormalTok{, }\DecValTok{60000000}\NormalTok{),}
  \AttributeTok{freq =} \FunctionTok{c}\NormalTok{(}\FloatTok{0.2}\NormalTok{, }\FloatTok{0.5}\NormalTok{, }\FloatTok{0.8}\NormalTok{, }\FloatTok{0.1}\NormalTok{, }\FloatTok{0.45}\NormalTok{, }\FloatTok{0.7}\NormalTok{)}
\NormalTok{)}

\CommentTok{\# Use the map}
\NormalTok{population }\OtherTok{\textless{}{-}} \FunctionTok{creating.diploid}\NormalTok{(}
  \AttributeTok{map =}\NormalTok{ my\_map,}
  \AttributeTok{nindi =} \DecValTok{100}\NormalTok{,}
  \AttributeTok{n.additive =} \DecValTok{50}
\NormalTok{)}
\end{Highlighting}
\end{Shaded}

\section{Importing Real Data}\label{sec-import-real-data}

\subsection{From VCF Files}\label{from-vcf-files}

VCF (Variant Call Format) is standard for genomic data:

\begin{Shaded}
\begin{Highlighting}[]
\CommentTok{\# Import VCF file}
\NormalTok{population }\OtherTok{\textless{}{-}} \FunctionTok{creating.diploid}\NormalTok{(}
  \AttributeTok{vcf =} \StringTok{"path/to/genotypes.vcf"}\NormalTok{,  }\CommentTok{\# Or .vcf.gz}
  \AttributeTok{n.additive =} \DecValTok{100}\NormalTok{,}
  \AttributeTok{vcf.maxsnp =} \DecValTok{10000}\NormalTok{,             }\CommentTok{\# Optional: limit SNPs}
  \AttributeTok{vcf.maxindi =} \DecValTok{500}               \CommentTok{\# Optional: limit individuals}
\NormalTok{)}
\end{Highlighting}
\end{Shaded}

\textbf{What gets imported:} - Phased or unphased genotypes - Chromosome
numbers - Base pair positions - Marker names (rsIDs)

\subsection{From PLINK Files}\label{from-plink-files}

PLINK (.ped/.map or .bed/.bim/.fam) is also widely used:

\begin{Shaded}
\begin{Highlighting}[]
\CommentTok{\# Import .ped and .map files}
\NormalTok{ped\_data }\OtherTok{\textless{}{-}} \FunctionTok{read.pedmap}\NormalTok{(}\StringTok{"path/to/data.ped"}\NormalTok{, }\StringTok{"path/to/data.map"}\NormalTok{)}

\CommentTok{\# Create population from imported data}
\NormalTok{population }\OtherTok{\textless{}{-}} \FunctionTok{creating.diploid}\NormalTok{(}
  \AttributeTok{dataset =}\NormalTok{ ped\_data}\SpecialCharTok{$}\NormalTok{dataset,  }\CommentTok{\# Genotype matrix}
  \AttributeTok{map =}\NormalTok{ ped\_data}\SpecialCharTok{$}\NormalTok{map,           }\CommentTok{\# Genetic map}
  \AttributeTok{n.additive =} \DecValTok{100}
\NormalTok{)}
\end{Highlighting}
\end{Shaded}

\subsection{Converting Genotypes to
Haplotypes}\label{converting-genotypes-to-haplotypes}

Imported genotype data needs proper format:

\begin{Shaded}
\begin{Highlighting}[]
\CommentTok{\# If you have a genotype matrix (individuals × SNPs coded 0/1/2)}
\CommentTok{\# Convert to haplotype format for MoBPS}

\CommentTok{\# Example: genotype matrix}
\NormalTok{geno\_matrix }\OtherTok{\textless{}{-}} \FunctionTok{matrix}\NormalTok{(}\FunctionTok{sample}\NormalTok{(}\DecValTok{0}\SpecialCharTok{:}\DecValTok{2}\NormalTok{, }\DecValTok{1000}\NormalTok{, }\AttributeTok{replace =} \ConstantTok{TRUE}\NormalTok{),}
                      \AttributeTok{nrow =} \DecValTok{10}\NormalTok{, }\AttributeTok{ncol =} \DecValTok{100}\NormalTok{)  }\CommentTok{\# 10 indi, 100 SNPs}

\CommentTok{\# Convert: each individual becomes 2 haplotypes}
\NormalTok{haplo\_matrix }\OtherTok{\textless{}{-}} \FunctionTok{matrix}\NormalTok{(}\DecValTok{0}\NormalTok{, }\AttributeTok{nrow =} \FunctionTok{ncol}\NormalTok{(geno\_matrix), }\AttributeTok{ncol =} \DecValTok{2} \SpecialCharTok{*} \FunctionTok{nrow}\NormalTok{(geno\_matrix))}

\ControlFlowTok{for}\NormalTok{ (i }\ControlFlowTok{in} \DecValTok{1}\SpecialCharTok{:}\FunctionTok{nrow}\NormalTok{(geno\_matrix)) \{}
  \ControlFlowTok{for}\NormalTok{ (j }\ControlFlowTok{in} \DecValTok{1}\SpecialCharTok{:}\FunctionTok{ncol}\NormalTok{(geno\_matrix)) \{}
    \ControlFlowTok{if}\NormalTok{ (geno\_matrix[i,j] }\SpecialCharTok{==} \DecValTok{0}\NormalTok{) \{}
\NormalTok{      haplo\_matrix[j, }\DecValTok{2}\SpecialCharTok{*}\NormalTok{i}\DecValTok{{-}1}\NormalTok{] }\OtherTok{\textless{}{-}} \DecValTok{0}
\NormalTok{      haplo\_matrix[j, }\DecValTok{2}\SpecialCharTok{*}\NormalTok{i] }\OtherTok{\textless{}{-}} \DecValTok{0}
\NormalTok{    \} }\ControlFlowTok{else} \ControlFlowTok{if}\NormalTok{ (geno\_matrix[i,j] }\SpecialCharTok{==} \DecValTok{1}\NormalTok{) \{}
\NormalTok{      haplo\_matrix[j, }\DecValTok{2}\SpecialCharTok{*}\NormalTok{i}\DecValTok{{-}1}\NormalTok{] }\OtherTok{\textless{}{-}} \DecValTok{0}
\NormalTok{      haplo\_matrix[j, }\DecValTok{2}\SpecialCharTok{*}\NormalTok{i] }\OtherTok{\textless{}{-}} \DecValTok{1}
\NormalTok{    \} }\ControlFlowTok{else}\NormalTok{ \{  }\CommentTok{\# == 2}
\NormalTok{      haplo\_matrix[j, }\DecValTok{2}\SpecialCharTok{*}\NormalTok{i}\DecValTok{{-}1}\NormalTok{] }\OtherTok{\textless{}{-}} \DecValTok{1}
\NormalTok{      haplo\_matrix[j, }\DecValTok{2}\SpecialCharTok{*}\NormalTok{i] }\OtherTok{\textless{}{-}} \DecValTok{1}
\NormalTok{    \}}
\NormalTok{  \}}
\NormalTok{\}}

\CommentTok{\# Use in MoBPS}
\NormalTok{population }\OtherTok{\textless{}{-}} \FunctionTok{creating.diploid}\NormalTok{(}
  \AttributeTok{dataset =}\NormalTok{ haplo\_matrix,}
  \AttributeTok{n.additive =} \DecValTok{20}
\NormalTok{)}
\end{Highlighting}
\end{Shaded}

\section{Sex Ratio Control}\label{sec-sex-control}

\subsection{Controlling Sex
Proportions}\label{controlling-sex-proportions}

\begin{Shaded}
\begin{Highlighting}[]
\CommentTok{\# Equal sex ratio (default)}
\NormalTok{pop\_equal }\OtherTok{\textless{}{-}} \FunctionTok{creating.diploid}\NormalTok{(}\AttributeTok{nsnp =} \DecValTok{1000}\NormalTok{, }\AttributeTok{nindi =} \DecValTok{100}\NormalTok{,}
                               \AttributeTok{sex.quota =} \FloatTok{0.5}\NormalTok{)  }\CommentTok{\# 50\% female}

\CommentTok{\# More females}
\NormalTok{pop\_female }\OtherTok{\textless{}{-}} \FunctionTok{creating.diploid}\NormalTok{(}\AttributeTok{nsnp =} \DecValTok{1000}\NormalTok{, }\AttributeTok{nindi =} \DecValTok{100}\NormalTok{,}
                                \AttributeTok{sex.quota =} \FloatTok{0.7}\NormalTok{)  }\CommentTok{\# 70\% female}

\CommentTok{\# Specify exactly}
\NormalTok{pop\_exact }\OtherTok{\textless{}{-}} \FunctionTok{creating.diploid}\NormalTok{(}
  \AttributeTok{nsnp =} \DecValTok{1000}\NormalTok{,}
  \AttributeTok{nindi =} \DecValTok{100}\NormalTok{,}
  \AttributeTok{sex.s =} \FunctionTok{c}\NormalTok{(}\FunctionTok{rep}\NormalTok{(}\DecValTok{1}\NormalTok{, }\DecValTok{30}\NormalTok{), }\FunctionTok{rep}\NormalTok{(}\DecValTok{2}\NormalTok{, }\DecValTok{70}\NormalTok{))  }\CommentTok{\# 30 males, 70 females}
\NormalTok{)}
\end{Highlighting}
\end{Shaded}

\subsection{One-Sex Mode}\label{one-sex-mode-1}

For plants or situations where sex doesn't matter:

\begin{Shaded}
\begin{Highlighting}[]
\CommentTok{\# All individuals in same group}
\NormalTok{population }\OtherTok{\textless{}{-}} \FunctionTok{creating.diploid}\NormalTok{(}
  \AttributeTok{nsnp =} \DecValTok{1000}\NormalTok{,}
  \AttributeTok{nindi =} \DecValTok{200}\NormalTok{,}
  \AttributeTok{one.sex.mode =} \ConstantTok{TRUE}  \CommentTok{\# Deactivate two{-}sex system}
\NormalTok{)}
\end{Highlighting}
\end{Shaded}

\section{Creating Multiple
Breeds/Populations}\label{sec-multiple-breeds}

\subsection{Sequential Addition}\label{sequential-addition}

Create distinct founder populations:

\begin{Shaded}
\begin{Highlighting}[]
\CommentTok{\# Create breed 1}
\NormalTok{population }\OtherTok{\textless{}{-}} \FunctionTok{creating.diploid}\NormalTok{(}
  \AttributeTok{nsnp =} \DecValTok{5000}\NormalTok{,}
  \AttributeTok{nindi =} \DecValTok{100}\NormalTok{,}
  \AttributeTok{n.additive =} \DecValTok{100}\NormalTok{,}
  \AttributeTok{founder.pool =} \DecValTok{1}\NormalTok{,           }\CommentTok{\# Mark as pool 1}
  \AttributeTok{name.cohort =} \StringTok{"Breed\_A"}
\NormalTok{)}

\CommentTok{\# Add breed 2 (different allele frequencies)}
\NormalTok{population }\OtherTok{\textless{}{-}} \FunctionTok{creating.diploid}\NormalTok{(}
  \AttributeTok{population =}\NormalTok{ population,    }\CommentTok{\# Add to existing population}
  \AttributeTok{nsnp =} \DecValTok{5000}\NormalTok{,}
  \AttributeTok{nindi =} \DecValTok{100}\NormalTok{,}
  \AttributeTok{n.additive =} \DecValTok{100}\NormalTok{,}
  \AttributeTok{founder.pool =} \DecValTok{2}\NormalTok{,           }\CommentTok{\# Mark as pool 2}
  \AttributeTok{name.cohort =} \StringTok{"Breed\_B"}\NormalTok{,}
  \AttributeTok{freq =} \StringTok{"diff"}               \CommentTok{\# Different frequencies}
\NormalTok{)}
\end{Highlighting}
\end{Shaded}

\textbf{Uses for founder pools:} - Model crossbreeding programs - Track
breed composition (admixture) - Assign breed-specific QTL effects -
Study heterosis

\subsection{Adding Chromosomes}\label{adding-chromosomes}

Add additional chromosomes to existing population:

\begin{Shaded}
\begin{Highlighting}[]
\CommentTok{\# Start with chromosome 1}
\NormalTok{pop }\OtherTok{\textless{}{-}} \FunctionTok{creating.diploid}\NormalTok{(}\AttributeTok{nsnp =} \DecValTok{1000}\NormalTok{, }\AttributeTok{nindi =} \DecValTok{100}\NormalTok{)}

\CommentTok{\# Add chromosome 2}
\NormalTok{pop }\OtherTok{\textless{}{-}} \FunctionTok{creating.diploid}\NormalTok{(}
  \AttributeTok{population =}\NormalTok{ pop,}
  \AttributeTok{nsnp =} \DecValTok{1500}\NormalTok{,}
  \AttributeTok{add.chromosome =} \ConstantTok{TRUE}  \CommentTok{\# Add, don\textquotesingle{}t replace}
\NormalTok{)}
\end{Highlighting}
\end{Shaded}

\section{Marker Positions}\label{sec-marker-positions}

\subsection{Even Spacing (Default)}\label{even-spacing-default}

\begin{Shaded}
\begin{Highlighting}[]
\CommentTok{\# Markers evenly distributed}
\NormalTok{pop }\OtherTok{\textless{}{-}} \FunctionTok{creating.diploid}\NormalTok{(}
  \AttributeTok{nsnp =} \DecValTok{1000}\NormalTok{,}
  \AttributeTok{chromosome.length =} \DecValTok{5}  \CommentTok{\# Spread evenly over 5M}
\NormalTok{)}
\end{Highlighting}
\end{Shaded}

\textbf{Best for:} Fast computation, when exact positions don't matter.

\subsection{From Base Pairs}\label{from-base-pairs}

Convert physical positions to Morgan:

\begin{Shaded}
\begin{Highlighting}[]
\CommentTok{\# Provide base pair positions}
\NormalTok{bp\_positions }\OtherTok{\textless{}{-}} \FunctionTok{seq}\NormalTok{(}\DecValTok{1}\NormalTok{, }\DecValTok{100000000}\NormalTok{, }\AttributeTok{length.out =} \DecValTok{5000}\NormalTok{)  }\CommentTok{\# 100 Mb}

\NormalTok{pop }\OtherTok{\textless{}{-}} \FunctionTok{creating.diploid}\NormalTok{(}
  \AttributeTok{nsnp =} \DecValTok{5000}\NormalTok{,}
  \AttributeTok{bp =}\NormalTok{ bp\_positions,}
  \AttributeTok{bpcm.conversion =} \DecValTok{1000000}   \CommentTok{\# 1 Mb = 1 cM (typical for mammals)}
\NormalTok{)}
\end{Highlighting}
\end{Shaded}

\textbf{Common conversion rates:} - Mammals: 1,000,000 bp/cM (=
100,000,000 bp/Morgan) - Chicken: 300,000 bp/cM (= 30,000,000 bp/Morgan)
- Varies by species and chromosome!

\subsection{Directly in Morgan}\label{directly-in-morgan}

\begin{Shaded}
\begin{Highlighting}[]
\CommentTok{\# Provide positions in Morgan directly}
\NormalTok{positions\_M }\OtherTok{\textless{}{-}} \FunctionTok{c}\NormalTok{(}\FloatTok{0.0}\NormalTok{, }\FloatTok{0.01}\NormalTok{, }\FloatTok{0.05}\NormalTok{, }\FloatTok{0.1}\NormalTok{, }\FloatTok{0.15}\NormalTok{, ...)  }\CommentTok{\# Custom positions}

\NormalTok{pop }\OtherTok{\textless{}{-}} \FunctionTok{creating.diploid}\NormalTok{(}
  \AttributeTok{nsnp =} \DecValTok{1000}\NormalTok{,}
  \AttributeTok{snp.position =}\NormalTok{ positions\_M,}
  \AttributeTok{position.scaling =} \ConstantTok{FALSE}  \CommentTok{\# Don\textquotesingle{}t rescale}
\NormalTok{)}
\end{Highlighting}
\end{Shaded}

\section{Advanced Options}\label{sec-advanced-options}

\subsection{Genotyping Arrays}\label{genotyping-arrays}

Simulate partial genotyping (chip data):

\begin{Shaded}
\begin{Highlighting}[]
\CommentTok{\# Not all SNPs genotyped}
\NormalTok{pop }\OtherTok{\textless{}{-}} \FunctionTok{creating.diploid}\NormalTok{(}
  \AttributeTok{nsnp =} \DecValTok{50000}\NormalTok{,          }\CommentTok{\# 50K total SNPs}
  \AttributeTok{nindi =} \DecValTok{1000}\NormalTok{,}
  \AttributeTok{genotyped.s =} \FunctionTok{rep}\NormalTok{(}\FunctionTok{c}\NormalTok{(}\DecValTok{1}\NormalTok{, }\DecValTok{0}\NormalTok{, }\DecValTok{0}\NormalTok{, }\DecValTok{0}\NormalTok{), }\DecValTok{12500}\NormalTok{),  }\CommentTok{\# Every 4th SNP genotyped}
  \AttributeTok{share.genotyped =} \FloatTok{0.8}  \CommentTok{\# 80\% of individuals genotyped}
\NormalTok{)}
\end{Highlighting}
\end{Shaded}

\textbf{Uses:} - Model cost of genotyping - Test effects of marker
density - Simulate imputation scenarios

\subsection{Size Scaling}\label{size-scaling}

When founders are related (real data), scale effective size:

\begin{Shaded}
\begin{Highlighting}[]
\NormalTok{pop }\OtherTok{\textless{}{-}} \FunctionTok{creating.diploid}\NormalTok{(}
  \AttributeTok{vcf =} \StringTok{"real\_data.vcf"}\NormalTok{,}
  \AttributeTok{size.scaling =} \FloatTok{0.7}  \CommentTok{\# Effective size is 70\% of actual}
\NormalTok{)}
\end{Highlighting}
\end{Shaded}

This affects: - Calculations of expected inbreeding - Expected
relationships - Effective population size estimates

\section{Practical Examples}\label{sec-examples-pop}

\subsection{Example 1: Cattle
Population}\label{example-1-cattle-population}

\begin{Shaded}
\begin{Highlighting}[]
\CommentTok{\# Realistic cattle breeding population}
\NormalTok{cattle }\OtherTok{\textless{}{-}} \FunctionTok{creating.diploid}\NormalTok{(}
  \AttributeTok{nsnp =} \DecValTok{50000}\NormalTok{,}
  \AttributeTok{nindi =} \DecValTok{500}\NormalTok{,}
  \AttributeTok{template.chip =} \StringTok{"cattle"}\NormalTok{,}
  \AttributeTok{n.additive =} \DecValTok{100}\NormalTok{,}
  \AttributeTok{n.dominant =} \DecValTok{20}\NormalTok{,}
  \AttributeTok{beta\_shape1 =} \FloatTok{0.5}\NormalTok{,      }\CommentTok{\# Some rare variants}
  \AttributeTok{beta\_shape2 =} \FloatTok{1.2}\NormalTok{,}
  \AttributeTok{share.genotyped =} \FloatTok{0.8}\NormalTok{,   }\CommentTok{\# 80\% genotyped}
  \AttributeTok{sex.quota =} \FloatTok{0.7}\NormalTok{,         }\CommentTok{\# More females (dairy)}
  \AttributeTok{var.target =} \DecValTok{100}\NormalTok{,}
  \AttributeTok{name.cohort =} \StringTok{"HolsteinFounders"}
\NormalTok{)}
\end{Highlighting}
\end{Shaded}

\subsection{Example 2: Maize Inbred
Lines}\label{example-2-maize-inbred-lines}

\begin{Shaded}
\begin{Highlighting}[]
\CommentTok{\# Fully homozygous inbred lines}
\NormalTok{maize }\OtherTok{\textless{}{-}} \FunctionTok{creating.diploid}\NormalTok{(}
  \AttributeTok{nsnp =} \DecValTok{10000}\NormalTok{,}
  \AttributeTok{nindi =} \DecValTok{20}\NormalTok{,}
  \AttributeTok{template.chip =} \StringTok{"maize"}\NormalTok{,}
  \AttributeTok{dataset =} \StringTok{"homorandom"}\NormalTok{,  }\CommentTok{\# Homozygous}
  \AttributeTok{one.sex.mode =} \ConstantTok{TRUE}\NormalTok{,      }\CommentTok{\# No sexes}
  \AttributeTok{n.additive =} \DecValTok{200}\NormalTok{,}
  \AttributeTok{name.cohort =} \StringTok{"InbredLines"}
\NormalTok{)}
\end{Highlighting}
\end{Shaded}

\subsection{Example 3: Crossbreeding
Setup}\label{example-3-crossbreeding-setup}

\begin{Shaded}
\begin{Highlighting}[]
\CommentTok{\# Breed A}
\NormalTok{cross\_pop }\OtherTok{\textless{}{-}} \FunctionTok{creating.diploid}\NormalTok{(}
  \AttributeTok{nsnp =} \DecValTok{10000}\NormalTok{, }\AttributeTok{nindi =} \DecValTok{100}\NormalTok{, }\AttributeTok{n.additive =} \DecValTok{100}\NormalTok{,}
  \AttributeTok{founder.pool =} \DecValTok{1}\NormalTok{, }\AttributeTok{name.cohort =} \StringTok{"BreedA"}\NormalTok{,}
  \AttributeTok{beta\_shape1 =} \DecValTok{2}\NormalTok{, }\AttributeTok{beta\_shape2 =} \DecValTok{2}  \CommentTok{\# Common variants}
\NormalTok{)}

\CommentTok{\# Breed B (different frequencies)}
\NormalTok{cross\_pop }\OtherTok{\textless{}{-}} \FunctionTok{creating.diploid}\NormalTok{(}
  \AttributeTok{population =}\NormalTok{ cross\_pop,}
  \AttributeTok{nsnp =} \DecValTok{10000}\NormalTok{, }\AttributeTok{nindi =} \DecValTok{100}\NormalTok{, }\AttributeTok{n.additive =} \DecValTok{100}\NormalTok{,}
  \AttributeTok{founder.pool =} \DecValTok{2}\NormalTok{, }\AttributeTok{name.cohort =} \StringTok{"BreedB"}\NormalTok{,}
  \AttributeTok{beta\_shape1 =} \FloatTok{0.8}\NormalTok{, }\AttributeTok{beta\_shape2 =} \DecValTok{2}\NormalTok{,  }\CommentTok{\# Rare variants}
  \AttributeTok{freq =} \StringTok{"diff"}  \CommentTok{\# Independent frequencies from Breed A}
\NormalTok{)}
\end{Highlighting}
\end{Shaded}

\section{Summary}\label{summary-3}

Key concepts from this chapter:

\begin{itemize}
\tightlist
\item
  ✅ \texttt{creating.diploid()} initializes founder populations
\item
  ✅ Control genomic structure (chromosomes, positions, allele
  frequencies)
\item
  ✅ Import real data from VCF/PLINK or simulate data
\item
  ✅ Use species templates for realistic chromosome structure
\item
  ✅ Create multiple breeds with founder pools
\item
  ✅ Control sex ratios and population composition
\end{itemize}

\section{What's Next?}\label{whats-next-3}

Now that you can create populations, let's design the \textbf{trait
architecture} - the genetic basis of the traits you want to select on.

Continue to \href{05-trait-architecture.qmd}{Chapter 5: Trait
Architecture}!

\chapter{Trait Architecture}\label{sec-trait-architecture}

\section{Learning Objectives}\label{learning-objectives-4}

\begin{itemize}
\tightlist
\item
  Design simple and complex trait architectures
\item
  Model additive, dominance, and epistatic effects
\item
  Create correlated traits
\item
  Add maternal/paternal effects
\item
  Implement breed-specific QTLs
\end{itemize}

\section{Overview}\label{overview}

\textbf{Trait architecture} defines the genetic basis of traits: - Which
markers are QTLs - Effect sizes of each QTL - Types of genetic effects
(additive, dominance, epistasis) - Correlations between traits

MoBPS provides both \textbf{automated} and \textbf{custom} approaches.

\section{Automated Trait Generation}\label{sec-auto-traits}

\subsection{Simple Additive Trait}\label{simple-additive-trait}

\begin{Shaded}
\begin{Highlighting}[]
\NormalTok{population }\OtherTok{\textless{}{-}} \FunctionTok{creating.diploid}\NormalTok{(}
  \AttributeTok{nsnp =} \DecValTok{10000}\NormalTok{,}
  \AttributeTok{nindi =} \DecValTok{200}\NormalTok{,}
  \AttributeTok{n.additive =} \DecValTok{100}\NormalTok{,      }\CommentTok{\# 100 additive QTLs}
  \AttributeTok{var.target =} \DecValTok{1}\NormalTok{,        }\CommentTok{\# Genetic variance = 1}
  \AttributeTok{mean.target =} \DecValTok{100}      \CommentTok{\# Mean breeding value = 100}
\NormalTok{)}
\end{Highlighting}
\end{Shaded}

\textbf{What happens:} - 100 random SNPs chosen as QTLs - Effects drawn
from Normal(0, σ²) - Automatically scaled to target variance

\subsection{Adding Dominance}\label{adding-dominance}

\begin{Shaded}
\begin{Highlighting}[]
\NormalTok{population }\OtherTok{\textless{}{-}} \FunctionTok{creating.diploid}\NormalTok{(}
  \AttributeTok{nsnp =} \DecValTok{10000}\NormalTok{,}
  \AttributeTok{nindi =} \DecValTok{200}\NormalTok{,}
  \AttributeTok{n.additive =} \DecValTok{100}\NormalTok{,}
  \AttributeTok{n.dominant =} \DecValTok{20}\NormalTok{,       }\CommentTok{\# 20 loci with dominance}
  \AttributeTok{var.additive =} \FloatTok{0.8}\NormalTok{,    }\CommentTok{\# 80\% of variance from additive}
  \AttributeTok{var.dominant =} \FloatTok{0.2}     \CommentTok{\# 20\% from dominance}
\NormalTok{)}
\end{Highlighting}
\end{Shaded}

\textbf{Dominance effect modeling:} - Random loci chosen for dominance -
Effects: genotype 0, 1, 2 each get independent effect - Can specify
\texttt{dominant.only.positive\ =\ TRUE} for heterosis modeling

\subsection{Epistatic Effects}\label{epistatic-effects}

\begin{Shaded}
\begin{Highlighting}[]
\NormalTok{population }\OtherTok{\textless{}{-}} \FunctionTok{creating.diploid}\NormalTok{(}
  \AttributeTok{nsnp =} \DecValTok{10000}\NormalTok{,}
  \AttributeTok{nindi =} \DecValTok{200}\NormalTok{,}
  \AttributeTok{n.additive =} \DecValTok{80}\NormalTok{,}
  \AttributeTok{n.qualitative =} \DecValTok{10}\NormalTok{,    }\CommentTok{\# Qualitative epistasis (ordered)}
  \AttributeTok{n.quantitative =} \DecValTok{5}     \CommentTok{\# Quantitative epistasis (unordered)}
\NormalTok{)}
\end{Highlighting}
\end{Shaded}

\textbf{Epistasis types:}

\begin{enumerate}
\def\labelenumi{\arabic{enumi}.}
\tightlist
\item
  \textbf{Qualitative:} 00 \textless{} 01 \textless{} 02 \textless{} 10
  \textless{} 11 \textless{} 12 \textless{} 20 \textless{} 21
  \textless{} 22 (ordered, selection always beneficial)
\item
  \textbf{Quantitative:} Mix of positive/negative interactions
\end{enumerate}

\subsection{Effect Size Distributions}\label{effect-size-distributions}

\begin{Shaded}
\begin{Highlighting}[]
\CommentTok{\# Gaussian effects (default)}
\NormalTok{pop\_gaussian }\OtherTok{\textless{}{-}} \FunctionTok{creating.diploid}\NormalTok{(}
  \AttributeTok{nsnp =} \DecValTok{5000}\NormalTok{, }\AttributeTok{nindi =} \DecValTok{200}\NormalTok{, }\AttributeTok{n.additive =} \DecValTok{100}\NormalTok{,}
  \AttributeTok{effect.distribution =} \StringTok{"gaussian"}
\NormalTok{)}

\CommentTok{\# Gamma distribution (for skewed effects)}
\NormalTok{pop\_gamma }\OtherTok{\textless{}{-}} \FunctionTok{creating.diploid}\NormalTok{(}
  \AttributeTok{nsnp =} \DecValTok{5000}\NormalTok{, }\AttributeTok{nindi =} \DecValTok{200}\NormalTok{, }\AttributeTok{n.additive =} \DecValTok{100}\NormalTok{,}
  \AttributeTok{effect.distribution =} \StringTok{"gamma"}\NormalTok{,}
  \AttributeTok{gamma.shape1 =} \FloatTok{0.4}\NormalTok{,    }\CommentTok{\# Shape parameter}
  \AttributeTok{gamma.shape2 =} \FloatTok{1.66}    \CommentTok{\# Rate parameter}
\NormalTok{)}

\CommentTok{\# Equal effects}
\NormalTok{pop\_equal }\OtherTok{\textless{}{-}} \FunctionTok{creating.diploid}\NormalTok{(}
  \AttributeTok{nsnp =} \DecValTok{5000}\NormalTok{, }\AttributeTok{nindi =} \DecValTok{200}\NormalTok{,}
  \AttributeTok{n.equal.additive =} \DecValTok{50}\NormalTok{,           }\CommentTok{\# 50 QTLs}
  \AttributeTok{effect.size.equal.add =} \FloatTok{0.5}      \CommentTok{\# Each effect = 0.5}
\NormalTok{)}
\end{Highlighting}
\end{Shaded}

\section{Multiple Traits}\label{sec-multiple-traits}

\subsection{Independent Traits}\label{independent-traits}

\begin{Shaded}
\begin{Highlighting}[]
\NormalTok{population }\OtherTok{\textless{}{-}} \FunctionTok{creating.diploid}\NormalTok{(}
  \AttributeTok{nsnp =} \DecValTok{10000}\NormalTok{,}
  \AttributeTok{nindi =} \DecValTok{200}\NormalTok{,}
  \AttributeTok{n.additive =} \FunctionTok{c}\NormalTok{(}\DecValTok{100}\NormalTok{, }\DecValTok{80}\NormalTok{, }\DecValTok{120}\NormalTok{),    }\CommentTok{\# 3 traits}
  \AttributeTok{var.target =} \FunctionTok{c}\NormalTok{(}\DecValTok{1}\NormalTok{, }\FloatTok{1.5}\NormalTok{, }\FloatTok{0.8}\NormalTok{),     }\CommentTok{\# Different variances}
  \AttributeTok{trait.name =} \FunctionTok{c}\NormalTok{(}\StringTok{"Yield"}\NormalTok{, }\StringTok{"Quality"}\NormalTok{, }\StringTok{"Disease\_Resistance"}\NormalTok{)}
\NormalTok{)}
\end{Highlighting}
\end{Shaded}

\subsection{Correlated Traits}\label{correlated-traits}

\begin{Shaded}
\begin{Highlighting}[]
\CommentTok{\# Define correlation matrix}
\NormalTok{cor\_matrix }\OtherTok{\textless{}{-}} \FunctionTok{matrix}\NormalTok{(}\FunctionTok{c}\NormalTok{(}
  \FloatTok{1.0}\NormalTok{,  }\FloatTok{0.5}\NormalTok{, }\SpecialCharTok{{-}}\FloatTok{0.3}\NormalTok{,   }\CommentTok{\# Yield}
  \FloatTok{0.5}\NormalTok{,  }\FloatTok{1.0}\NormalTok{, }\SpecialCharTok{{-}}\FloatTok{0.2}\NormalTok{,   }\CommentTok{\# Quality}
 \SpecialCharTok{{-}}\FloatTok{0.3}\NormalTok{, }\SpecialCharTok{{-}}\FloatTok{0.2}\NormalTok{,  }\FloatTok{1.0}    \CommentTok{\# Disease Resistance}
\NormalTok{), }\AttributeTok{nrow =} \DecValTok{3}\NormalTok{, }\AttributeTok{byrow =} \ConstantTok{TRUE}\NormalTok{)}

\NormalTok{population }\OtherTok{\textless{}{-}} \FunctionTok{creating.diploid}\NormalTok{(}
  \AttributeTok{nsnp =} \DecValTok{10000}\NormalTok{,}
  \AttributeTok{nindi =} \DecValTok{200}\NormalTok{,}
  \AttributeTok{n.additive =} \FunctionTok{c}\NormalTok{(}\DecValTok{100}\NormalTok{, }\DecValTok{100}\NormalTok{, }\DecValTok{100}\NormalTok{),}
  \AttributeTok{trait.cor =}\NormalTok{ cor\_matrix,          }\CommentTok{\# Correlations}
  \AttributeTok{trait.name =} \FunctionTok{c}\NormalTok{(}\StringTok{"Yield"}\NormalTok{, }\StringTok{"Quality"}\NormalTok{, }\StringTok{"Disease\_Resistance"}\NormalTok{)}
\NormalTok{)}
\end{Highlighting}
\end{Shaded}

\textbf{How it works:} - Same markers used as QTLs for correlated traits
- Effects drawn from multivariate normal distribution - Correlations
independent of LD structure

\begin{tcolorbox}[enhanced jigsaw, rightrule=.15mm, bottomrule=.15mm, bottomtitle=1mm, titlerule=0mm, colbacktitle=quarto-callout-warning-color!10!white, colframe=quarto-callout-warning-color-frame, opacityback=0, title=\textcolor{quarto-callout-warning-color}{\faExclamationTriangle}\hspace{0.5em}{Correlation vs.~Covariance}, coltitle=black, left=2mm, leftrule=.75mm, breakable, toptitle=1mm, toprule=.15mm, opacitybacktitle=0.6, arc=.35mm, colback=white]

MoBPS uses \textbf{correlation} matrices, not covariance matrices.
Standardize your traits first!

\end{tcolorbox}

\section{Custom Trait Architecture}\label{sec-custom-traits}

\subsection{Single-Locus Effects
(real.bv.add)}\label{single-locus-effects-real.bv.add}

Specify exactly which SNPs are QTLs and their effects:

\begin{Shaded}
\begin{Highlighting}[]
\CommentTok{\# Create custom effects matrix}
\CommentTok{\# Columns: SNP\_number, chromosome, effect\_0, effect\_1, effect\_2}
\NormalTok{custom\_effects }\OtherTok{\textless{}{-}} \FunctionTok{rbind}\NormalTok{(}
  \FunctionTok{c}\NormalTok{(}\DecValTok{120}\NormalTok{,  }\DecValTok{1}\NormalTok{, }\SpecialCharTok{{-}}\FloatTok{1.0}\NormalTok{,  }\FloatTok{0.0}\NormalTok{,  }\FloatTok{1.0}\NormalTok{),   }\CommentTok{\# Additive effect}
  \FunctionTok{c}\NormalTok{(}\DecValTok{42}\NormalTok{,   }\DecValTok{5}\NormalTok{,  }\FloatTok{0.0}\NormalTok{,  }\FloatTok{0.0}\NormalTok{,  }\FloatTok{2.0}\NormalTok{),   }\CommentTok{\# Recessive}
  \FunctionTok{c}\NormalTok{(}\DecValTok{17}\NormalTok{,  }\DecValTok{22}\NormalTok{,  }\FloatTok{0.1}\NormalTok{,  }\FloatTok{0.1}\NormalTok{,  }\FloatTok{0.1}\NormalTok{)    }\CommentTok{\# No effect (neutral)}
\NormalTok{)}

\NormalTok{population }\OtherTok{\textless{}{-}} \FunctionTok{creating.diploid}\NormalTok{(}
  \AttributeTok{nsnp =} \DecValTok{10000}\NormalTok{,}
  \AttributeTok{nindi =} \DecValTok{200}\NormalTok{,}
  \AttributeTok{real.bv.add =}\NormalTok{ custom\_effects}
\NormalTok{)}
\end{Highlighting}
\end{Shaded}

\textbf{Effect coding:} - Column 1: SNP number (1 to nsnp) - Column 2:
Chromosome number - Columns 3-5: Effects for genotypes 0, 1, 2

\subsection{Two-Locus Epistasis
(real.bv.mult)}\label{two-locus-epistasis-real.bv.mult}

\begin{Shaded}
\begin{Highlighting}[]
\CommentTok{\# Epistatic interaction between two loci}
\CommentTok{\# 9 effects for all genotype combinations}
\NormalTok{epistatic\_effects }\OtherTok{\textless{}{-}} \FunctionTok{rbind}\NormalTok{(}
  \FunctionTok{c}\NormalTok{(}\DecValTok{144}\NormalTok{, }\DecValTok{1}\NormalTok{, }\DecValTok{145}\NormalTok{, }\DecValTok{1}\NormalTok{,  }\CommentTok{\# SNPs 144 and 145 on chr 1}
    \FloatTok{1.0}\NormalTok{, }\FloatTok{0.0}\NormalTok{, }\FloatTok{0.0}\NormalTok{,   }\CommentTok{\# Effects when SNP1 = 0}
    \FloatTok{0.0}\NormalTok{, }\FloatTok{0.0}\NormalTok{, }\FloatTok{0.0}\NormalTok{,   }\CommentTok{\# Effects when SNP1 = 1}
    \FloatTok{0.0}\NormalTok{, }\FloatTok{0.0}\NormalTok{, }\FloatTok{0.0}\NormalTok{),  }\CommentTok{\# Effects when SNP1 = 2}
  \FunctionTok{c}\NormalTok{(}\DecValTok{4}\NormalTok{, }\DecValTok{3}\NormalTok{, }\DecValTok{188}\NormalTok{, }\DecValTok{5}\NormalTok{,    }\CommentTok{\# SNP 4 on chr 3, SNP 188 on chr 5}
    \FloatTok{0.37}\NormalTok{, }\FloatTok{0.16}\NormalTok{, }\FloatTok{1.35}\NormalTok{,}
    \FloatTok{1.99}\NormalTok{, }\FloatTok{1.58}\NormalTok{, }\FloatTok{2.51}\NormalTok{,}
    \FloatTok{0.38}\NormalTok{, }\FloatTok{2.12}\NormalTok{, }\FloatTok{0.98}\NormalTok{)}
\NormalTok{)}

\NormalTok{population }\OtherTok{\textless{}{-}} \FunctionTok{creating.diploid}\NormalTok{(}
  \AttributeTok{nsnp =} \DecValTok{10000}\NormalTok{,}
  \AttributeTok{nindi =} \DecValTok{200}\NormalTok{,}
  \AttributeTok{real.bv.mult =}\NormalTok{ epistatic\_effects}
\NormalTok{)}
\end{Highlighting}
\end{Shaded}

\subsection{Multi-Locus Epistasis
(real.bv.dice)}\label{multi-locus-epistasis-real.bv.dice}

For effects involving 3+ loci:

\begin{Shaded}
\begin{Highlighting}[]
\CommentTok{\# Complex epistatic network}
\NormalTok{dice\_list }\OtherTok{\textless{}{-}} \FunctionTok{list}\NormalTok{(}
  \CommentTok{\# First network: 3 SNPs}
  \AttributeTok{location =} \FunctionTok{list}\NormalTok{(}
    \FunctionTok{matrix}\NormalTok{(}\FunctionTok{c}\NormalTok{(}\DecValTok{11}\NormalTok{, }\DecValTok{1}\NormalTok{,}
             \DecValTok{12}\NormalTok{, }\DecValTok{1}\NormalTok{,}
             \DecValTok{16}\NormalTok{, }\DecValTok{2}\NormalTok{), }\AttributeTok{ncol =} \DecValTok{2}\NormalTok{, }\AttributeTok{byrow =} \ConstantTok{TRUE}\NormalTok{)}
\NormalTok{  ),}
  \CommentTok{\# Effects for all 27 combinations (3\^{}3)}
  \AttributeTok{effects =} \FunctionTok{list}\NormalTok{(}
    \FunctionTok{runif}\NormalTok{(}\DecValTok{27}\NormalTok{, }\SpecialCharTok{{-}}\DecValTok{1}\NormalTok{, }\DecValTok{1}\NormalTok{)  }\CommentTok{\# Random effects}
\NormalTok{  )}
\NormalTok{)}

\NormalTok{population }\OtherTok{\textless{}{-}} \FunctionTok{creating.diploid}\NormalTok{(}
  \AttributeTok{nsnp =} \DecValTok{10000}\NormalTok{,}
  \AttributeTok{nindi =} \DecValTok{200}\NormalTok{,}
  \AttributeTok{real.bv.dice =}\NormalTok{ dice\_list}
\NormalTok{)}
\end{Highlighting}
\end{Shaded}

\section{Special Trait Types}\label{sec-special-traits}

\subsection{Maternal/Paternal Effects}\label{maternalpaternal-effects}

\begin{Shaded}
\begin{Highlighting}[]
\CommentTok{\# Trait influenced by mother\textquotesingle{}s genotype}
\NormalTok{population }\OtherTok{\textless{}{-}} \FunctionTok{creating.diploid}\NormalTok{(}
  \AttributeTok{nsnp =} \DecValTok{5000}\NormalTok{,}
  \AttributeTok{nindi =} \DecValTok{100}\NormalTok{,}
  \AttributeTok{n.additive =} \FunctionTok{c}\NormalTok{(}\DecValTok{100}\NormalTok{, }\DecValTok{50}\NormalTok{),}
  \AttributeTok{is.maternal =} \FunctionTok{c}\NormalTok{(}\ConstantTok{FALSE}\NormalTok{, }\ConstantTok{TRUE}\NormalTok{),  }\CommentTok{\# Trait 2 is maternal}
  \AttributeTok{trait.name =} \FunctionTok{c}\NormalTok{(}\StringTok{"Direct"}\NormalTok{, }\StringTok{"Maternal"}\NormalTok{)}
\NormalTok{)}

\CommentTok{\# Combine into single trait}
\NormalTok{population }\OtherTok{\textless{}{-}} \FunctionTok{add.combi}\NormalTok{(}
\NormalTok{  population,}
  \AttributeTok{trait =} \DecValTok{3}\NormalTok{,}
  \AttributeTok{combi.weights =} \FunctionTok{c}\NormalTok{(}\FloatTok{0.7}\NormalTok{, }\FloatTok{0.3}\NormalTok{),  }\CommentTok{\# 70\% direct, 30\% maternal}
  \AttributeTok{trait.name =} \StringTok{"Combined"}
\NormalTok{)}
\end{Highlighting}
\end{Shaded}

\textbf{Use cases:} - Birth weight (maternal nutrition) - Milk
production (maternal effects on calves) - Egg quality (maternal
contribution)

\subsection{Breed-Specific QTLs}\label{breed-specific-qtls}

Model QTLs that only have effects in specific breeds:

\begin{Shaded}
\begin{Highlighting}[]
\CommentTok{\# Create two breeds}
\NormalTok{pop }\OtherTok{\textless{}{-}} \FunctionTok{creating.diploid}\NormalTok{(}
  \AttributeTok{nsnp =} \DecValTok{5000}\NormalTok{, }\AttributeTok{nindi =} \DecValTok{100}\NormalTok{,}
  \AttributeTok{founder.pool =} \DecValTok{1}\NormalTok{,}
  \AttributeTok{name.cohort =} \StringTok{"Breed\_A"}
\NormalTok{)}

\NormalTok{pop }\OtherTok{\textless{}{-}} \FunctionTok{creating.diploid}\NormalTok{(}
  \AttributeTok{population =}\NormalTok{ pop, }\AttributeTok{nsnp =} \DecValTok{5000}\NormalTok{, }\AttributeTok{nindi =} \DecValTok{100}\NormalTok{,}
  \AttributeTok{founder.pool =} \DecValTok{2}\NormalTok{,}
  \AttributeTok{name.cohort =} \StringTok{"Breed\_B"}
\NormalTok{)}

\CommentTok{\# Add trait only active in breed A}
\NormalTok{pop }\OtherTok{\textless{}{-}} \FunctionTok{creating.trait}\NormalTok{(}
\NormalTok{  pop,}
  \AttributeTok{n.additive =} \DecValTok{100}\NormalTok{,}
  \AttributeTok{trait.pool =} \DecValTok{1}  \CommentTok{\# Only affects pool 1 (Breed A)}
\NormalTok{)}
\end{Highlighting}
\end{Shaded}

\subsection{Non-QTL Polygenic Traits}\label{non-qtl-polygenic-traits}

Simulate infinitesimal model:

\begin{Shaded}
\begin{Highlighting}[]
\NormalTok{population }\OtherTok{\textless{}{-}} \FunctionTok{creating.diploid}\NormalTok{(}
  \AttributeTok{nsnp =} \DecValTok{10000}\NormalTok{,}
  \AttributeTok{nindi =} \DecValTok{200}\NormalTok{,}
  \AttributeTok{n.traits =} \DecValTok{2}\NormalTok{,              }\CommentTok{\# Total traits}
  \AttributeTok{n.additive =} \DecValTok{100}\NormalTok{,          }\CommentTok{\# Trait 1 has QTLs}
  \AttributeTok{polygenic.variance =} \DecValTok{1}     \CommentTok{\# Trait 2 is polygenic}
\NormalTok{)}
\end{Highlighting}
\end{Shaded}

\textbf{Trait 2 characteristics:} - Based on parental average +
inbreeding - No specific QTLs - Useful for: very highly polygenic traits

\section{Adding Traits to Existing Populations}\label{sec-add-traits}

Use \texttt{creating.trait()} to add traits after population creation:

\begin{Shaded}
\begin{Highlighting}[]
\CommentTok{\# Create population}
\NormalTok{pop }\OtherTok{\textless{}{-}} \FunctionTok{creating.diploid}\NormalTok{(}\AttributeTok{nsnp =} \DecValTok{10000}\NormalTok{, }\AttributeTok{nindi =} \DecValTok{200}\NormalTok{)}

\CommentTok{\# Later, add a new trait}
\NormalTok{pop }\OtherTok{\textless{}{-}} \FunctionTok{creating.trait}\NormalTok{(}
\NormalTok{  pop,}
  \AttributeTok{n.additive =} \DecValTok{80}\NormalTok{,}
  \AttributeTok{n.dominant =} \DecValTok{20}\NormalTok{,}
  \AttributeTok{trait.name =} \StringTok{"NewTrait"}
\NormalTok{)}
\end{Highlighting}
\end{Shaded}

\textbf{Why add traits later?} - Different traits need different
architectures - Want to replace/update trait definitions - Building
complex multi-trait models step-by-step

\section{Combining Traits}\label{sec-combine-traits}

Create composite traits from existing traits:

\begin{Shaded}
\begin{Highlighting}[]
\CommentTok{\# Create base traits}
\NormalTok{pop }\OtherTok{\textless{}{-}} \FunctionTok{creating.diploid}\NormalTok{(}
  \AttributeTok{nsnp =} \DecValTok{5000}\NormalTok{, }\AttributeTok{nindi =} \DecValTok{200}\NormalTok{,}
  \AttributeTok{n.additive =} \FunctionTok{c}\NormalTok{(}\DecValTok{100}\NormalTok{, }\DecValTok{80}\NormalTok{, }\DecValTok{60}\NormalTok{),}
  \AttributeTok{trait.name =} \FunctionTok{c}\NormalTok{(}\StringTok{"Milk"}\NormalTok{, }\StringTok{"Fat"}\NormalTok{, }\StringTok{"Protein"}\NormalTok{)}
\NormalTok{)}

\CommentTok{\# Create selection index}
\NormalTok{pop }\OtherTok{\textless{}{-}} \FunctionTok{add.combi}\NormalTok{(}
\NormalTok{  pop,}
  \AttributeTok{trait =} \DecValTok{4}\NormalTok{,  }\CommentTok{\# New trait number}
  \AttributeTok{combi.weights =} \FunctionTok{c}\NormalTok{(}\DecValTok{1}\NormalTok{, }\DecValTok{10}\NormalTok{, }\DecValTok{5}\NormalTok{),  }\CommentTok{\# Economic weights}
  \AttributeTok{trait.name =} \StringTok{"EconomicIndex"}
\NormalTok{)}
\end{Highlighting}
\end{Shaded}

\section{Trait Standardization}\label{sec-standardization}

Ensure traits have desired mean and variance:

\begin{Shaded}
\begin{Highlighting}[]
\NormalTok{population }\OtherTok{\textless{}{-}} \FunctionTok{bv.standardization}\NormalTok{(}
\NormalTok{  population,}
  \AttributeTok{trait =} \DecValTok{1}\NormalTok{,              }\CommentTok{\# Which trait}
  \AttributeTok{mean.target =} \DecValTok{100}\NormalTok{,      }\CommentTok{\# Target mean}
  \AttributeTok{var.target =} \DecValTok{1}\NormalTok{,         }\CommentTok{\# Target variance}
  \AttributeTok{set.zero =} \ConstantTok{TRUE}         \CommentTok{\# Force 0 allele effect = 0}
\NormalTok{)}
\end{Highlighting}
\end{Shaded}

\textbf{When to use:} - After custom trait specification - To compare
traits on same scale - Before economic index creation

\section{Practical Examples}\label{sec-examples-traits}

\subsection{Example 1: Dairy Cattle (Multiple
Traits)}\label{example-1-dairy-cattle-multiple-traits}

\begin{Shaded}
\begin{Highlighting}[]
\CommentTok{\# Correlated production traits}
\NormalTok{cor\_matrix }\OtherTok{\textless{}{-}} \FunctionTok{matrix}\NormalTok{(}\FunctionTok{c}\NormalTok{(}
  \FloatTok{1.0}\NormalTok{,  }\FloatTok{0.6}\NormalTok{,  }\FloatTok{0.5}\NormalTok{,   }\CommentTok{\# Milk}
  \FloatTok{0.6}\NormalTok{,  }\FloatTok{1.0}\NormalTok{,  }\FloatTok{0.7}\NormalTok{,   }\CommentTok{\# Fat}
  \FloatTok{0.5}\NormalTok{,  }\FloatTok{0.7}\NormalTok{,  }\FloatTok{1.0}    \CommentTok{\# Protein}
\NormalTok{), }\AttributeTok{nrow =} \DecValTok{3}\NormalTok{, }\AttributeTok{byrow =} \ConstantTok{TRUE}\NormalTok{)}

\NormalTok{dairy }\OtherTok{\textless{}{-}} \FunctionTok{creating.diploid}\NormalTok{(}
  \AttributeTok{nsnp =} \DecValTok{50000}\NormalTok{,}
  \AttributeTok{nindi =} \DecValTok{500}\NormalTok{,}
  \AttributeTok{template.chip =} \StringTok{"cattle"}\NormalTok{,}
  \AttributeTok{n.additive =} \FunctionTok{c}\NormalTok{(}\DecValTok{200}\NormalTok{, }\DecValTok{150}\NormalTok{, }\DecValTok{150}\NormalTok{),}
  \AttributeTok{n.dominant =} \FunctionTok{c}\NormalTok{(}\DecValTok{20}\NormalTok{, }\DecValTok{10}\NormalTok{, }\DecValTok{10}\NormalTok{),}
  \AttributeTok{trait.cor =}\NormalTok{ cor\_matrix,}
  \AttributeTok{var.target =} \FunctionTok{c}\NormalTok{(}\DecValTok{2000}\NormalTok{, }\DecValTok{50}\NormalTok{, }\DecValTok{30}\NormalTok{),}
  \AttributeTok{mean.target =} \FunctionTok{c}\NormalTok{(}\DecValTok{8000}\NormalTok{, }\DecValTok{350}\NormalTok{, }\DecValTok{320}\NormalTok{),}
  \AttributeTok{trait.name =} \FunctionTok{c}\NormalTok{(}\StringTok{"Milk\_kg"}\NormalTok{, }\StringTok{"Fat\_kg"}\NormalTok{, }\StringTok{"Protein\_kg"}\NormalTok{)}
\NormalTok{)}
\end{Highlighting}
\end{Shaded}

\subsection{Example 2: Maize Hybrid
Performance}\label{example-2-maize-hybrid-performance}

\begin{Shaded}
\begin{Highlighting}[]
\CommentTok{\# Model heterosis with dominance}
\NormalTok{maize }\OtherTok{\textless{}{-}} \FunctionTok{creating.diploid}\NormalTok{(}
  \AttributeTok{nsnp =} \DecValTok{10000}\NormalTok{,}
  \AttributeTok{nindi =} \DecValTok{20}\NormalTok{,}
  \AttributeTok{template.chip =} \StringTok{"maize"}\NormalTok{,}
  \AttributeTok{dataset =} \StringTok{"homorandom"}\NormalTok{,      }\CommentTok{\# Inbred lines}
  \AttributeTok{n.additive =} \DecValTok{200}\NormalTok{,}
  \AttributeTok{n.dominant =} \DecValTok{50}\NormalTok{,}
  \AttributeTok{dominant.only.positive =} \ConstantTok{TRUE}\NormalTok{,  }\CommentTok{\# Heterosis = positive dominance}
  \AttributeTok{var.additive =} \FloatTok{0.7}\NormalTok{,}
  \AttributeTok{var.dominant =} \FloatTok{0.3}\NormalTok{,}
  \AttributeTok{trait.name =} \StringTok{"GrainYield"}
\NormalTok{)}
\end{Highlighting}
\end{Shaded}

\subsection{Example 3: Disease Resistance (Binary
Trait)}\label{example-3-disease-resistance-binary-trait}

\begin{Shaded}
\begin{Highlighting}[]
\CommentTok{\# Create underlying continuous trait}
\NormalTok{pop }\OtherTok{\textless{}{-}} \FunctionTok{creating.diploid}\NormalTok{(}
  \AttributeTok{nsnp =} \DecValTok{5000}\NormalTok{, }\AttributeTok{nindi =} \DecValTok{500}\NormalTok{,}
  \AttributeTok{n.additive =} \DecValTok{50}\NormalTok{,}
  \AttributeTok{mean.target =} \DecValTok{0}\NormalTok{,}
  \AttributeTok{var.target =} \DecValTok{1}\NormalTok{,}
  \AttributeTok{trait.name =} \StringTok{"Liability"}
\NormalTok{)}

\CommentTok{\# Apply threshold transformation}
\NormalTok{pop }\OtherTok{\textless{}{-}} \FunctionTok{creating.phenotypic.transform}\NormalTok{(}
\NormalTok{  pop,}
  \AttributeTok{trait =} \DecValTok{1}\NormalTok{,}
  \AttributeTok{trafo.function =} \ControlFlowTok{function}\NormalTok{(x) \{ }\FunctionTok{return}\NormalTok{(x }\SpecialCharTok{\textgreater{}} \DecValTok{0}\NormalTok{) \}  }\CommentTok{\# 50\% prevalence}
\NormalTok{)}
\end{Highlighting}
\end{Shaded}

\section{Summary}\label{summary-4}

\begin{itemize}
\tightlist
\item
  ✅ Automated trait generation for common architectures
\item
  ✅ Custom effects for precise control
\item
  ✅ Multiple correlated traits supported
\item
  ✅ Dominance and epistasis modeling
\item
  ✅ Special traits: maternal, breed-specific, polygenic
\item
  ✅ Flexible trait combination and standardization
\end{itemize}

\section{What's Next?}\label{whats-next-4}

With realistic populations and traits, let's build more complex
scenarios with LD structure and data import.

Continue to \href{06-advanced-population-setup.qmd}{Chapter 6: Advanced
Population Setup}!

\chapter{Advanced Population Setup}\label{sec-advanced-population-setup}

\section{Learning Objectives}\label{learning-objectives-5}

\begin{itemize}
\tightlist
\item
  Generate realistic LD structure with \texttt{founder.simulation()}
\item
  Import and use real genomic/phenotypic data
\item
  Create populations with selection signatures
\item
  Model population history
\end{itemize}

\section{Creating Realistic LD Structure}\label{sec-ld-structure}

Real populations have \textbf{linkage disequilibrium (LD)} from
historical selection and drift. The \texttt{founder.simulation()}
function creates this automatically:

\begin{Shaded}
\begin{Highlighting}[]
\CommentTok{\# Create population with realistic LD}
\NormalTok{population }\OtherTok{\textless{}{-}} \FunctionTok{creating.diploid}\NormalTok{(}\AttributeTok{nsnp =} \DecValTok{10000}\NormalTok{, }\AttributeTok{nindi =} \DecValTok{100}\NormalTok{)}

\NormalTok{population }\OtherTok{\textless{}{-}} \FunctionTok{founder.simulation}\NormalTok{(}
\NormalTok{  population,}
  \AttributeTok{generations =} \DecValTok{100}\NormalTok{,        }\CommentTok{\# 100 generations of random mating}
  \AttributeTok{breeding.size =} \DecValTok{200}\NormalTok{,      }\CommentTok{\# Population size per generation}
  \AttributeTok{mutation.rate =} \FloatTok{0.0001}    \CommentTok{\# Mutation rate per site}
\NormalTok{)}
\end{Highlighting}
\end{Shaded}

\textbf{What happens:} - Simulates historical random mating - Creates
realistic LD decay with distance - Generates allele frequency spectrum -
Produces heterozygosity patterns

\subsection{Controlling LD Levels}\label{controlling-ld-levels}

\begin{Shaded}
\begin{Highlighting}[]
\CommentTok{\# Low LD (large historical population)}
\NormalTok{pop\_low\_ld }\OtherTok{\textless{}{-}} \FunctionTok{founder.simulation}\NormalTok{(}
  \FunctionTok{creating.diploid}\NormalTok{(}\AttributeTok{nsnp =} \DecValTok{5000}\NormalTok{, }\AttributeTok{nindi =} \DecValTok{100}\NormalTok{),}
  \AttributeTok{generations =} \DecValTok{1000}\NormalTok{,}
  \AttributeTok{breeding.size =} \DecValTok{1000}
\NormalTok{)}

\CommentTok{\# High LD (bottleneck)}
\NormalTok{pop\_high\_ld }\OtherTok{\textless{}{-}} \FunctionTok{founder.simulation}\NormalTok{(}
  \FunctionTok{creating.diploid}\NormalTok{(}\AttributeTok{nsnp =} \DecValTok{5000}\NormalTok{, }\AttributeTok{nindi =} \DecValTok{100}\NormalTok{),}
  \AttributeTok{generations =} \DecValTok{100}\NormalTok{,}
  \AttributeTok{breeding.size =} \DecValTok{20}        \CommentTok{\# Small Ne = high LD}
\NormalTok{)}
\end{Highlighting}
\end{Shaded}

\section{Modeling Selection Signatures}\label{sec-selection-signatures}

Simulate hard sweeps (selective sweeps):

\begin{Shaded}
\begin{Highlighting}[]
\CommentTok{\# Create population}
\NormalTok{pop }\OtherTok{\textless{}{-}} \FunctionTok{creating.diploid}\NormalTok{(}
  \AttributeTok{nsnp =} \DecValTok{10000}\NormalTok{,}
  \AttributeTok{nindi =} \DecValTok{200}\NormalTok{,}
  \AttributeTok{chr.nr =} \DecValTok{1}\NormalTok{,}
  \AttributeTok{n.additive =} \DecValTok{100}
\NormalTok{)}

\CommentTok{\# Historical selection on trait 1}
\NormalTok{pop }\OtherTok{\textless{}{-}} \FunctionTok{founder.simulation}\NormalTok{(}
\NormalTok{  pop,}
  \AttributeTok{generations =} \DecValTok{50}\NormalTok{,}
  \AttributeTok{breeding.size =} \DecValTok{200}\NormalTok{,}
  \AttributeTok{selection.criteria =} \StringTok{"bv"}\NormalTok{,     }\CommentTok{\# Select on breeding values}
  \AttributeTok{selection.size =} \FunctionTok{c}\NormalTok{(}\DecValTok{20}\NormalTok{, }\DecValTok{20}\NormalTok{),    }\CommentTok{\# Strong selection}
  \AttributeTok{selection.m.gen =} \StringTok{"all"}\NormalTok{,}
  \AttributeTok{selection.f.gen =} \StringTok{"all"}
\NormalTok{)}
\end{Highlighting}
\end{Shaded}

\textbf{Creates:} - Reduced diversity around selected QTLs - Extended
haplotype homozygosity - Realistic post-selection LD patterns

\section{Importing Real Phenotypic Data}\label{sec-import-pheno}

Use real phenotypes to estimate trait architecture:

\begin{Shaded}
\begin{Highlighting}[]
\CommentTok{\# Example: Load data from MoBPS example population}
\FunctionTok{data}\NormalTok{(ex\_pop)}

\CommentTok{\# Extract information}
\NormalTok{pheno }\OtherTok{\textless{}{-}} \FunctionTok{get.pheno}\NormalTok{(ex\_pop, }\AttributeTok{gen =} \DecValTok{1}\SpecialCharTok{:}\DecValTok{5}\NormalTok{)}
\NormalTok{geno }\OtherTok{\textless{}{-}} \FunctionTok{get.geno}\NormalTok{(ex\_pop, }\AttributeTok{gen =} \DecValTok{1}\SpecialCharTok{:}\DecValTok{5}\NormalTok{)}
\NormalTok{map }\OtherTok{\textless{}{-}} \FunctionTok{get.map}\NormalTok{(ex\_pop, }\AttributeTok{use.snp.nr =} \ConstantTok{TRUE}\NormalTok{)}

\CommentTok{\# Estimate marker effects from data}
\NormalTok{estimated\_effects }\OtherTok{\textless{}{-}} \FunctionTok{effect.estimate.add}\NormalTok{(geno, pheno, map)}

\CommentTok{\# Create new population with estimated architecture}
\NormalTok{population }\OtherTok{\textless{}{-}} \FunctionTok{creating.diploid}\NormalTok{(}
  \AttributeTok{dataset =} \FunctionTok{get.haplo}\NormalTok{(ex\_pop, }\AttributeTok{gen =} \DecValTok{1}\NormalTok{),}
  \AttributeTok{map =}\NormalTok{ map,}
  \AttributeTok{real.bv.add =}\NormalTok{ estimated\_effects}
\NormalTok{)}
\end{Highlighting}
\end{Shaded}

\textbf{Workflow:} 1. Import real genotypes + phenotypes 2. Estimate
marker effects (GWAS/GBLUP) 3. Use estimated effects in simulation

\section{Creating Age-Structured Populations}\label{sec-age-structure}

Model overlapping generations:

\begin{Shaded}
\begin{Highlighting}[]
\CommentTok{\# Create founders (year 0)}
\NormalTok{pop }\OtherTok{\textless{}{-}} \FunctionTok{creating.diploid}\NormalTok{(}
  \AttributeTok{nsnp =} \DecValTok{5000}\NormalTok{, }\AttributeTok{nindi =} \DecValTok{100}\NormalTok{,}
  \AttributeTok{n.additive =} \DecValTok{100}\NormalTok{,}
  \AttributeTok{name.cohort =} \StringTok{"Year\_0"}
\NormalTok{)}

\CommentTok{\# Year 1: Some old + new animals}
\NormalTok{pop }\OtherTok{\textless{}{-}} \FunctionTok{breeding.diploid}\NormalTok{(}
\NormalTok{  pop,}
  \AttributeTok{selection.m.cohorts =} \StringTok{"Year\_0"}\NormalTok{,}
  \AttributeTok{selection.f.cohorts =} \StringTok{"Year\_0"}\NormalTok{,}
  \AttributeTok{selection.size =} \FunctionTok{c}\NormalTok{(}\DecValTok{20}\NormalTok{, }\DecValTok{40}\NormalTok{),    }\CommentTok{\# Keep some old animals}
  \AttributeTok{breeding.size =} \FunctionTok{c}\NormalTok{(}\DecValTok{30}\NormalTok{, }\DecValTok{60}\NormalTok{),     }\CommentTok{\# Plus new calves}
  \AttributeTok{name.cohort =} \StringTok{"Year\_1"}
\NormalTok{)}

\CommentTok{\# Select from BOTH cohorts in Year 2}
\NormalTok{pop }\OtherTok{\textless{}{-}} \FunctionTok{breeding.diploid}\NormalTok{(}
\NormalTok{  pop,}
  \AttributeTok{selection.m.cohorts =} \FunctionTok{c}\NormalTok{(}\StringTok{"Year\_0"}\NormalTok{, }\StringTok{"Year\_1"}\NormalTok{),}
  \AttributeTok{selection.f.cohorts =} \FunctionTok{c}\NormalTok{(}\StringTok{"Year\_0"}\NormalTok{, }\StringTok{"Year\_1"}\NormalTok{),}
  \AttributeTok{selection.size =} \FunctionTok{c}\NormalTok{(}\DecValTok{20}\NormalTok{, }\DecValTok{40}\NormalTok{),}
  \AttributeTok{breeding.size =} \FunctionTok{c}\NormalTok{(}\DecValTok{30}\NormalTok{, }\DecValTok{60}\NormalTok{),}
  \AttributeTok{name.cohort =} \StringTok{"Year\_2"}
\NormalTok{)}
\end{Highlighting}
\end{Shaded}

\textbf{Applications:} - Dairy cattle (multi-parity cows) - Progeny
testing schemes - Conservation programs

\section{Multi-Stage Pedigree Simulation}\label{sec-pedigree-simulation}

Create populations from pedigree structure only:

\begin{Shaded}
\begin{Highlighting}[]
\CommentTok{\# Generate pedigree{-}based population}
\NormalTok{pop }\OtherTok{\textless{}{-}} \FunctionTok{pedigree.simulation}\NormalTok{(}
  \AttributeTok{population =} \FunctionTok{creating.diploid}\NormalTok{(}\AttributeTok{nsnp =} \DecValTok{1000}\NormalTok{, }\AttributeTok{nindi =} \DecValTok{10}\NormalTok{),}
  \AttributeTok{breeding.size =} \DecValTok{100}\NormalTok{,}
  \AttributeTok{generations =} \DecValTok{5}\NormalTok{,}
  \AttributeTok{selection.size =} \FunctionTok{c}\NormalTok{(}\DecValTok{5}\NormalTok{, }\DecValTok{10}\NormalTok{)}
\NormalTok{)}
\end{Highlighting}
\end{Shaded}

\textbf{Use cases:} - Quick exploration without full genomics - Very
large pedigrees - When genotypes not needed

\section{Example: Import VCF + GWAS Results}\label{sec-example-vcf-gwas}

Complete workflow with real data:

\begin{Shaded}
\begin{Highlighting}[]
\CommentTok{\# Step 1: Import VCF file}
\NormalTok{pop }\OtherTok{\textless{}{-}} \FunctionTok{creating.diploid}\NormalTok{(}
  \AttributeTok{vcf =} \StringTok{"path/to/cattle\_genotypes.vcf.gz"}\NormalTok{,}
  \AttributeTok{vcf.maxsnp =} \DecValTok{50000}\NormalTok{,     }\CommentTok{\# Limit to 50K SNPs}
  \AttributeTok{vcf.maxindi =} \DecValTok{1000}      \CommentTok{\# Limit to 1000 animals}
\NormalTok{)}

\CommentTok{\# Step 2: Load phenotypes (from external file)}
\NormalTok{pheno\_data }\OtherTok{\textless{}{-}} \FunctionTok{read.csv}\NormalTok{(}\StringTok{"phenotypes.csv"}\NormalTok{)}
\CommentTok{\# Assume columns: AnimalID, Trait1, Trait2, Trait3}

\CommentTok{\# Step 3: Match phenotypes to individuals}
\CommentTok{\# (You\textquotesingle{}ll need to map VCF sample IDs to phenotype IDs)}

\CommentTok{\# Step 4: Estimate effects using your favorite method}
\CommentTok{\# Example with rrBLUP}
\FunctionTok{library}\NormalTok{(rrBLUP)}
\NormalTok{geno\_matrix }\OtherTok{\textless{}{-}} \FunctionTok{get.geno}\NormalTok{(pop, }\AttributeTok{gen =} \DecValTok{1}\NormalTok{)}
\NormalTok{effects }\OtherTok{\textless{}{-}} \FunctionTok{mixed.solve}\NormalTok{(pheno\_data}\SpecialCharTok{$}\NormalTok{Trait1, }\AttributeTok{Z =} \FunctionTok{t}\NormalTok{(geno\_matrix))}

\CommentTok{\# Step 5: Create trait architecture from estimated effects}
\NormalTok{qtl\_effects }\OtherTok{\textless{}{-}} \FunctionTok{cbind}\NormalTok{(}
  \DecValTok{1}\SpecialCharTok{:}\FunctionTok{length}\NormalTok{(effects}\SpecialCharTok{$}\NormalTok{u),  }\CommentTok{\# SNP number}
  \DecValTok{1}\NormalTok{,                    }\CommentTok{\# Chromosome (placeholder)}
  \SpecialCharTok{{-}}\NormalTok{effects}\SpecialCharTok{$}\NormalTok{u,           }\CommentTok{\# Effect for allele 0}
  \DecValTok{0}\NormalTok{,                    }\CommentTok{\# Effect for allele 1}
\NormalTok{  effects}\SpecialCharTok{$}\NormalTok{u             }\CommentTok{\# Effect for allele 2}
\NormalTok{)}

\NormalTok{pop }\OtherTok{\textless{}{-}} \FunctionTok{creating.trait}\NormalTok{(}
\NormalTok{  pop,}
  \AttributeTok{real.bv.add =}\NormalTok{ qtl\_effects,}
  \AttributeTok{trait.name =} \StringTok{"EstimatedTrait"}
\NormalTok{)}

\CommentTok{\# Now ready to simulate breeding program!}
\end{Highlighting}
\end{Shaded}

\section{Summary}\label{summary-5}

\begin{itemize}
\tightlist
\item
  ✅ \texttt{founder.simulation()} creates realistic LD structure
\item
  ✅ Model selection signatures and population history
\item
  ✅ Import real data (genotypes, phenotypes, GWAS results)
\item
  ✅ Create age-structured populations
\item
  ✅ Flexible approaches for complex scenarios
\end{itemize}

Continue to \href{07-population-examples.qmd}{Chapter 7: Population
Examples}!

\chapter{Population Examples}\label{sec-population-examples}

\section{Learning Objectives}\label{learning-objectives-6}

\begin{itemize}
\tightlist
\item
  See complete worked examples of population creation
\item
  Learn patterns for common scenarios
\item
  Understand how to combine features
\end{itemize}

\section{Example 1: Inbred Lines (DH
Production)}\label{sec-example-inbred}

\begin{Shaded}
\begin{Highlighting}[]
\CommentTok{\# Heterozygous founders}
\NormalTok{pop }\OtherTok{\textless{}{-}} \FunctionTok{creating.diploid}\NormalTok{(}
  \AttributeTok{nindi =} \DecValTok{50}\NormalTok{,}
  \AttributeTok{nsnp =} \DecValTok{5000}\NormalTok{,}
  \AttributeTok{sex.quota =} \DecValTok{0}\NormalTok{,          }\CommentTok{\# All "male" (no sex distinction)}
  \AttributeTok{n.additive =} \DecValTok{100}
\NormalTok{)}

\CommentTok{\# Create DH lines (instant homozygosity)}
\NormalTok{pop }\OtherTok{\textless{}{-}} \FunctionTok{breeding.diploid}\NormalTok{(}
\NormalTok{  pop,}
  \AttributeTok{breeding.size =} \FunctionTok{c}\NormalTok{(}\DecValTok{200}\NormalTok{, }\DecValTok{0}\NormalTok{),}
  \AttributeTok{dh.mating =} \ConstantTok{TRUE}\NormalTok{,        }\CommentTok{\# Doubled haploid technology}
  \AttributeTok{name.cohort =} \StringTok{"DH\_Lines"}
\NormalTok{)}

\CommentTok{\# Check homozygosity}
\FunctionTok{table}\NormalTok{(}\FunctionTok{get.geno}\NormalTok{(pop, }\AttributeTok{gen =} \DecValTok{2}\NormalTok{))  }\CommentTok{\# Should be mostly 0 and 2, few 1s}
\end{Highlighting}
\end{Shaded}

\textbf{Compare with selfing:}

\begin{Shaded}
\begin{Highlighting}[]
\CommentTok{\# 5 generations of selfing}
\NormalTok{pop\_selfing }\OtherTok{\textless{}{-}}\NormalTok{ pop}
\ControlFlowTok{for}\NormalTok{ (i }\ControlFlowTok{in} \DecValTok{1}\SpecialCharTok{:}\DecValTok{5}\NormalTok{) \{}
\NormalTok{  pop\_selfing }\OtherTok{\textless{}{-}} \FunctionTok{breeding.diploid}\NormalTok{(}
\NormalTok{    pop\_selfing,}
    \AttributeTok{breeding.size =} \FunctionTok{c}\NormalTok{(}\DecValTok{200}\NormalTok{, }\DecValTok{0}\NormalTok{),}
    \AttributeTok{selfing.mating =} \ConstantTok{TRUE}
\NormalTok{  )}
\NormalTok{\}}

\CommentTok{\# Check reduction in heterozygosity}
\ControlFlowTok{for}\NormalTok{ (gen }\ControlFlowTok{in} \DecValTok{1}\SpecialCharTok{:}\DecValTok{6}\NormalTok{) \{}
  \FunctionTok{print}\NormalTok{(}\FunctionTok{table}\NormalTok{(}\FunctionTok{get.geno}\NormalTok{(pop\_selfing, }\AttributeTok{gen =}\NormalTok{ gen)))}
\NormalTok{\}}
\CommentTok{\# \textasciitilde{}50\% reduction per generation}
\end{Highlighting}
\end{Shaded}

\section{Example 2: MAGIC Population (Maize)}\label{sec-example-magic}

Multi-parent Advanced Generation Inter-Cross:

\begin{Shaded}
\begin{Highlighting}[]
\CommentTok{\# 20 fully homozygous founder lines}
\NormalTok{pop }\OtherTok{\textless{}{-}} \FunctionTok{creating.diploid}\NormalTok{(}
  \AttributeTok{nindi =} \DecValTok{20}\NormalTok{,}
  \AttributeTok{sex.quota =} \DecValTok{0}\NormalTok{,}
  \AttributeTok{template.chip =} \StringTok{"maize"}\NormalTok{,}
  \AttributeTok{dataset =} \StringTok{"homorandom"}\NormalTok{,}
  \AttributeTok{n.additive =} \DecValTok{200}
\NormalTok{)}

\CommentTok{\# Round 1: All pairwise crosses (190 F1s)}
\NormalTok{pop }\OtherTok{\textless{}{-}} \FunctionTok{breeding.diploid}\NormalTok{(}
\NormalTok{  pop,}
  \AttributeTok{breeding.size =} \FunctionTok{c}\NormalTok{(}\DecValTok{190}\NormalTok{, }\DecValTok{0}\NormalTok{),}
  \AttributeTok{breeding.all.combination =} \ConstantTok{TRUE}\NormalTok{,  }\CommentTok{\# All combinations}
  \AttributeTok{selection.size =} \FunctionTok{c}\NormalTok{(}\DecValTok{20}\NormalTok{, }\DecValTok{0}\NormalTok{),}
  \AttributeTok{max.offspring =} \DecValTok{19}                \CommentTok{\# Each parent in 19 crosses}
\NormalTok{)}

\CommentTok{\# Rounds 2{-}4: Random inter{-}mating}
\ControlFlowTok{for}\NormalTok{ (round }\ControlFlowTok{in} \DecValTok{1}\SpecialCharTok{:}\DecValTok{3}\NormalTok{) \{}
\NormalTok{  pop }\OtherTok{\textless{}{-}} \FunctionTok{breeding.diploid}\NormalTok{(}
\NormalTok{    pop,}
    \AttributeTok{breeding.size =} \FunctionTok{c}\NormalTok{(}\DecValTok{190}\NormalTok{, }\DecValTok{0}\NormalTok{),}
    \AttributeTok{selection.size =} \FunctionTok{c}\NormalTok{(}\DecValTok{190}\NormalTok{, }\DecValTok{0}\NormalTok{),}
    \AttributeTok{same.sex.activ =} \ConstantTok{TRUE}\NormalTok{,}
    \AttributeTok{same.sex.sex =} \DecValTok{0}\NormalTok{,}
    \AttributeTok{max.offspring =} \DecValTok{2}
\NormalTok{  )}
\NormalTok{\}}

\CommentTok{\# Final: Create RILs by selfing}
\NormalTok{pop }\OtherTok{\textless{}{-}} \FunctionTok{breeding.diploid}\NormalTok{(}
\NormalTok{  pop,}
  \AttributeTok{breeding.size =} \FunctionTok{c}\NormalTok{(}\DecValTok{380}\NormalTok{, }\DecValTok{0}\NormalTok{),}
  \AttributeTok{selfing.mating =} \ConstantTok{TRUE}
\NormalTok{)}
\end{Highlighting}
\end{Shaded}

\section{Example 3: Crossbreeding
Setup}\label{sec-example-crossbreeding}

Two-breed cross with breed-specific QTLs:

\begin{Shaded}
\begin{Highlighting}[]
\CommentTok{\# Breed A}
\NormalTok{pop }\OtherTok{\textless{}{-}} \FunctionTok{creating.diploid}\NormalTok{(}
  \AttributeTok{nsnp =} \DecValTok{5000}\NormalTok{, }\AttributeTok{nindi =} \DecValTok{100}\NormalTok{,}
  \AttributeTok{founder.pool =} \DecValTok{1}\NormalTok{,}
  \AttributeTok{name.cohort =} \StringTok{"Breed\_A"}\NormalTok{,}
  \AttributeTok{n.additive =} \DecValTok{0}  \CommentTok{\# Add traits later}
\NormalTok{)}

\CommentTok{\# Breed B}
\NormalTok{pop }\OtherTok{\textless{}{-}} \FunctionTok{creating.diploid}\NormalTok{(}
  \AttributeTok{population =}\NormalTok{ pop,}
  \AttributeTok{nsnp =} \DecValTok{5000}\NormalTok{, }\AttributeTok{nindi =} \DecValTok{100}\NormalTok{,}
  \AttributeTok{founder.pool =} \DecValTok{2}\NormalTok{,}
  \AttributeTok{name.cohort =} \StringTok{"Breed\_B"}\NormalTok{,}
  \AttributeTok{freq =} \StringTok{"diff"}  \CommentTok{\# Different allele frequencies}
\NormalTok{)}

\CommentTok{\# Trait only in Breed A}
\NormalTok{pop }\OtherTok{\textless{}{-}} \FunctionTok{creating.trait}\NormalTok{(pop, }\AttributeTok{n.additive =} \DecValTok{100}\NormalTok{, }\AttributeTok{trait.pool =} \DecValTok{1}\NormalTok{)}

\CommentTok{\# Trait only in Breed B}
\NormalTok{pop }\OtherTok{\textless{}{-}} \FunctionTok{creating.trait}\NormalTok{(pop, }\AttributeTok{n.additive =} \DecValTok{100}\NormalTok{, }\AttributeTok{trait.pool =} \DecValTok{2}\NormalTok{)}

\CommentTok{\# Get QTL effects}
\NormalTok{qtl\_eff }\OtherTok{\textless{}{-}} \FunctionTok{get.qtl.effects}\NormalTok{(pop)}

\CommentTok{\# Combine into joint trait}
\NormalTok{combined\_effects }\OtherTok{\textless{}{-}} \FunctionTok{rbind}\NormalTok{(qtl\_eff[[}\DecValTok{1}\NormalTok{]][[}\DecValTok{1}\NormalTok{]], qtl\_eff[[}\DecValTok{1}\NormalTok{]][[}\DecValTok{2}\NormalTok{]])}

\NormalTok{pop }\OtherTok{\textless{}{-}} \FunctionTok{creating.trait}\NormalTok{(}
\NormalTok{  pop,}
  \AttributeTok{real.bv.add =}\NormalTok{ combined\_effects,}
  \AttributeTok{replace.traits =} \ConstantTok{TRUE}\NormalTok{,}
  \AttributeTok{trait.name =} \StringTok{"Combined"}
\NormalTok{)}

\CommentTok{\# Cross the breeds}
\NormalTok{pop }\OtherTok{\textless{}{-}} \FunctionTok{breeding.diploid}\NormalTok{(}
\NormalTok{  pop,}
  \AttributeTok{selection.m.cohorts =} \StringTok{"Breed\_A"}\NormalTok{,}
  \AttributeTok{selection.f.cohorts =} \StringTok{"Breed\_B"}\NormalTok{,}
  \AttributeTok{breeding.size =} \FunctionTok{c}\NormalTok{(}\DecValTok{100}\NormalTok{, }\DecValTok{100}\NormalTok{),}
  \AttributeTok{name.cohort =} \StringTok{"F1\_Cross"}
\NormalTok{)}

\CommentTok{\# Visualize breed composition}
\NormalTok{pools }\OtherTok{\textless{}{-}} \FunctionTok{get.pool}\NormalTok{(pop, }\AttributeTok{gen =} \DecValTok{3}\NormalTok{, }\AttributeTok{plot =} \ConstantTok{TRUE}\NormalTok{)}
\end{Highlighting}
\end{Shaded}

\section{Example 4: Maternal Effects}\label{sec-example-maternal}

Birth weight with direct and maternal components:

\begin{Shaded}
\begin{Highlighting}[]
\CommentTok{\# Create two traits: direct and maternal}
\NormalTok{pop }\OtherTok{\textless{}{-}} \FunctionTok{creating.diploid}\NormalTok{(}
  \AttributeTok{nsnp =} \DecValTok{5000}\NormalTok{,}
  \AttributeTok{nindi =} \DecValTok{200}\NormalTok{,}
  \AttributeTok{n.additive =} \FunctionTok{c}\NormalTok{(}\DecValTok{100}\NormalTok{, }\DecValTok{100}\NormalTok{),}
  \AttributeTok{var.target =} \FunctionTok{c}\NormalTok{(}\DecValTok{100}\NormalTok{, }\DecValTok{20}\NormalTok{),}
  \AttributeTok{trait.cor =} \FunctionTok{matrix}\NormalTok{(}\FunctionTok{c}\NormalTok{(}\DecValTok{1}\NormalTok{, }\FloatTok{0.5}\NormalTok{, }\FloatTok{0.5}\NormalTok{, }\DecValTok{1}\NormalTok{), }\AttributeTok{nrow =} \DecValTok{2}\NormalTok{),}
  \AttributeTok{is.maternal =} \FunctionTok{c}\NormalTok{(}\ConstantTok{FALSE}\NormalTok{, }\ConstantTok{TRUE}\NormalTok{),}
  \AttributeTok{trait.name =} \FunctionTok{c}\NormalTok{(}\StringTok{"Direct"}\NormalTok{, }\StringTok{"Maternal"}\NormalTok{)}
\NormalTok{)}

\CommentTok{\# Combine into observed trait}
\NormalTok{pop }\OtherTok{\textless{}{-}} \FunctionTok{add.combi}\NormalTok{(}
\NormalTok{  pop,}
  \AttributeTok{trait =} \DecValTok{3}\NormalTok{,}
  \AttributeTok{combi.weights =} \FunctionTok{c}\NormalTok{(}\DecValTok{1}\NormalTok{, }\DecValTok{1}\NormalTok{),}
  \AttributeTok{trait.name =} \StringTok{"BirthWeight"}
\NormalTok{)}

\CommentTok{\# Generate offspring}
\NormalTok{pop }\OtherTok{\textless{}{-}} \FunctionTok{breeding.diploid}\NormalTok{(pop, }\AttributeTok{breeding.size =} \FunctionTok{c}\NormalTok{(}\DecValTok{100}\NormalTok{, }\DecValTok{100}\NormalTok{))}

\CommentTok{\# Check: maternal trait depends on dam\textquotesingle{}s genotype}
\NormalTok{bvs }\OtherTok{\textless{}{-}} \FunctionTok{get.bv}\NormalTok{(pop, }\AttributeTok{gen =} \DecValTok{2}\NormalTok{)}
\NormalTok{pedi }\OtherTok{\textless{}{-}} \FunctionTok{get.pedigree}\NormalTok{(pop, }\AttributeTok{gen =} \DecValTok{2}\NormalTok{)}
\FunctionTok{head}\NormalTok{(}\FunctionTok{cbind}\NormalTok{(pedi[, }\DecValTok{3}\NormalTok{], bvs[}\DecValTok{2}\NormalTok{, ]))  }\CommentTok{\# Dam ID vs maternal BV}
\end{Highlighting}
\end{Shaded}

\section{Example 5: Import VCF + Create
Trait}\label{sec-example-vcf-import}

\begin{Shaded}
\begin{Highlighting}[]
\CommentTok{\# Import real genotypes}
\NormalTok{pop }\OtherTok{\textless{}{-}} \FunctionTok{creating.diploid}\NormalTok{(}
  \AttributeTok{vcf =} \StringTok{"holstein\_1000genomes.vcf.gz"}\NormalTok{,}
  \AttributeTok{vcf.maxsnp =} \DecValTok{50000}\NormalTok{,}
  \AttributeTok{vcf.maxindi =} \DecValTok{500}
\NormalTok{)}

\CommentTok{\# Add simulated trait}
\NormalTok{pop }\OtherTok{\textless{}{-}} \FunctionTok{creating.trait}\NormalTok{(}
\NormalTok{  pop,}
  \AttributeTok{n.additive =} \DecValTok{200}\NormalTok{,}
  \AttributeTok{n.dominant =} \DecValTok{30}\NormalTok{,}
  \AttributeTok{var.target =} \DecValTok{100}\NormalTok{,}
  \AttributeTok{mean.target =} \DecValTok{8000}\NormalTok{,}
  \AttributeTok{trait.name =} \StringTok{"MilkYield"}
\NormalTok{)}

\CommentTok{\# Generate LD structure}
\NormalTok{pop }\OtherTok{\textless{}{-}} \FunctionTok{founder.simulation}\NormalTok{(}
\NormalTok{  pop,}
  \AttributeTok{generations =} \DecValTok{100}\NormalTok{,}
  \AttributeTok{breeding.size =} \DecValTok{300}
\NormalTok{)}
\end{Highlighting}
\end{Shaded}

\section{Example 6: Hard Sweep (Selection
Signature)}\label{sec-example-hard-sweep}

Simulate past selection:

\begin{Shaded}
\begin{Highlighting}[]
\CommentTok{\# Create base population}
\NormalTok{pop }\OtherTok{\textless{}{-}} \FunctionTok{creating.diploid}\NormalTok{(}
  \AttributeTok{nsnp =} \DecValTok{10000}\NormalTok{,}
  \AttributeTok{nindi =} \DecValTok{200}\NormalTok{,}
  \AttributeTok{chr.nr =} \DecValTok{5}\NormalTok{,}
  \AttributeTok{n.additive =} \DecValTok{100}\NormalTok{,}
  \AttributeTok{var.target =} \DecValTok{1}
\NormalTok{)}

\CommentTok{\# Simulate 50 generations of STRONG selection}
\NormalTok{pop }\OtherTok{\textless{}{-}} \FunctionTok{founder.simulation}\NormalTok{(}
\NormalTok{  pop,}
  \AttributeTok{generations =} \DecValTok{50}\NormalTok{,}
  \AttributeTok{breeding.size =} \DecValTok{200}\NormalTok{,}
  \AttributeTok{selection.criteria =} \StringTok{"bv"}\NormalTok{,}
  \AttributeTok{selection.size =} \FunctionTok{c}\NormalTok{(}\DecValTok{5}\NormalTok{, }\DecValTok{5}\NormalTok{),      }\CommentTok{\# Top 2.5\%!}
  \AttributeTok{phenotyping =} \StringTok{"all"}\NormalTok{,}
  \AttributeTok{heritability =} \FloatTok{0.8}
\NormalTok{)}

\CommentTok{\# Now selective sweep visible in allele frequencies}
\NormalTok{qtls }\OtherTok{\textless{}{-}} \FunctionTok{get.qtl}\NormalTok{(pop)}
\NormalTok{freq\_gen1 }\OtherTok{\textless{}{-}} \FunctionTok{get.allele.freq}\NormalTok{(pop, }\AttributeTok{gen =} \DecValTok{1}\NormalTok{, }\AttributeTok{snp =}\NormalTok{ qtls[[}\DecValTok{1}\NormalTok{]])}
\NormalTok{freq\_gen51 }\OtherTok{\textless{}{-}} \FunctionTok{get.allele.freq}\NormalTok{(pop, }\AttributeTok{gen =} \DecValTok{51}\NormalTok{, }\AttributeTok{snp =}\NormalTok{ qtls[[}\DecValTok{1}\NormalTok{]])}

\FunctionTok{cbind}\NormalTok{(freq\_gen1, freq\_gen51)  }\CommentTok{\# QTL frequencies changed}
\end{Highlighting}
\end{Shaded}

\section{Example 7: GxE
(Genotype-by-Environment)}\label{sec-example-gxe}

\begin{Shaded}
\begin{Highlighting}[]
\CommentTok{\# Manual approach: Create correlated "traits" for each environment}
\NormalTok{pop }\OtherTok{\textless{}{-}} \FunctionTok{creating.diploid}\NormalTok{(}
  \AttributeTok{nsnp =} \DecValTok{5000}\NormalTok{,}
  \AttributeTok{nindi =} \DecValTok{200}\NormalTok{,}
  \AttributeTok{n.additive =} \FunctionTok{c}\NormalTok{(}\DecValTok{100}\NormalTok{, }\DecValTok{100}\NormalTok{),}
  \AttributeTok{trait.cor =} \FunctionTok{matrix}\NormalTok{(}\FunctionTok{c}\NormalTok{(}\DecValTok{1}\NormalTok{, }\FloatTok{0.6}\NormalTok{, }\FloatTok{0.6}\NormalTok{, }\DecValTok{1}\NormalTok{), }\AttributeTok{nrow =} \DecValTok{2}\NormalTok{),  }\CommentTok{\# Moderate correlation}
  \AttributeTok{trait.name =} \FunctionTok{c}\NormalTok{(}\StringTok{"Env1"}\NormalTok{, }\StringTok{"Env2"}\NormalTok{)}
\NormalTok{)}

\CommentTok{\# Automated GxE model}
\NormalTok{pop }\OtherTok{\textless{}{-}} \FunctionTok{creating.diploid}\NormalTok{(}
  \AttributeTok{nsnp =} \DecValTok{5000}\NormalTok{,}
  \AttributeTok{nindi =} \DecValTok{200}\NormalTok{,}
  \AttributeTok{n.additive =} \DecValTok{100}\NormalTok{,}
  \AttributeTok{gxe.model =} \ConstantTok{TRUE}\NormalTok{,}
  \AttributeTok{gxe.covariance =} \FunctionTok{matrix}\NormalTok{(}\FunctionTok{c}\NormalTok{(}\DecValTok{1}\NormalTok{, }\FloatTok{0.6}\NormalTok{, }\FloatTok{0.6}\NormalTok{, }\FloatTok{1.2}\NormalTok{), }\AttributeTok{nrow =} \DecValTok{2}\NormalTok{)}
\NormalTok{)}
\end{Highlighting}
\end{Shaded}

\section{Example 8: Import PLINK Files}\label{sec-example-plink}

\begin{Shaded}
\begin{Highlighting}[]
\CommentTok{\# Read PLINK .ped and .map files}
\NormalTok{ped\_data }\OtherTok{\textless{}{-}} \FunctionTok{read.pedmap}\NormalTok{(}
  \AttributeTok{path\_ped =} \StringTok{"mydata.ped"}\NormalTok{,}
  \AttributeTok{path\_map =} \StringTok{"mydata.map"}
\NormalTok{)}

\CommentTok{\# Create population}
\NormalTok{pop }\OtherTok{\textless{}{-}} \FunctionTok{creating.diploid}\NormalTok{(}
  \AttributeTok{dataset =}\NormalTok{ ped\_data}\SpecialCharTok{$}\NormalTok{dataset,}
  \AttributeTok{map =}\NormalTok{ ped\_data}\SpecialCharTok{$}\NormalTok{map,}
  \AttributeTok{n.additive =} \DecValTok{150}\NormalTok{,}
  \AttributeTok{trait.name =} \StringTok{"ImportedTrait"}
\NormalTok{)}
\end{Highlighting}
\end{Shaded}

\section{Summary}\label{summary-6}

These examples demonstrate:

\begin{itemize}
\tightlist
\item
  ✅ Inbred line creation (DH vs selfing)
\item
  ✅ MAGIC population schemes
\item
  ✅ Crossbreeding with breed-specific QTLs
\item
  ✅ Maternal effects modeling
\item
  ✅ VCF/PLINK data import
\item
  ✅ Selection signatures
\item
  ✅ Genotype-by-environment interactions
\end{itemize}

\section{What's Next?}\label{whats-next-5}

With populations created, let's simulate \textbf{breeding actions} -
selection, mating, and culling.

Continue to \href{08-selection-strategies.qmd}{Chapter 8: Selection
Strategies}!

\part{Selection, Mating \& Breeding}

\chapter{Selection Strategies}\label{sec-selection-strategies}

\section{Overview}\label{overview-1}

Selection is controlled through \texttt{breeding.diploid()} parameters.
Key options:

\begin{itemize}
\tightlist
\item
  \texttt{selection.criteria} - What to select on
\item
  \texttt{selection.size} - How many to select
\item
  \texttt{selection.m.gen} / \texttt{selection.f.gen} - Which
  generations
\item
  \texttt{selection.index} - Multi-trait selection
\end{itemize}

\section{Random Selection}\label{random-selection}

\begin{Shaded}
\begin{Highlighting}[]
\NormalTok{pop }\OtherTok{\textless{}{-}} \FunctionTok{breeding.diploid}\NormalTok{(}
\NormalTok{  pop,}
  \AttributeTok{selection.criteria =} \StringTok{"random"}\NormalTok{,}
  \AttributeTok{selection.size =} \FunctionTok{c}\NormalTok{(}\DecValTok{20}\NormalTok{, }\DecValTok{20}\NormalTok{),}
  \AttributeTok{breeding.size =} \FunctionTok{c}\NormalTok{(}\DecValTok{100}\NormalTok{, }\DecValTok{100}\NormalTok{)}
\NormalTok{)}
\end{Highlighting}
\end{Shaded}

\section{Phenotypic Selection}\label{phenotypic-selection}

\begin{Shaded}
\begin{Highlighting}[]
\CommentTok{\# Phenotype first}
\NormalTok{pop }\OtherTok{\textless{}{-}} \FunctionTok{breeding.diploid}\NormalTok{(pop, }\AttributeTok{phenotyping =} \StringTok{"all"}\NormalTok{, }\AttributeTok{heritability =} \FloatTok{0.5}\NormalTok{)}

\CommentTok{\# Select on phenotypes}
\NormalTok{pop }\OtherTok{\textless{}{-}} \FunctionTok{breeding.diploid}\NormalTok{(}
\NormalTok{  pop,}
  \AttributeTok{selection.criteria =} \StringTok{"pheno"}\NormalTok{,}
  \AttributeTok{selection.size =} \FunctionTok{c}\NormalTok{(}\DecValTok{10}\NormalTok{, }\DecValTok{10}\NormalTok{),}
  \AttributeTok{breeding.size =} \FunctionTok{c}\NormalTok{(}\DecValTok{50}\NormalTok{, }\DecValTok{50}\NormalTok{)}
\NormalTok{)}
\end{Highlighting}
\end{Shaded}

\section{Genomic Selection (BVE)}\label{genomic-selection-bve}

\begin{Shaded}
\begin{Highlighting}[]
\CommentTok{\# Calculate BVEs}
\NormalTok{pop }\OtherTok{\textless{}{-}} \FunctionTok{breeding.diploid}\NormalTok{(pop, }\AttributeTok{bve =} \ConstantTok{TRUE}\NormalTok{, }\AttributeTok{bve.gen =} \DecValTok{1}\SpecialCharTok{:}\DecValTok{3}\NormalTok{)}

\CommentTok{\# Select on estimated breeding values}
\NormalTok{pop }\OtherTok{\textless{}{-}} \FunctionTok{breeding.diploid}\NormalTok{(}
\NormalTok{  pop,}
  \AttributeTok{selection.criteria =} \StringTok{"bve"}\NormalTok{,}
  \AttributeTok{selection.size =} \FunctionTok{c}\NormalTok{(}\DecValTok{10}\NormalTok{, }\DecValTok{10}\NormalTok{),}
  \AttributeTok{breeding.size =} \FunctionTok{c}\NormalTok{(}\DecValTok{50}\NormalTok{, }\DecValTok{50}\NormalTok{)}
\NormalTok{)}
\end{Highlighting}
\end{Shaded}

\section{Selection on True Breeding
Values}\label{selection-on-true-breeding-values}

\begin{Shaded}
\begin{Highlighting}[]
\CommentTok{\# Select on known BVs (simulation only!)}
\NormalTok{pop }\OtherTok{\textless{}{-}} \FunctionTok{breeding.diploid}\NormalTok{(}
\NormalTok{  pop,}
  \AttributeTok{selection.criteria =} \StringTok{"bv"}\NormalTok{,}
  \AttributeTok{selection.size =} \FunctionTok{c}\NormalTok{(}\DecValTok{5}\NormalTok{, }\DecValTok{5}\NormalTok{),}
  \AttributeTok{breeding.size =} \FunctionTok{c}\NormalTok{(}\DecValTok{100}\NormalTok{, }\DecValTok{100}\NormalTok{)}
\NormalTok{)}
\end{Highlighting}
\end{Shaded}

\section{Multi-Trait Selection Index}\label{multi-trait-selection-index}

\begin{Shaded}
\begin{Highlighting}[]
\CommentTok{\# Define economic weights}
\NormalTok{pop }\OtherTok{\textless{}{-}} \FunctionTok{breeding.diploid}\NormalTok{(}
\NormalTok{  pop,}
  \AttributeTok{selection.criteria =} \StringTok{"bve"}\NormalTok{,}
  \AttributeTok{selection.index =} \FunctionTok{c}\NormalTok{(}\DecValTok{1}\NormalTok{, }\DecValTok{10}\NormalTok{, }\DecValTok{5}\NormalTok{),  }\CommentTok{\# Weights for traits 1,2,3}
  \AttributeTok{selection.size =} \FunctionTok{c}\NormalTok{(}\DecValTok{10}\NormalTok{, }\DecValTok{10}\NormalTok{),}
  \AttributeTok{breeding.size =} \FunctionTok{c}\NormalTok{(}\DecValTok{50}\NormalTok{, }\DecValTok{50}\NormalTok{)}
\NormalTok{)}
\end{Highlighting}
\end{Shaded}

\section{Truncation Selection with Variable
Intensity}\label{truncation-selection-with-variable-intensity}

\begin{Shaded}
\begin{Highlighting}[]
\CommentTok{\# Select top males, more females}
\NormalTok{pop }\OtherTok{\textless{}{-}} \FunctionTok{breeding.diploid}\NormalTok{(}
\NormalTok{  pop,}
  \AttributeTok{selection.criteria =} \StringTok{"bve"}\NormalTok{,}
  \AttributeTok{selection.size =} \FunctionTok{c}\NormalTok{(}\DecValTok{5}\NormalTok{, }\DecValTok{20}\NormalTok{),     }\CommentTok{\# 5\% males, 20\% females}
  \AttributeTok{selection.m.gen =} \DecValTok{4}\NormalTok{,}
  \AttributeTok{selection.f.gen =} \DecValTok{4}\NormalTok{,}
  \AttributeTok{breeding.size =} \FunctionTok{c}\NormalTok{(}\DecValTok{50}\NormalTok{, }\DecValTok{50}\NormalTok{)}
\NormalTok{)}
\end{Highlighting}
\end{Shaded}

\section{Threshold Selection}\label{threshold-selection}

\begin{Shaded}
\begin{Highlighting}[]
\CommentTok{\# Only individuals above threshold}
\NormalTok{pop }\OtherTok{\textless{}{-}} \FunctionTok{breeding.diploid}\NormalTok{(}
\NormalTok{  pop,}
  \AttributeTok{selection.criteria =} \StringTok{"bve"}\NormalTok{,}
  \AttributeTok{selection.m.threshold =} \DecValTok{105}\NormalTok{,   }\CommentTok{\# BVE \textgreater{} 105}
  \AttributeTok{selection.f.threshold =} \DecValTok{100}\NormalTok{,   }\CommentTok{\# BVE \textgreater{} 100}
  \AttributeTok{breeding.size =} \FunctionTok{c}\NormalTok{(}\DecValTok{50}\NormalTok{, }\DecValTok{50}\NormalTok{)}
\NormalTok{)}
\end{Highlighting}
\end{Shaded}

\section{Optimal Genetic Contribution
(OGC)}\label{optimal-genetic-contribution-ogc}

Balance genetic gain with diversity:

\begin{Shaded}
\begin{Highlighting}[]
\CommentTok{\# Requires optiSel package}
\FunctionTok{library}\NormalTok{(optiSel)}

\NormalTok{pop }\OtherTok{\textless{}{-}} \FunctionTok{breeding.diploid}\NormalTok{(}
\NormalTok{  pop,}
  \AttributeTok{selection.algorithm =} \StringTok{"ogc"}\NormalTok{,}
  \AttributeTok{ogc.target.inbreeding =} \FloatTok{0.01}\NormalTok{,  }\CommentTok{\# Max 1\% inbreeding increase}
  \AttributeTok{selection.size =} \FunctionTok{c}\NormalTok{(}\DecValTok{20}\NormalTok{, }\DecValTok{20}\NormalTok{),}
  \AttributeTok{breeding.size =} \FunctionTok{c}\NormalTok{(}\DecValTok{100}\NormalTok{, }\DecValTok{100}\NormalTok{)}
\NormalTok{)}
\end{Highlighting}
\end{Shaded}

\section{Custom Selection Function}\label{custom-selection-function}

\begin{Shaded}
\begin{Highlighting}[]
\CommentTok{\# Define custom selection}
\NormalTok{my\_selector }\OtherTok{\textless{}{-}} \ControlFlowTok{function}\NormalTok{(population, gen, sex) \{}
\NormalTok{  bv }\OtherTok{\textless{}{-}} \FunctionTok{get.bv}\NormalTok{(population, }\AttributeTok{gen =}\NormalTok{ gen, }\AttributeTok{sex =}\NormalTok{ sex)}
\NormalTok{  pheno }\OtherTok{\textless{}{-}} \FunctionTok{get.pheno}\NormalTok{(population, }\AttributeTok{gen =}\NormalTok{ gen, }\AttributeTok{sex =}\NormalTok{ sex)}

  \CommentTok{\# Custom criterion: BV + 0.5*pheno}
\NormalTok{  score }\OtherTok{\textless{}{-}}\NormalTok{ bv[}\DecValTok{1}\NormalTok{,] }\SpecialCharTok{+} \FloatTok{0.5} \SpecialCharTok{*}\NormalTok{ pheno[}\DecValTok{1}\NormalTok{,]}

  \CommentTok{\# Return indices of selected individuals}
  \FunctionTok{return}\NormalTok{(}\FunctionTok{order}\NormalTok{(score, }\AttributeTok{decreasing =} \ConstantTok{TRUE}\NormalTok{)[}\DecValTok{1}\SpecialCharTok{:}\DecValTok{10}\NormalTok{])}
\NormalTok{\}}

\CommentTok{\# Use in breeding}
\NormalTok{pop }\OtherTok{\textless{}{-}} \FunctionTok{breeding.diploid}\NormalTok{(}
\NormalTok{  pop,}
  \AttributeTok{selection.m.function =}\NormalTok{ my\_selector,}
  \AttributeTok{selection.f.function =}\NormalTok{ my\_selector,}
  \AttributeTok{breeding.size =} \FunctionTok{c}\NormalTok{(}\DecValTok{50}\NormalTok{, }\DecValTok{50}\NormalTok{)}
\NormalTok{)}
\end{Highlighting}
\end{Shaded}

\section{Summary}\label{summary-7}

\begin{itemize}
\tightlist
\item
  Random, phenotypic, genomic, and true BV selection
\item
  Multi-trait selection indices
\item
  Threshold and OGC methods
\item
  Custom selection functions
\end{itemize}

Continue to \href{09-mating-designs.qmd}{Chapter 9: Mating Designs}!

\chapter{Mating Designs}\label{sec-mating-designs}

\section{Random Mating (Default)}\label{random-mating-default}

\begin{Shaded}
\begin{Highlighting}[]
\NormalTok{pop }\OtherTok{\textless{}{-}} \FunctionTok{breeding.diploid}\NormalTok{(}
\NormalTok{  pop,}
  \AttributeTok{selection.size =} \FunctionTok{c}\NormalTok{(}\DecValTok{10}\NormalTok{, }\DecValTok{20}\NormalTok{),}
  \AttributeTok{breeding.size =} \FunctionTok{c}\NormalTok{(}\DecValTok{100}\NormalTok{, }\DecValTok{100}\NormalTok{)}
\NormalTok{)}
\CommentTok{\# Parents randomly paired}
\end{Highlighting}
\end{Shaded}

\section{Assortative Mating}\label{assortative-mating}

\begin{Shaded}
\begin{Highlighting}[]
\CommentTok{\# Positive: Best × Best}
\NormalTok{pop }\OtherTok{\textless{}{-}} \FunctionTok{breeding.diploid}\NormalTok{(}
\NormalTok{  pop,}
  \AttributeTok{selection.size =} \FunctionTok{c}\NormalTok{(}\DecValTok{10}\NormalTok{, }\DecValTok{10}\NormalTok{),}
  \AttributeTok{mating.assortative =} \DecValTok{1}\NormalTok{,        }\CommentTok{\# Perfect positive}
  \AttributeTok{breeding.size =} \FunctionTok{c}\NormalTok{(}\DecValTok{100}\NormalTok{, }\DecValTok{100}\NormalTok{)}
\NormalTok{)}

\CommentTok{\# Negative: Best × Worst}
\NormalTok{pop }\OtherTok{\textless{}{-}} \FunctionTok{breeding.diploid}\NormalTok{(}
\NormalTok{  pop,}
  \AttributeTok{mating.assortative =} \SpecialCharTok{{-}}\DecValTok{1}\NormalTok{,       }\CommentTok{\# Perfect negative}
  \AttributeTok{breeding.size =} \FunctionTok{c}\NormalTok{(}\DecValTok{100}\NormalTok{, }\DecValTok{100}\NormalTok{)}
\NormalTok{)}
\end{Highlighting}
\end{Shaded}

\section{Avoiding Inbreeding}\label{avoiding-inbreeding}

\begin{Shaded}
\begin{Highlighting}[]
\CommentTok{\# Avoid full{-}sib and half{-}sib matings}
\NormalTok{pop }\OtherTok{\textless{}{-}} \FunctionTok{breeding.diploid}\NormalTok{(}
\NormalTok{  pop,}
  \AttributeTok{selection.size =} \FunctionTok{c}\NormalTok{(}\DecValTok{5}\NormalTok{, }\DecValTok{20}\NormalTok{),}
  \AttributeTok{avoid.mating.fullsib =} \ConstantTok{TRUE}\NormalTok{,}
  \AttributeTok{avoid.mating.halfsib =} \ConstantTok{TRUE}\NormalTok{,}
  \AttributeTok{breeding.size =} \FunctionTok{c}\NormalTok{(}\DecValTok{100}\NormalTok{, }\DecValTok{100}\NormalTok{)}
\NormalTok{)}
\end{Highlighting}
\end{Shaded}

\section{Maximum Offspring Per
Parent}\label{maximum-offspring-per-parent}

\begin{Shaded}
\begin{Highlighting}[]
\CommentTok{\# Limit offspring per sire}
\NormalTok{pop }\OtherTok{\textless{}{-}} \FunctionTok{breeding.diploid}\NormalTok{(}
\NormalTok{  pop,}
  \AttributeTok{selection.size =} \FunctionTok{c}\NormalTok{(}\DecValTok{5}\NormalTok{, }\DecValTok{20}\NormalTok{),}
  \AttributeTok{max.offspring =} \DecValTok{10}\NormalTok{,             }\CommentTok{\# Max 10 offspring per parent}
  \AttributeTok{breeding.size =} \FunctionTok{c}\NormalTok{(}\DecValTok{100}\NormalTok{, }\DecValTok{100}\NormalTok{)}
\NormalTok{)}

\CommentTok{\# Different limits by sex}
\NormalTok{pop }\OtherTok{\textless{}{-}} \FunctionTok{breeding.diploid}\NormalTok{(}
\NormalTok{  pop,}
  \AttributeTok{selection.size =} \FunctionTok{c}\NormalTok{(}\DecValTok{5}\NormalTok{, }\DecValTok{20}\NormalTok{),}
  \AttributeTok{max.offspring.m =} \DecValTok{20}\NormalTok{,           }\CommentTok{\# Males can have 20}
  \AttributeTok{max.offspring.f =} \DecValTok{5}\NormalTok{,            }\CommentTok{\# Females max 5}
  \AttributeTok{breeding.size =} \FunctionTok{c}\NormalTok{(}\DecValTok{100}\NormalTok{, }\DecValTok{100}\NormalTok{)}
\NormalTok{)}
\end{Highlighting}
\end{Shaded}

\section{All Combinations (Diallel)}\label{all-combinations-diallel}

\begin{Shaded}
\begin{Highlighting}[]
\CommentTok{\# Test all crosses}
\NormalTok{pop }\OtherTok{\textless{}{-}} \FunctionTok{breeding.diploid}\NormalTok{(}
\NormalTok{  pop,}
  \AttributeTok{selection.size =} \FunctionTok{c}\NormalTok{(}\DecValTok{10}\NormalTok{, }\DecValTok{10}\NormalTok{),}
  \AttributeTok{breeding.all.combination =} \ConstantTok{TRUE}\NormalTok{,  }\CommentTok{\# 10×10 = 100 crosses}
  \AttributeTok{breeding.size =} \FunctionTok{c}\NormalTok{(}\DecValTok{100}\NormalTok{, }\DecValTok{0}\NormalTok{)}
\NormalTok{)}
\end{Highlighting}
\end{Shaded}

\section{Targeted/Fixed Matings}\label{targetedfixed-matings}

\begin{Shaded}
\begin{Highlighting}[]
\CommentTok{\# Specify exact matings}
\CommentTok{\# Matrix: [generation, sex, individual, generation, sex, individual]}
\NormalTok{matings }\OtherTok{\textless{}{-}} \FunctionTok{rbind}\NormalTok{(}
  \FunctionTok{c}\NormalTok{(}\DecValTok{2}\NormalTok{, }\DecValTok{1}\NormalTok{, }\DecValTok{5}\NormalTok{,  }\DecValTok{2}\NormalTok{, }\DecValTok{2}\NormalTok{, }\DecValTok{10}\NormalTok{),  }\CommentTok{\# Male 5 × Female 10 from gen 2}
  \FunctionTok{c}\NormalTok{(}\DecValTok{2}\NormalTok{, }\DecValTok{1}\NormalTok{, }\DecValTok{8}\NormalTok{,  }\DecValTok{2}\NormalTok{, }\DecValTok{2}\NormalTok{, }\DecValTok{15}\NormalTok{),  }\CommentTok{\# Male 8 × Female 15}
  \FunctionTok{c}\NormalTok{(}\DecValTok{2}\NormalTok{, }\DecValTok{1}\NormalTok{, }\DecValTok{3}\NormalTok{,  }\DecValTok{2}\NormalTok{, }\DecValTok{2}\NormalTok{, }\DecValTok{20}\NormalTok{)   }\CommentTok{\# Male 3 × Female 20}
\NormalTok{)}

\NormalTok{pop }\OtherTok{\textless{}{-}} \FunctionTok{breeding.diploid}\NormalTok{(}
\NormalTok{  pop,}
  \AttributeTok{multiple.s =}\NormalTok{ matings,}
  \AttributeTok{copy.individual.m =} \DecValTok{5}\NormalTok{,           }\CommentTok{\# 5 offspring per mating}
  \AttributeTok{copy.individual.f =} \DecValTok{5}
\NormalTok{)}
\end{Highlighting}
\end{Shaded}

\section{Plant Breeding: Selfing}\label{plant-breeding-selfing}

\begin{Shaded}
\begin{Highlighting}[]
\NormalTok{pop }\OtherTok{\textless{}{-}} \FunctionTok{breeding.diploid}\NormalTok{(}
\NormalTok{  pop,}
  \AttributeTok{selection.size =} \FunctionTok{c}\NormalTok{(}\DecValTok{100}\NormalTok{, }\DecValTok{0}\NormalTok{),}
  \AttributeTok{selfing.mating =} \ConstantTok{TRUE}\NormalTok{,}
  \AttributeTok{breeding.size =} \FunctionTok{c}\NormalTok{(}\DecValTok{100}\NormalTok{, }\DecValTok{0}\NormalTok{)}
\NormalTok{)}
\end{Highlighting}
\end{Shaded}

\section{Plant Breeding: Cloning}\label{plant-breeding-cloning}

\begin{Shaded}
\begin{Highlighting}[]
\NormalTok{pop }\OtherTok{\textless{}{-}} \FunctionTok{breeding.diploid}\NormalTok{(}
\NormalTok{  pop,}
  \AttributeTok{selection.size =} \FunctionTok{c}\NormalTok{(}\DecValTok{50}\NormalTok{, }\DecValTok{0}\NormalTok{),}
  \AttributeTok{copy.individual.m =} \DecValTok{10}\NormalTok{,   }\CommentTok{\# 10 clones each}
  \AttributeTok{breeding.size =} \FunctionTok{c}\NormalTok{(}\DecValTok{500}\NormalTok{, }\DecValTok{0}\NormalTok{)}
\NormalTok{)}
\end{Highlighting}
\end{Shaded}

\section{Same-Sex Mating}\label{same-sex-mating}

\begin{Shaded}
\begin{Highlighting}[]
\CommentTok{\# For plants without sex distinction}
\NormalTok{pop }\OtherTok{\textless{}{-}} \FunctionTok{breeding.diploid}\NormalTok{(}
\NormalTok{  pop,}
  \AttributeTok{selection.size =} \FunctionTok{c}\NormalTok{(}\DecValTok{100}\NormalTok{, }\DecValTok{0}\NormalTok{),}
  \AttributeTok{same.sex.activ =} \ConstantTok{TRUE}\NormalTok{,}
  \AttributeTok{same.sex.sex =} \DecValTok{0}\NormalTok{,              }\CommentTok{\# Use "males" as both parents}
  \AttributeTok{breeding.size =} \FunctionTok{c}\NormalTok{(}\DecValTok{200}\NormalTok{, }\DecValTok{0}\NormalTok{)}
\NormalTok{)}
\end{Highlighting}
\end{Shaded}

\section{Summary}\label{summary-8}

\begin{itemize}
\tightlist
\item
  Random vs assortative mating
\item
  Inbreeding avoidance
\item
  Offspring number control
\item
  Fixed/targeted matings
\item
  Plant-specific: selfing, cloning
\end{itemize}

Continue to \href{10-breeding-simulation.qmd}{Chapter 10: Breeding
Simulation}!

\chapter{Breeding Program Simulation}\label{sec-breeding-simulation}

\section{Complete Breeding Cycle}\label{complete-breeding-cycle}

\begin{Shaded}
\begin{Highlighting}[]
\CommentTok{\# Typical cycle}
\NormalTok{pop }\OtherTok{\textless{}{-}} \FunctionTok{breeding.diploid}\NormalTok{(}
\NormalTok{  pop,}
  \CommentTok{\# 1. Phenotype}
  \AttributeTok{phenotyping =} \StringTok{"all"}\NormalTok{,}
  \AttributeTok{phenotyping.gen =} \DecValTok{5}\NormalTok{,}
  \AttributeTok{heritability =} \FloatTok{0.5}\NormalTok{,}

  \CommentTok{\# 2. Genotype (if needed)}
  \AttributeTok{genotyping.gen =} \DecValTok{5}\NormalTok{,}
  \AttributeTok{share.genotyped =} \FloatTok{0.8}\NormalTok{,}

  \CommentTok{\# 3. Predict breeding values}
  \AttributeTok{bve =} \ConstantTok{TRUE}\NormalTok{,}
  \AttributeTok{bve.gen =} \DecValTok{1}\SpecialCharTok{:}\DecValTok{5}\NormalTok{,}

  \CommentTok{\# 4. Select parents}
  \AttributeTok{selection.criteria =} \StringTok{"bve"}\NormalTok{,}
  \AttributeTok{selection.size =} \FunctionTok{c}\NormalTok{(}\DecValTok{10}\NormalTok{, }\DecValTok{20}\NormalTok{),}
  \AttributeTok{selection.gen =} \DecValTok{5}\NormalTok{,}

  \CommentTok{\# 5. Mate and create offspring}
  \AttributeTok{breeding.size =} \FunctionTok{c}\NormalTok{(}\DecValTok{50}\NormalTok{, }\DecValTok{50}\NormalTok{),}
  \AttributeTok{name.cohort =} \StringTok{"Gen6"}
\NormalTok{)}
\end{Highlighting}
\end{Shaded}

\section{Multi-Stage Selection}\label{multi-stage-selection}

\begin{Shaded}
\begin{Highlighting}[]
\CommentTok{\# Stage 1: Phenotypic pre{-}selection}
\NormalTok{pop }\OtherTok{\textless{}{-}} \FunctionTok{breeding.diploid}\NormalTok{(}
\NormalTok{  pop,}
  \AttributeTok{phenotyping =} \StringTok{"all"}\NormalTok{,}
  \AttributeTok{phenotyping.gen =} \DecValTok{3}\NormalTok{,}
  \AttributeTok{heritability =} \FloatTok{0.3}
\NormalTok{)}

\NormalTok{pop }\OtherTok{\textless{}{-}} \FunctionTok{set.class}\NormalTok{(pop, }\AttributeTok{gen =} \DecValTok{3}\NormalTok{, }\AttributeTok{class =} \DecValTok{1}\NormalTok{)  }\CommentTok{\# Mark all as candidates}

\CommentTok{\# Select top 50\% phenotypically}
\NormalTok{pheno }\OtherTok{\textless{}{-}} \FunctionTok{get.pheno}\NormalTok{(pop, }\AttributeTok{gen =} \DecValTok{3}\NormalTok{)}
\NormalTok{cutoff }\OtherTok{\textless{}{-}} \FunctionTok{quantile}\NormalTok{(pheno, }\FloatTok{0.5}\NormalTok{)}
\CommentTok{\# ... mark individuals above cutoff as class 2 ...}

\CommentTok{\# Stage 2: Genotype and genomic selection on survivors}
\NormalTok{pop }\OtherTok{\textless{}{-}} \FunctionTok{breeding.diploid}\NormalTok{(}
\NormalTok{  pop,}
  \AttributeTok{genotyping.gen =} \DecValTok{3}\NormalTok{,}
  \AttributeTok{genotyping.class =} \DecValTok{2}\NormalTok{,     }\CommentTok{\# Only genotype class 2}
  \AttributeTok{bve =} \ConstantTok{TRUE}\NormalTok{,}
  \AttributeTok{selection.criteria =} \StringTok{"bve"}\NormalTok{,}
  \AttributeTok{selection.class =} \DecValTok{2}\NormalTok{,       }\CommentTok{\# Select from class 2}
  \AttributeTok{selection.size =} \FunctionTok{c}\NormalTok{(}\DecValTok{10}\NormalTok{, }\DecValTok{10}\NormalTok{),}
  \AttributeTok{breeding.size =} \FunctionTok{c}\NormalTok{(}\DecValTok{100}\NormalTok{, }\DecValTok{100}\NormalTok{)}
\NormalTok{)}
\end{Highlighting}
\end{Shaded}

\section{Progeny Testing}\label{progeny-testing}

\begin{Shaded}
\begin{Highlighting}[]
\CommentTok{\# Year 1: Create test progeny}
\NormalTok{pop }\OtherTok{\textless{}{-}} \FunctionTok{breeding.diploid}\NormalTok{(}
\NormalTok{  pop,}
  \AttributeTok{selection.m.gen =} \DecValTok{5}\NormalTok{,}
  \AttributeTok{selection.f.gen =} \DecValTok{5}\NormalTok{,}
  \AttributeTok{selection.size =} \FunctionTok{c}\NormalTok{(}\DecValTok{50}\NormalTok{, }\DecValTok{50}\NormalTok{),}
  \AttributeTok{breeding.size =} \FunctionTok{c}\NormalTok{(}\DecValTok{500}\NormalTok{, }\DecValTok{500}\NormalTok{),    }\CommentTok{\# Many test progeny}
  \AttributeTok{name.cohort =} \StringTok{"TestProgeny"}
\NormalTok{)}

\CommentTok{\# Year 2: Phenotype progeny}
\NormalTok{pop }\OtherTok{\textless{}{-}} \FunctionTok{breeding.diploid}\NormalTok{(}
\NormalTok{  pop,}
  \AttributeTok{phenotyping.cohorts =} \StringTok{"TestProgeny"}\NormalTok{,}
  \AttributeTok{phenotyping =} \StringTok{"all"}\NormalTok{,}
  \AttributeTok{heritability =} \FloatTok{0.5}
\NormalTok{)}

\CommentTok{\# Calculate parent averages from progeny}
\NormalTok{pop }\OtherTok{\textless{}{-}} \FunctionTok{breeding.diploid}\NormalTok{(}
\NormalTok{  pop,}
  \AttributeTok{bve =} \ConstantTok{TRUE}\NormalTok{,}
  \AttributeTok{bve.gen =} \DecValTok{6}  \CommentTok{\# Use progeny phenotypes}
\NormalTok{)}

\CommentTok{\# Year 3: Select bulls based on daughter performance}
\NormalTok{pop }\OtherTok{\textless{}{-}} \FunctionTok{breeding.diploid}\NormalTok{(}
\NormalTok{  pop,}
  \AttributeTok{selection.criteria =} \StringTok{"bve"}\NormalTok{,}
  \AttributeTok{selection.m.gen =} \DecValTok{5}\NormalTok{,      }\CommentTok{\# Original candidate bulls}
  \AttributeTok{selection.size =} \FunctionTok{c}\NormalTok{(}\DecValTok{5}\NormalTok{, }\DecValTok{20}\NormalTok{),}
  \AttributeTok{breeding.size =} \FunctionTok{c}\NormalTok{(}\DecValTok{100}\NormalTok{, }\DecValTok{100}\NormalTok{)}
\NormalTok{)}
\end{Highlighting}
\end{Shaded}

\section{Nucleus Breeding Program}\label{nucleus-breeding-program}

\begin{Shaded}
\begin{Highlighting}[]
\CommentTok{\# Create nucleus}
\NormalTok{pop }\OtherTok{\textless{}{-}} \FunctionTok{creating.diploid}\NormalTok{(}
  \AttributeTok{nsnp =} \DecValTok{10000}\NormalTok{, }\AttributeTok{nindi =} \DecValTok{100}\NormalTok{,}
  \AttributeTok{n.additive =} \DecValTok{100}\NormalTok{,}
  \AttributeTok{name.cohort =} \StringTok{"Nucleus"}
\NormalTok{)}

\CommentTok{\# Nucleus selection (intensive)}
\NormalTok{pop }\OtherTok{\textless{}{-}} \FunctionTok{breeding.diploid}\NormalTok{(}
\NormalTok{  pop,}
  \AttributeTok{phenotyping.cohorts =} \StringTok{"Nucleus"}\NormalTok{,}
  \AttributeTok{heritability =} \FloatTok{0.5}\NormalTok{,}
  \AttributeTok{bve =} \ConstantTok{TRUE}\NormalTok{,}
  \AttributeTok{selection.criteria =} \StringTok{"bve"}\NormalTok{,}
  \AttributeTok{selection.size =} \FunctionTok{c}\NormalTok{(}\DecValTok{5}\NormalTok{, }\DecValTok{10}\NormalTok{),}
  \AttributeTok{breeding.size =} \FunctionTok{c}\NormalTok{(}\DecValTok{50}\NormalTok{, }\DecValTok{50}\NormalTok{),}
  \AttributeTok{name.cohort =} \StringTok{"Nucleus\_Gen2"}
\NormalTok{)}

\CommentTok{\# Commercial multiplier (uses nucleus genetics)}
\NormalTok{pop }\OtherTok{\textless{}{-}} \FunctionTok{breeding.diploid}\NormalTok{(}
\NormalTok{  pop,}
  \AttributeTok{selection.m.cohorts =} \StringTok{"Nucleus\_Gen2"}\NormalTok{,}
  \AttributeTok{selection.f.cohorts =} \StringTok{"Nucleus"}\NormalTok{,    }\CommentTok{\# Or commercial dams}
  \AttributeTok{selection.size =} \FunctionTok{c}\NormalTok{(}\DecValTok{10}\NormalTok{, }\DecValTok{100}\NormalTok{),}
  \AttributeTok{breeding.size =} \FunctionTok{c}\NormalTok{(}\DecValTok{500}\NormalTok{, }\DecValTok{500}\NormalTok{),}
  \AttributeTok{name.cohort =} \StringTok{"Commercial"}
\NormalTok{)}
\end{Highlighting}
\end{Shaded}

\section{Summary}\label{summary-9}

\begin{itemize}
\tightlist
\item
  Complete breeding cycles
\item
  Multi-stage selection
\item
  Progeny testing schemes
\item
  Nucleus breeding programs
\end{itemize}

Continue to \href{11-plant-vs-animal.qmd}{Chapter 11: Plant vs Animal
Breeding}!

\chapter{Plant vs Animal Breeding}\label{sec-plant-vs-animal}

\section{Plant-Specific Features}\label{plant-specific-features}

\subsection{Selfing}\label{selfing}

\begin{Shaded}
\begin{Highlighting}[]
\NormalTok{pop }\OtherTok{\textless{}{-}} \FunctionTok{breeding.diploid}\NormalTok{(pop, }\AttributeTok{selfing.mating =} \ConstantTok{TRUE}\NormalTok{, }\AttributeTok{breeding.size =} \FunctionTok{c}\NormalTok{(}\DecValTok{200}\NormalTok{, }\DecValTok{0}\NormalTok{))}
\end{Highlighting}
\end{Shaded}

\subsection{Doubled Haploids (DH)}\label{doubled-haploids-dh}

\begin{Shaded}
\begin{Highlighting}[]
\NormalTok{pop }\OtherTok{\textless{}{-}} \FunctionTok{breeding.diploid}\NormalTok{(pop, }\AttributeTok{dh.mating =} \ConstantTok{TRUE}\NormalTok{, }\AttributeTok{breeding.size =} \FunctionTok{c}\NormalTok{(}\DecValTok{200}\NormalTok{, }\DecValTok{0}\NormalTok{))}
\end{Highlighting}
\end{Shaded}

\subsection{Cloning}\label{cloning}

\begin{Shaded}
\begin{Highlighting}[]
\NormalTok{pop }\OtherTok{\textless{}{-}} \FunctionTok{breeding.diploid}\NormalTok{(pop, }\AttributeTok{copy.individual.m =} \DecValTok{10}\NormalTok{, }\AttributeTok{breeding.size =} \FunctionTok{c}\NormalTok{(}\DecValTok{1000}\NormalTok{, }\DecValTok{0}\NormalTok{))}
\end{Highlighting}
\end{Shaded}

\subsection{One-Sex Mode}\label{one-sex-mode-2}

\begin{Shaded}
\begin{Highlighting}[]
\NormalTok{pop }\OtherTok{\textless{}{-}} \FunctionTok{creating.diploid}\NormalTok{(}\AttributeTok{nsnp =} \DecValTok{5000}\NormalTok{, }\AttributeTok{nindi =} \DecValTok{200}\NormalTok{, }\AttributeTok{one.sex.mode =} \ConstantTok{TRUE}\NormalTok{)}
\end{Highlighting}
\end{Shaded}

\section{Animal-Specific Features}\label{animal-specific-features}

\subsection{Litter Size Control}\label{litter-size-control}

\begin{Shaded}
\begin{Highlighting}[]
\NormalTok{pop }\OtherTok{\textless{}{-}} \FunctionTok{breeding.diploid}\NormalTok{(}
\NormalTok{  pop,}
  \AttributeTok{selection.size =} \FunctionTok{c}\NormalTok{(}\DecValTok{10}\NormalTok{, }\DecValTok{20}\NormalTok{),}
  \AttributeTok{litter.size =} \DecValTok{8}\NormalTok{,              }\CommentTok{\# 8 piglets per litter}
  \AttributeTok{breeding.size =} \FunctionTok{c}\NormalTok{(}\DecValTok{160}\NormalTok{, }\DecValTok{160}\NormalTok{)   }\CommentTok{\# 20 litters}
\NormalTok{)}
\end{Highlighting}
\end{Shaded}

\subsection{Multiple Parities}\label{multiple-parities}

\begin{Shaded}
\begin{Highlighting}[]
\CommentTok{\# Sows have multiple litters}
\ControlFlowTok{for}\NormalTok{ (parity }\ControlFlowTok{in} \DecValTok{1}\SpecialCharTok{:}\DecValTok{3}\NormalTok{) \{}
\NormalTok{  pop }\OtherTok{\textless{}{-}} \FunctionTok{breeding.diploid}\NormalTok{(}
\NormalTok{    pop,}
    \AttributeTok{selection.f.gen =} \DecValTok{5}\NormalTok{,        }\CommentTok{\# Same females each time}
    \AttributeTok{selection.m.gen =} \DecValTok{6}\NormalTok{,        }\CommentTok{\# New males}
    \AttributeTok{breeding.size =} \FunctionTok{c}\NormalTok{(}\DecValTok{100}\NormalTok{, }\DecValTok{100}\NormalTok{)}
\NormalTok{  )}
\NormalTok{\}}
\end{Highlighting}
\end{Shaded}

\section{Summary}\label{summary-10}

\begin{itemize}
\tightlist
\item
  Plants: selfing, DH, cloning, no sex
\item
  Animals: litter size, multiple parities
\end{itemize}

Continue to \href{12-snp-panels.qmd}{Chapter 12: SNP Panels}!

\part{Genotyping \& Genomic Analysis}

\chapter{SNP Panels \& Genotyping}\label{sec-snp-panels}

\section{Partial Genotyping}\label{partial-genotyping}

\begin{Shaded}
\begin{Highlighting}[]
\CommentTok{\# Only genotype 80\% of individuals}
\NormalTok{pop }\OtherTok{\textless{}{-}} \FunctionTok{breeding.diploid}\NormalTok{(}
\NormalTok{  pop,}
  \AttributeTok{genotyping.gen =} \DecValTok{3}\NormalTok{,}
  \AttributeTok{share.genotyped =} \FloatTok{0.8}
\NormalTok{)}

\CommentTok{\# Check who\textquotesingle{}s genotyped}
\NormalTok{genotyped\_status }\OtherTok{\textless{}{-}} \FunctionTok{get.genotyped}\NormalTok{(pop, }\AttributeTok{gen =} \DecValTok{3}\NormalTok{)}
\end{Highlighting}
\end{Shaded}

\section{Genotyping Arrays (Chips)}\label{genotyping-arrays-chips}

\begin{Shaded}
\begin{Highlighting}[]
\CommentTok{\# Define which SNPs are on chip}
\NormalTok{chip\_snps }\OtherTok{\textless{}{-}} \FunctionTok{seq}\NormalTok{(}\DecValTok{1}\NormalTok{, }\DecValTok{50000}\NormalTok{, }\AttributeTok{by =} \DecValTok{10}\NormalTok{)  }\CommentTok{\# Every 10th SNP}

\NormalTok{pop }\OtherTok{\textless{}{-}} \FunctionTok{creating.diploid}\NormalTok{(}
  \AttributeTok{nsnp =} \DecValTok{50000}\NormalTok{,}
  \AttributeTok{nindi =} \DecValTok{500}\NormalTok{,}
  \AttributeTok{genotyped.s =} \FunctionTok{rep}\NormalTok{(}\DecValTok{0}\NormalTok{, }\DecValTok{50000}\NormalTok{)  }\CommentTok{\# Start with none genotyped}
\NormalTok{)}

\CommentTok{\# Set chip SNPs as genotyped}
\NormalTok{pop}\SpecialCharTok{$}\NormalTok{info}\SpecialCharTok{$}\NormalTok{snp}\SpecialCharTok{$}\NormalTok{genotyped[chip\_snps] }\OtherTok{\textless{}{-}} \DecValTok{1}
\end{Highlighting}
\end{Shaded}

\section{Genotyping Costs}\label{genotyping-costs}

\begin{Shaded}
\begin{Highlighting}[]
\CommentTok{\# Track costs}
\NormalTok{costs }\OtherTok{\textless{}{-}} \FunctionTok{compute.costs.cohorts}\NormalTok{(}
\NormalTok{  pop,}
  \AttributeTok{cohorts =} \StringTok{"Gen5"}\NormalTok{,}
  \AttributeTok{cost.genotyping =} \DecValTok{50}\NormalTok{,      }\CommentTok{\# $50 per sample}
  \AttributeTok{cost.phenotyping =} \DecValTok{200}     \CommentTok{\# $200 per phenotype}
\NormalTok{)}

\FunctionTok{print}\NormalTok{(costs}\SpecialCharTok{$}\NormalTok{total.costs)}
\end{Highlighting}
\end{Shaded}

\section{Summary}\label{summary-11}

\begin{itemize}
\tightlist
\item
  Partial genotyping strategies
\item
  SNP chip simulation
\item
  Cost tracking
\end{itemize}

Continue to \href{13-genomic-prediction.qmd}{Chapter 13: Genomic
Prediction}!

\chapter{Genomic Prediction}\label{sec-genomic-prediction}

\section{Internal GBLUP (Known h²)}\label{internal-gblup-known-huxb2}

\begin{Shaded}
\begin{Highlighting}[]
\NormalTok{pop }\OtherTok{\textless{}{-}} \FunctionTok{breeding.diploid}\NormalTok{(}
\NormalTok{  pop,}
  \AttributeTok{bve =} \ConstantTok{TRUE}\NormalTok{,}
  \AttributeTok{bve.gen =} \DecValTok{1}\SpecialCharTok{:}\DecValTok{5}\NormalTok{,}
  \AttributeTok{heritability =} \FloatTok{0.5}     \CommentTok{\# Must provide h²}
\NormalTok{)}
\end{Highlighting}
\end{Shaded}

\section{rrBLUP Package}\label{rrblup-package}

\begin{Shaded}
\begin{Highlighting}[]
\NormalTok{pop }\OtherTok{\textless{}{-}} \FunctionTok{breeding.diploid}\NormalTok{(}
\NormalTok{  pop,}
  \AttributeTok{bve =} \ConstantTok{TRUE}\NormalTok{,}
  \AttributeTok{bve.gen =} \DecValTok{1}\SpecialCharTok{:}\DecValTok{5}\NormalTok{,}
  \AttributeTok{bve.rrblup =} \ConstantTok{TRUE}     \CommentTok{\# Uses rrBLUP::mixed.solve}
\NormalTok{)}
\end{Highlighting}
\end{Shaded}

\section{Bayesian Methods (BGLR)}\label{bayesian-methods-bglr}

\begin{Shaded}
\begin{Highlighting}[]
\NormalTok{pop }\OtherTok{\textless{}{-}} \FunctionTok{breeding.diploid}\NormalTok{(}
\NormalTok{  pop,}
  \AttributeTok{bve =} \ConstantTok{TRUE}\NormalTok{,}
  \AttributeTok{bve.gen =} \DecValTok{1}\SpecialCharTok{:}\DecValTok{5}\NormalTok{,}
  \AttributeTok{bve.bglr =} \ConstantTok{TRUE}\NormalTok{,}
  \AttributeTok{bglr.model =} \StringTok{"BayesB"}   \CommentTok{\# BayesA, BayesB, BayesC, etc.}
\NormalTok{)}
\end{Highlighting}
\end{Shaded}

\section{BLUPF90}\label{blupf90}

\begin{Shaded}
\begin{Highlighting}[]
\NormalTok{pop }\OtherTok{\textless{}{-}} \FunctionTok{breeding.diploid}\NormalTok{(}
\NormalTok{  pop,}
  \AttributeTok{bve =} \ConstantTok{TRUE}\NormalTok{,}
  \AttributeTok{bve.gen =} \DecValTok{1}\SpecialCharTok{:}\DecValTok{5}\NormalTok{,}
  \AttributeTok{bve.blupf90 =} \ConstantTok{TRUE}\NormalTok{,}
  \AttributeTok{blupf90.path =} \StringTok{"/path/to/blupf90"}
\NormalTok{)}
\end{Highlighting}
\end{Shaded}

\section{Cross-Validation}\label{cross-validation}

\begin{Shaded}
\begin{Highlighting}[]
\CommentTok{\# Assess prediction accuracy}
\NormalTok{accuracy }\OtherTok{\textless{}{-}} \FunctionTok{analyze.bv}\NormalTok{(}
\NormalTok{  pop,}
  \AttributeTok{gen =} \DecValTok{5}\NormalTok{,}
  \AttributeTok{cohorts =} \StringTok{"TestSet"}\NormalTok{,}
  \AttributeTok{bve.gen =} \DecValTok{1}\SpecialCharTok{:}\DecValTok{4}         \CommentTok{\# Training set}
\NormalTok{)}

\FunctionTok{print}\NormalTok{(accuracy}\SpecialCharTok{$}\NormalTok{correlation)  }\CommentTok{\# Accuracy}
\FunctionTok{print}\NormalTok{(accuracy}\SpecialCharTok{$}\NormalTok{bias)          }\CommentTok{\# Bias}
\end{Highlighting}
\end{Shaded}

\section{Summary}\label{summary-12}

\begin{itemize}
\tightlist
\item
  Multiple BVE methods
\item
  Integration with R packages and external software
\item
  Accuracy assessment
\end{itemize}

Continue to \href{14-population-structure.qmd}{Chapter 14: Population
Structure}!

\chapter{Population Structure Analysis}\label{sec-population-structure}

\section{Inbreeding Coefficients}\label{inbreeding-coefficients}

\begin{Shaded}
\begin{Highlighting}[]
\CommentTok{\# Pedigree{-}based}
\NormalTok{inb\_ped }\OtherTok{\textless{}{-}} \FunctionTok{inbreeding.exp}\NormalTok{(pop, }\AttributeTok{gen =} \DecValTok{5}\NormalTok{)}

\CommentTok{\# Genomic}
\NormalTok{inb\_gen }\OtherTok{\textless{}{-}} \FunctionTok{inbreeding.emp}\NormalTok{(pop, }\AttributeTok{gen =} \DecValTok{5}\NormalTok{)}

\FunctionTok{mean}\NormalTok{(inb\_ped)}
\FunctionTok{mean}\NormalTok{(inb\_gen)}
\end{Highlighting}
\end{Shaded}

\section{Kinship Matrices}\label{kinship-matrices}

\begin{Shaded}
\begin{Highlighting}[]
\CommentTok{\# Expected (pedigree)}
\NormalTok{K\_exp }\OtherTok{\textless{}{-}} \FunctionTok{kinship.exp}\NormalTok{(pop, }\AttributeTok{gen =} \FunctionTok{c}\NormalTok{(}\DecValTok{5}\NormalTok{, }\DecValTok{6}\NormalTok{))}

\CommentTok{\# Empirical (genomic)}
\NormalTok{K\_emp }\OtherTok{\textless{}{-}} \FunctionTok{kinship.emp}\NormalTok{(pop, }\AttributeTok{gen =} \FunctionTok{c}\NormalTok{(}\DecValTok{5}\NormalTok{, }\DecValTok{6}\NormalTok{))}

\CommentTok{\# Average kinship}
\FunctionTok{mean}\NormalTok{(K\_emp[}\FunctionTok{upper.tri}\NormalTok{(K\_emp)])}
\end{Highlighting}
\end{Shaded}

\section{PCA Analysis}\label{pca-analysis}

\begin{Shaded}
\begin{Highlighting}[]
\CommentTok{\# Visualize population structure}
\FunctionTok{get.pca}\NormalTok{(pop, }\AttributeTok{gen =} \DecValTok{1}\SpecialCharTok{:}\DecValTok{10}\NormalTok{, }\AttributeTok{pca.color.breeding =} \ConstantTok{TRUE}\NormalTok{)}
\end{Highlighting}
\end{Shaded}

\section{Admixture Analysis}\label{admixture-analysis}

\begin{Shaded}
\begin{Highlighting}[]
\CommentTok{\# K=3 ancestral populations}
\FunctionTok{get.admixture}\NormalTok{(pop, }\AttributeTok{gen =} \DecValTok{5}\NormalTok{, }\AttributeTok{K =} \DecValTok{3}\NormalTok{, }\AttributeTok{plot =} \ConstantTok{TRUE}\NormalTok{)}
\end{Highlighting}
\end{Shaded}

\section{Dendrograms}\label{dendrograms}

\begin{Shaded}
\begin{Highlighting}[]
\CommentTok{\# Hierarchical clustering}
\FunctionTok{get.dendrogram}\NormalTok{(pop, }\AttributeTok{gen =} \DecValTok{5}\NormalTok{, }\AttributeTok{type =} \StringTok{"kinship"}\NormalTok{)}
\end{Highlighting}
\end{Shaded}

\section{Effective Population Size}\label{effective-population-size}

\begin{Shaded}
\begin{Highlighting}[]
\NormalTok{Ne }\OtherTok{\textless{}{-}} \FunctionTok{get.effective.size}\NormalTok{(pop, }\AttributeTok{gen =} \DecValTok{5}\NormalTok{)}
\FunctionTok{print}\NormalTok{(Ne)}
\end{Highlighting}
\end{Shaded}

\section{Summary}\label{summary-13}

\begin{itemize}
\tightlist
\item
  Inbreeding and kinship analysis
\item
  PCA and admixture visualization
\item
  Effective population size
\end{itemize}

Continue to \href{15-data-export-import.qmd}{Chapter 15: Data
Export/Import}!

\part{Advanced Topics \& Reference}

\chapter{Data Export \& Import}\label{sec-data-export-import}

\section{Export Genotypes}\label{export-genotypes}

\subsection{VCF Format}\label{vcf-format}

\begin{Shaded}
\begin{Highlighting}[]
\FunctionTok{get.vcf}\NormalTok{(pop, }\AttributeTok{gen =} \DecValTok{5}\NormalTok{, }\AttributeTok{path =} \StringTok{"generation5.vcf"}\NormalTok{)}
\end{Highlighting}
\end{Shaded}

\subsection{PLINK Format}\label{plink-format}

\begin{Shaded}
\begin{Highlighting}[]
\FunctionTok{get.plink}\NormalTok{(pop, }\AttributeTok{gen =} \DecValTok{5}\NormalTok{, }\AttributeTok{path =} \StringTok{"generation5"}\NormalTok{)  }\CommentTok{\# Creates .bed/.bim/.fam}
\FunctionTok{get.pedmap}\NormalTok{(pop, }\AttributeTok{gen =} \DecValTok{5}\NormalTok{, }\AttributeTok{path =} \StringTok{"generation5"}\NormalTok{) }\CommentTok{\# Creates .ped/.map}
\end{Highlighting}
\end{Shaded}

\section{Export Pedigrees}\label{export-pedigrees}

\begin{Shaded}
\begin{Highlighting}[]
\CommentTok{\# Standard pedigree (ID, Sire, Dam)}
\NormalTok{ped }\OtherTok{\textless{}{-}} \FunctionTok{get.pedigree}\NormalTok{(pop, }\AttributeTok{gen =} \DecValTok{5}\NormalTok{)}
\FunctionTok{write.csv}\NormalTok{(ped, }\StringTok{"pedigree.csv"}\NormalTok{)}

\CommentTok{\# Extended pedigree (more info)}
\NormalTok{ped2 }\OtherTok{\textless{}{-}} \FunctionTok{get.pedigree2}\NormalTok{(pop, }\AttributeTok{gen =} \DecValTok{5}\NormalTok{)  }\CommentTok{\# Includes sex, cohort, etc.}
\end{Highlighting}
\end{Shaded}

\section{Export Phenotypes}\label{export-phenotypes}

\begin{Shaded}
\begin{Highlighting}[]
\NormalTok{pheno }\OtherTok{\textless{}{-}} \FunctionTok{get.pheno}\NormalTok{(pop, }\AttributeTok{gen =} \DecValTok{1}\SpecialCharTok{:}\DecValTok{5}\NormalTok{)}
\FunctionTok{write.csv}\NormalTok{(}\FunctionTok{t}\NormalTok{(pheno), }\StringTok{"phenotypes.csv"}\NormalTok{)}
\end{Highlighting}
\end{Shaded}

\section{Export Breeding Values}\label{export-breeding-values}

\begin{Shaded}
\begin{Highlighting}[]
\NormalTok{bv }\OtherTok{\textless{}{-}} \FunctionTok{get.bv}\NormalTok{(pop, }\AttributeTok{gen =} \DecValTok{1}\SpecialCharTok{:}\DecValTok{10}\NormalTok{)}
\NormalTok{bve }\OtherTok{\textless{}{-}} \FunctionTok{get.bve}\NormalTok{(pop, }\AttributeTok{gen =} \DecValTok{1}\SpecialCharTok{:}\DecValTok{10}\NormalTok{)}

\NormalTok{data\_export }\OtherTok{\textless{}{-}} \FunctionTok{data.frame}\NormalTok{(}
  \AttributeTok{ID =} \FunctionTok{colnames}\NormalTok{(bv),}
  \AttributeTok{TrueBV =}\NormalTok{ bv[}\DecValTok{1}\NormalTok{,],}
  \AttributeTok{EstimatedBV =}\NormalTok{ bve[}\DecValTok{1}\NormalTok{,]}
\NormalTok{)}
\FunctionTok{write.csv}\NormalTok{(data\_export, }\StringTok{"breeding\_values.csv"}\NormalTok{)}
\end{Highlighting}
\end{Shaded}

\section{Import External BVEs}\label{import-external-bves}

\begin{Shaded}
\begin{Highlighting}[]
\CommentTok{\# Calculate BVEs externally, then import}
\NormalTok{external\_bve }\OtherTok{\textless{}{-}} \FunctionTok{read.csv}\NormalTok{(}\StringTok{"external\_predictions.csv"}\NormalTok{)}

\NormalTok{pop }\OtherTok{\textless{}{-}} \FunctionTok{insert.bve}\NormalTok{(}
\NormalTok{  pop,}
  \AttributeTok{bve =}\NormalTok{ external\_bve}\SpecialCharTok{$}\NormalTok{EBV,}
  \AttributeTok{gen =} \DecValTok{5}
\NormalTok{)}
\end{Highlighting}
\end{Shaded}

\section{Summary}\label{summary-14}

\begin{itemize}
\tightlist
\item
  Export to VCF, PLINK formats
\item
  Pedigree and phenotype export
\item
  Import external predictions
\end{itemize}

Continue to \href{16-performance-memory.qmd}{Chapter 16: Performance \&
Memory}!

\chapter{Performance \& Memory Management}\label{sec-performance-memory}

\section{Memory Management}\label{memory-management}

\subsection{Clean Up Old Generations}\label{clean-up-old-generations}

\begin{Shaded}
\begin{Highlighting}[]
\CommentTok{\# Remove data from old generations}
\NormalTok{pop }\OtherTok{\textless{}{-}} \FunctionTok{clean.up}\NormalTok{(pop, }\AttributeTok{gen =} \DecValTok{1}\SpecialCharTok{:}\DecValTok{5}\NormalTok{)  }\CommentTok{\# Keep gens 6+}
\end{Highlighting}
\end{Shaded}

\subsection{New Base Generation}\label{new-base-generation}

\begin{Shaded}
\begin{Highlighting}[]
\CommentTok{\# Reset generation numbering, free memory}
\NormalTok{pop }\OtherTok{\textless{}{-}} \FunctionTok{new.base.generation}\NormalTok{(pop, }\AttributeTok{base.gen =} \DecValTok{10}\NormalTok{)}
\CommentTok{\# Gen 10 becomes new "gen 1"}
\end{Highlighting}
\end{Shaded}

\section{Computational Speed}\label{computational-speed}

\subsection{Use miraculix}\label{use-miraculix}

\begin{Shaded}
\begin{Highlighting}[]
\CommentTok{\# Install for fast matrix operations}
\NormalTok{devtools}\SpecialCharTok{::}\FunctionTok{install\_github}\NormalTok{(}\StringTok{"tpook92/miraculix"}\NormalTok{)}

\CommentTok{\# Automatically used when available}
\end{Highlighting}
\end{Shaded}

\subsection{Optimize Cores}\label{optimize-cores}

\begin{Shaded}
\begin{Highlighting}[]
\CommentTok{\# Find optimal number of cores for your system}
\NormalTok{optimal }\OtherTok{\textless{}{-}} \FunctionTok{optimize.cores}\NormalTok{(pop)}
\end{Highlighting}
\end{Shaded}

\subsection{Reduce Marker Number}\label{reduce-marker-number}

\begin{Shaded}
\begin{Highlighting}[]
\CommentTok{\# Fewer SNPs = faster}
\NormalTok{pop }\OtherTok{\textless{}{-}} \FunctionTok{creating.diploid}\NormalTok{(}
  \AttributeTok{nsnp =} \DecValTok{1000}\NormalTok{,     }\CommentTok{\# Instead of 50000}
  \AttributeTok{nindi =} \DecValTok{1000}\NormalTok{,}
  \AttributeTok{n.additive =} \DecValTok{100}
\NormalTok{)}
\end{Highlighting}
\end{Shaded}

\subsection{Disable Copying}\label{disable-copying}

\begin{Shaded}
\begin{Highlighting}[]
\CommentTok{\# Don\textquotesingle{}t store full breeding values in memory}
\NormalTok{pop }\OtherTok{\textless{}{-}} \FunctionTok{breeding.diploid}\NormalTok{(}
\NormalTok{  pop,}
  \AttributeTok{copy.breeding.values =} \ConstantTok{FALSE}
\NormalTok{)}
\end{Highlighting}
\end{Shaded}

\section{Parallel Processing}\label{parallel-processing}

\begin{Shaded}
\begin{Highlighting}[]
\CommentTok{\# Run multiple simulations in parallel}
\FunctionTok{library}\NormalTok{(parallel)}

\NormalTok{results }\OtherTok{\textless{}{-}} \FunctionTok{mclapply}\NormalTok{(}\DecValTok{1}\SpecialCharTok{:}\DecValTok{100}\NormalTok{, }\ControlFlowTok{function}\NormalTok{(i) \{}
\NormalTok{  pop }\OtherTok{\textless{}{-}} \FunctionTok{creating.diploid}\NormalTok{(}\AttributeTok{nsnp =} \DecValTok{5000}\NormalTok{, }\AttributeTok{nindi =} \DecValTok{200}\NormalTok{)}
  \CommentTok{\# ... run simulation ...}
  \FunctionTok{return}\NormalTok{(}\FunctionTok{mean}\NormalTok{(}\FunctionTok{get.bv}\NormalTok{(pop, }\AttributeTok{gen =} \DecValTok{10}\NormalTok{)))}
\NormalTok{\}, }\AttributeTok{mc.cores =} \DecValTok{4}\NormalTok{)}
\end{Highlighting}
\end{Shaded}

\section{Summary}\label{summary-15}

\begin{itemize}
\tightlist
\item
  Memory management techniques
\item
  Speed optimization
\item
  Parallel simulation
\end{itemize}

Continue to \href{17-function-reference.qmd}{Chapter 17: Function
Reference}!

\chapter{Function Reference Quick Guide}\label{sec-function-reference}

\section{Core Functions}\label{core-functions}

\begin{longtable}[]{@{}ll@{}}
\toprule\noalign{}
Function & Purpose \\
\midrule\noalign{}
\endhead
\bottomrule\noalign{}
\endlastfoot
\texttt{creating.diploid()} & Create founder population \\
\texttt{breeding.diploid()} & Simulate breeding actions \\
\texttt{creating.trait()} & Add traits to population \\
\end{longtable}

\section{Extraction Functions (get.*)}\label{extraction-functions-get.}

\begin{longtable}[]{@{}ll@{}}
\toprule\noalign{}
Function & Returns \\
\midrule\noalign{}
\endhead
\bottomrule\noalign{}
\endlastfoot
\texttt{get.bv()} & True breeding values \\
\texttt{get.bve()} & Estimated breeding values \\
\texttt{get.pheno()} & Phenotypes \\
\texttt{get.geno()} & Genotypes (0/1/2 matrix) \\
\texttt{get.haplo()} & Haplotypes \\
\texttt{get.pedigree()} & Pedigree (ID, Sire, Dam) \\
\texttt{get.map()} & Genetic map \\
\texttt{get.qtl()} & QTL positions \\
\end{longtable}

\section{Analysis Functions}\label{analysis-functions}

\begin{longtable}[]{@{}ll@{}}
\toprule\noalign{}
Function & Purpose \\
\midrule\noalign{}
\endhead
\bottomrule\noalign{}
\endlastfoot
\texttt{kinship.emp()} & Genomic kinship matrix \\
\texttt{kinship.exp()} & Pedigree kinship matrix \\
\texttt{inbreeding.emp()} & Genomic inbreeding \\
\texttt{inbreeding.exp()} & Pedigree inbreeding \\
\texttt{get.pca()} & PCA plot \\
\texttt{get.admixture()} & Admixture analysis \\
\texttt{bv.development()} & Genetic trend plot \\
\texttt{analyze.bv()} & Prediction accuracy \\
\end{longtable}

\section{Utility Functions}\label{utility-functions}

\begin{longtable}[]{@{}ll@{}}
\toprule\noalign{}
Function & Purpose \\
\midrule\noalign{}
\endhead
\bottomrule\noalign{}
\endlastfoot
\texttt{summary()} & Population summary \\
\texttt{get.size()} & Number of individuals \\
\texttt{set.class()} & Assign class to individuals \\
\texttt{clean.up()} & Free memory \\
\texttt{new.base.generation()} & Reset generation numbering \\
\end{longtable}

\section{Export Functions}\label{export-functions}

\begin{longtable}[]{@{}ll@{}}
\toprule\noalign{}
Function & Purpose \\
\midrule\noalign{}
\endhead
\bottomrule\noalign{}
\endlastfoot
\texttt{get.vcf()} & Export VCF file \\
\texttt{get.plink()} & Export PLINK files \\
\texttt{get.pedmap()} & Export ped/map files \\
\end{longtable}

\section{For Full Documentation}\label{for-full-documentation}

\begin{Shaded}
\begin{Highlighting}[]
\CommentTok{\# Function help}
\NormalTok{?creating.diploid}
\NormalTok{?breeding.diploid}

\CommentTok{\# Package help}
\FunctionTok{help}\NormalTok{(}\AttributeTok{package =} \StringTok{"MoBPS"}\NormalTok{)}
\end{Highlighting}
\end{Shaded}

\section{Summary}\label{summary-16}

Quick reference for common MoBPS functions. See the full manual for
complete parameter lists and details.

\bookmarksetup{startatroot}

\chapter*{References}\label{references}
\addcontentsline{toc}{chapter}{References}

\markboth{References}{References}

\section*{Key MoBPS Publications}\label{key-mobps-publications}
\addcontentsline{toc}{section}{Key MoBPS Publications}

\markright{Key MoBPS Publications}

\textbf{Main Paper:}

Pook, T., Freudenthal, J., Korte, A., \& Simianer, H. (2020). Using
local convolutional neural networks for genomic prediction.
\emph{Frontiers in Genetics}, 11, 561497.

Pook, T., Schlather, M., \& Simianer, H. (2020). MoBPS - Modular
Breeding Program Simulator. \emph{G3: Genes, Genomes, Genetics}, 10(6),
1915-1918.

\section*{Additional Resources}\label{additional-resources}
\addcontentsline{toc}{section}{Additional Resources}

\markright{Additional Resources}

\begin{itemize}
\tightlist
\item
  \textbf{GitHub Repository:} \url{https://github.com/tpook92/MoBPS}
\item
  \textbf{Web Interface:} \url{https://www.mobps.de}
\item
  \textbf{CRAN Page:} \url{https://CRAN.R-project.org/package=MoBPS}
\end{itemize}

\section*{Contact}\label{contact}
\addcontentsline{toc}{section}{Contact}

\markright{Contact}

For questions, bug reports, or collaborations:

\begin{itemize}
\tightlist
\item
  \textbf{Torsten Pook:}
  \href{mailto:Torsten.pook@wur.nl}{\nolinkurl{Torsten.pook@wur.nl}}
\item
  Wageningen University \& Research
\item
  Animal Breeding and Genomics Group
\end{itemize}

\section*{Citation}\label{citation}
\addcontentsline{toc}{section}{Citation}

\markright{Citation}

If you use MoBPS in your research, please cite:

\begin{verbatim}
@article{pook2020mobps,
  title={MoBPS—Modular breeding program simulator},
  author={Pook, Torsten and Schlather, Martin and Simianer, Henner},
  journal={G3: Genes, Genomes, Genetics},
  volume={10},
  number={6},
  pages={1915--1918},
  year={2020},
  publisher={Oxford University Press}
}
\end{verbatim}

\section*{Related Software}\label{related-software}
\addcontentsline{toc}{section}{Related Software}

\markright{Related Software}

\begin{itemize}
\tightlist
\item
  \textbf{AlphaSimR:} Another breeding simulation package
\item
  \textbf{QMSim:} Quantitative and molecular simulation
\item
  \textbf{Genomic Selection packages:} rrBLUP, BGLR, sommer, BLUPF90
\end{itemize}


\backmatter


\end{document}
